\chapter{The non-Archimedean Monge--Amp\`ere operator}
\index{Monge--Ampere measure!BFJ}\index{nonpluripolar product!BFJ}
\label{chap:bfj-ma}

The psh functions of \Cref{chap:bfj-psh} are defined by positivity on
all models. We now turn that positivity into a measure. There are two steps,
and it is important not to merge them. For model functions the measure is an
intersection product on one model. For bounded psh functions it is the unique
monotone extension of those atomic measures. The first step is algebraic; the
second is a non-Archimedean Bedford--Taylor theorem.

Throughout the chapter we retain the original BFJ hypotheses: \(K\) is
complete and discretely valued, its residue field has characteristic zero,
\(X\) is smooth and projective of dimension \(n\), and
\(\theta\in\mathcal Z^{1,1}(X)\) has ample generic numerical class
\(\{\theta\}\in N^1(X)\), in the sense of \eqref{eq:ch20-generic-class}. Put
\[
                         V=(\{\theta\}^n)>0.
\]
Thus our Monge--Amp\`ere measures have mass \(V\), not mass one. The main
construction and all monotone-continuity statements in this chapter are
\cite[\S\S2.7 and~3]{BFJ15}.

\section{Orientation}
\label{sec:ch24-orientation}

\subsection{What has to be built, and why}
\label{sec:ch24-why}

By the end of \Cref{chap:bfj-envelopes} the linear theory is complete.  We
have a space \(\mathcal Z^{1,1}(X)\) of closed model forms
(\Cref{def:ch20-closed-model-form}), an operator \(\ddc\) from model
functions to closed model forms with known kernel and image
(\Cref{thm:ch20-ddc-exact}), a class \(\PSH(X,\theta)\) of potentials that is
stable under the operations one wants (\Cref{thm:ch21-stability}), a
compactness theorem modulo constants (\Cref{thm:bfj-psh-compactness}), a
decreasing regularization theorem (\Cref{thm:bfj-decreasing-regularization})
and a continuity theorem for envelopes (\Cref{thm:bfj-envelope-regularity}).
What is entirely missing is a \emph{nonlinear} invariant: a way of assigning
to a potential a measure whose total mass is an intersection number and whose
local behaviour reflects the curvature of the potential.  That invariant is
the Monge--Amp\`ere measure, and constructing it is the business of this
chapter.

At model level there is no difficulty and no choice.  A closed model form
determined on \(\cX\) restricts to each component \(E\) of the special fibre
\(\cX_s=\sum_Eb_EE\); an \(n\)-fold product of such restrictions is a number,
computed on the projective \(k\)-variety \(E\); weighting by the multiplicity
\(b_E\) and putting the result at the divisorial point \(x_E\) produces an
atomic measure.  That this is \emph{forced}---by multilinearity together with
the requirement that the total mass be the generic intersection
number---was already argued in \Cref{sec:metrics-cl-measure}, where the same
measure appeared under the name of Chambert--Loir.  The only work at this
level is bookkeeping: independence of the model, symmetry, positivity, mass,
and the identity that converts an integral of a model function into a
vertical intersection number.

The difficulty begins one step later.  A bounded \(\theta\)-psh function is
not determined on any single model; by
\Cref{thm:bfj-decreasing-regularization} it is only a decreasing limit of a
\emph{net} of \(\theta\)-psh model functions, and by
\Cref{ex:ch22-no-countable-cofinal} a sequence will not in general do.  One
would like to define its Monge--Amp\`ere measure as the weak limit of the
model measures along such a net.  Three separate things then have to be
proved, and they organize \Cref{sec:ch24-bounded}.

\begin{enumerate}
\item\label{itm:ch24-plan:exists}
 \emph{The limit has to exist.}  Weak compactness of the space of positive
 measures of mass \(V\) (\Cref{prop:measure-compactness}) gives convergent
 subnets, not convergence.  Convergence has to come from a monotonicity
 property of the model measures, and the only monotone quantities available
 are \emph{numbers}, not measures.
\item\label{itm:ch24-plan:unique}
 \emph{The limit has to be independent of the net.}  Two decreasing nets of
 model functions with the same limit are not comparable index by index; the
 comparison has to be made by a compactness argument on \(X^{\An}\)
 (\Cref{lem:ch24-sandwich}), which is where Dini's theorem for nets,
 \Cref{prop:dini-nets}, enters.
\item\label{itm:ch24-plan:symmetric}
 \emph{The result has to be symmetric and multilinear.}  Any construction
 that treats one slot at a time produces an object that is manifestly
 well defined but not manifestly symmetric.  Symmetry is recovered by
 factoring the whole construction through a single \emph{symmetric} scalar
 invariant, the mixed energy pairing of \Cref{def:ch24-mixed-energy}.
\end{enumerate}

The device that settles all three at once is the following.  Instead of
extending the measure-valued operator directly, we extend the
\((n+1)\)-variable scalar
\[
 \mathfrak E(\varphi_0,\ldots,\varphi_n),
\]
the mixed energy pairing, which at model level is the sum of all vertical
intersection numbers built from the potentials.  It is symmetric, monotone
in each variable, and equal to
\(\mathfrak E(0,\varphi_1,\ldots,\varphi_n)
 +\int\varphi_0\,\MA_\theta(\varphi_1,\ldots,\varphi_n)\).
Monotone quantities have monotone limits, so \(\mathfrak E\) extends without
any compactness argument beyond \Cref{lem:monotone-net-converges}, and the
measure is then recovered from it by the Riesz representation theorem
(\Cref{thm:riesz-interface}).  Symmetry, multilinearity, monotone
continuity, and the weighted monotone continuity
\eqref{eq:weighted-monotone-ma} all become statements about
\(\mathfrak E\), where they are visible.  This is the strategy of
\Cref{sec:ch24-energy-bounded,sec:ch24-extension-theorem}.

\subsection{Four Monge--Amp\`ere products, kept apart}
\label{sec:ch24-four-products}

A single symbol \(\MA_\theta\) is used in this book for four different
objects.  They agree where their domains overlap---that is a theorem, not a
convention---but they are constructed by four different mechanisms and enjoy
four different continuity properties.  The governing rule of
\Cref{part:bfj} is that no statement may be made without saying which of them
is meant.

\begin{center}
\small
\begin{tabular}{p{.16\textwidth}|p{.21\textwidth}|p{.24\textwidth}|p{.24\textwidth}}
\toprule
\textbf{Product}&\textbf{Domain}&\textbf{Mechanism}&\textbf{Continuity available}\\
\midrule
model
 &\(\theta\)-psh model functions
 &restriction to components of \(\cX_s\), \eqref{eq:model-mixed-measure}
 &none needed; the measure is a finite sum of Dirac masses\\
continuous
 &continuous \(\theta\)-psh, i.e.\ uniform limits of model ones
 &uniform density and the mass bound \(V\), \Cref{prop:ch24-continuous-uniform}
 &continuous for the uniform topology on potentials\\
bounded
 &bounded \(\theta\)-psh
 &monotone extension, \Cref{thm:bounded-ma-extension}
 &continuous along decreasing nets with bounded limit, in the weighted form
  \eqref{eq:weighted-monotone-ma}\\
finite energy
 &\(\cE^1(X,\theta)\), \Cref{def:e1}
 &truncation and a nonpluripolar convention, \Cref{def:nonpluripolar-ma}
 &continuous along decreasing nets, and along strong convergence;
  \Cref{thm:e1-mixed-products}\\
\bottomrule
\end{tabular}
\end{center}

This chapter constructs the first three.  The fourth is
\Cref{chap:bfj-energy} and is \emph{not} obtained by taking a limit of the
constructions here: an unbounded potential is a decreasing limit of bounded
ones, but the measures need not converge to anything of mass \(V\), and mass
may be lost to the polar set.  \Cref{prop:bounded-ma-no-polar-mass} is the
precise statement that no mass is lost \emph{as long as the potentials stay
bounded}, and \Cref{def:nonpluripolar-ma} is the convention that
reinstates the loss for unbounded ones.  A fifth product, the
Chambert--Loir--Ducros Bedford--Taylor product of \Cref{chap:cld-bedford-taylor},
is a local object of a different nature and is compared with the present one
only on the model and continuous loci; see
\Cref{sec:ch24-comparison} and the delimitation table
\Cref{sec:ch16-delimitation-table}.

\begin{warning}{One symbol, four operators}
Whenever \(\MA_\theta\) occurs below, the reader should be able to say
without looking back which row of the table is in force.  The three
constructions of this chapter are compatible by
\Cref{thm:bounded-ma-extension}\,(1) and
\Cref{prop:ch24-continuous-uniform}; compatibility with the fourth is
\Cref{thm:e1-mixed-products}, and compatibility with the CLD product is
\Cref{thm:bfj-cld-ma-common-locus}, whose scope clause is not decorative.
No statement of this chapter concerns unbounded potentials except
\Cref{prop:bounded-ma-no-polar-mass}, in which the unbounded function is
integrated \emph{against} a bounded product and is never itself an entry.
\end{warning}

\subsection{The payoff, stated in advance}
\label{sec:ch24-payoff}

Five statements proved below are used continually in
\Crefrange{chap:bfj-capacity}{chap:bfj-examples}, and it is worth seeing them
before the constructions.

\begin{enumerate}
\item \emph{The mixed model measure is a positive atomic measure whose total
 mass is the mixed intersection number \((\{\theta_1\}\cdots\{\theta_n\})\)}
 (\Cref{prop:model-ma-basic}, \Cref{thm:ch24-mass}).  All later mass
 bookkeeping---capacity in \Cref{chap:bfj-capacity}, energy in
 \Cref{chap:bfj-energy}, the variational functional in
 \Cref{chap:bfj-finite-energy-measures}---is this one number.
\item \emph{Integrating a model function against a mixed model measure gives
 a vertical intersection number} (\eqref{eq:model-test-intersection}, which
 is \eqref{eq:metrics-cl-pairing}).  Every computation in this chapter is
 an application of this identity.
\item \emph{The mixed energy pairing is symmetric, monotone, and continuous
 along decreasing nets} (\Cref{prop:ch24-energy-basic},
 \Cref{thm:ch24-energy-extension}).  This is the engine; the extension
 theorem is a corollary of it.
\item \emph{The Monge--Amp\`ere operator extends uniquely to bounded psh
 potentials, with weighted monotone continuity}
 (\Cref{thm:bounded-ma-extension}).  This is the non-Archimedean
 Bedford--Taylor theorem, and it is what makes \(\MA_\theta\) an operator
 rather than a formula.
\item \emph{Monge--Amp\`ere measures of bounded potentials are local on the
 set where two potentials differ} (\Cref{prop:model-maximum-principle} at
 model level, \Cref{thm:bounded-maximum-principle} in general), and satisfy
 the comparison principle (\Cref{prop:ch24-model-comparison} at model
 level).  Locality is what converts the operator into a tool for uniqueness
 statements.
\end{enumerate}

A sixth statement is negative and belongs with the others.  Nothing in this
chapter identifies a bounded \(\theta\)-psh function with a locally CLD-psh
function, or the bounded monotone product with a CLD Bedford--Taylor product
outside the model and continuous loci.  The exact boundary is
\Cref{thm:ch19-cld-bfj-comparison}\,(C), quoted in
\Cref{sec:ch24-scope-comparison}.

\subsection{Standing hypotheses}
\label{sec:ch24-hypotheses}

Throughout the chapter we work under the original BFJ signature
\eqref{eq:ch20-standing}, which we repeat in the form used here.

\begin{equation}
\label{eq:ch24-standing}
\left\{
\begin{aligned}
 &\text{\(K\) is complete and \emph{discretely} valued, with residue field
   \(k\) of characteristic \(0\);}\\
 &\text{a uniformizer \(\pi\) is normalized by \(v_K(\pi)=1\); }\\
 &\text{\(X\) is a smooth, connected, projective \(K\)-variety of dimension
   \(n\geq1\), \(R=K^\circ\);}\\
 &\text{models \(\cX\) are normal, flat and projective over \(R\) with
   \(\cX_K\simeq X\);}\\
 &\text{\(\theta\in\mathcal Z^{1,1}(X)\) is a closed model form whose generic
   class \(\{\theta\}\) is ample.}
\end{aligned}
\right.
\end{equation}

By \Cref{prop:ch20-positive-representative} we may and do assume, whenever it
is convenient, that \(\theta\) is itself \(\cX\)-positive on some model, so
that in particular \(0\in\PSH(X,\theta)\) and constants are
\(\theta\)-psh (\Cref{prop:ch21-constants}).  Replacing \(\theta\) by
\(\theta+\ddc\psi\) and every potential \(\varphi\) by \(\varphi-\psi\)
changes nothing in the statements below, by
\Cref{cor:ch21-reference-independence}; the normalization is recorded
explicitly at each place where it is used.  We write

\begin{equation}
\label{eq:ch24-volume}
                V=(\{\theta\}^n)>0
\end{equation}

for the common total mass of all the measures constructed here.  BFJ
normalize \(V=1\) in \cite{BFJ15}; we do not, and every displayed formula
carries the factor.  The probability normalization \(V^{-1}\MA_\theta\) is
used nowhere in this chapter and is recorded in
\eqref{eq:appJ-prob-ma} only so that formulas quoted from the literature can
be converted.

Two normalizations are used constantly and are quoted so that no proof has
to rederive them.  The first is the sign of a model function,
\eqref{eq:ch20-sign}: if \(D\) is a vertical Cartier divisor with local
equation \(f_D\), then \(\varphi_D=-\log|f_D|=v_\bullet(f_D)\), so effective
vertical divisors give nonnegative model functions and
\(\varphi_{\cX_s}=1\).  The second is the divisorial normalization
\(x_E=b_E^{-1}\ord_E\) of \Cref{rem:divisorial-normalization}, which is why
the multiplicity \(b_E\) appears in \eqref{eq:model-mixed-measure} and
cancels in \eqref{eq:model-test-intersection}.

\subsection{What is proved here and what is imported}
\label{sec:ch24-imported}

Everything in
\Crefrange{sec:ch24-model-measure}{sec:ch24-comparison} is proved, with the
following exceptions.

\begin{enumerate}
\item\label{itm:ch24-import:intersection}
 \emph{Intersection theory of Cartier divisors on a scheme}, in the form of
 commutativity of the first Chern class operators and the projection
 formula, \cite[Chapters~1--2]{Ful}, used through
 \Cref{prop:appF-mixed-cartier}, \Cref{thm:appF-projection},
 \Cref{lem:vertical-intersection-degree} and
 \Cref{thm:model-projection-formula}.  This is the only input to
 \Cref{lem:ch24-vertical-numbers}, and hence to the whole of
 \Cref{sec:ch24-model-measure}.
\item\label{itm:ch24-import:kleiman}
 \emph{Openness of the ample cone} on a projective scheme over a field,
 which is Kleiman's criterion in the form already imported in
 \Cref{chap:bfj-model-forms}
 (\Cref{sec:ch20-imported}, item~\cref{itm:ch20-import:kleiman}); see
 \cite[\S1.4]{PAG1}.  It is used once, in \Cref{lem:ch24-psh-difference},
 to write an arbitrary model function as a multiple of a difference of two
 \(\theta\)-psh model functions.
\item\label{itm:ch24-import:regularization}
 \emph{Decreasing regularization}, \Cref{thm:bfj-decreasing-regularization},
 proved in \Cref{chap:bfj-compactness} from the envelope-continuity theorem
 \Cref{thm:bfj-envelope-regularity}.  This is where residue characteristic
 zero enters the present chapter, and it enters only here; see
 \Cref{sec:ch24-scope}.
\item\label{itm:ch24-import:density}
 \emph{Density of model functions} in \(C^0(X^{\An})\),
 \Cref{prop:metrics-model-function-density}, and the Hodge-type inequality
 \Cref{thm:ch20-mixed-hodge}.  Both are proved earlier in this book.
\item\label{itm:ch24-import:cld}
 \emph{The CLD comparison} \Cref{thm:ch19-cld-bfj-comparison}, proved in
 \Cref{chap:cld-examples}.  It is quoted, not reproved, in
 \Cref{sec:ch24-comparison}; the part of it that this chapter uses,
 statement~(A), is purely model-theoretic and involves no analysis, so no
 circularity arises.
\end{enumerate}

No result of \Cref{chap:bfj-capacity} or later is used.  In particular the
locality statements proved here are the model-level ones; their bounded
counterparts require quasicontinuity and are proved in
\Cref{chap:bfj-capacity}.  This separation is deliberate and is tabulated in
\Cref{sec:ch24-needs-capacity}.

\subsection{Plan}
\label{sec:ch24-plan}

\Cref{sec:ch24-model-measure} constructs the mixed model measure, proves its
independence of the model, its symmetry, positivity and total mass,
identifies it with the Chambert--Loir measure of
\Cref{sec:metrics-cl-measure}, and proves the model integration-by-parts
formula together with its Hodge-type consequences.
\Cref{sec:ch24-formulae} introduces the mixed energy pairing, proves the
intersection-theoretic formulae that will drive the extension, and verifies
the total-mass statement in mixed form.
\Cref{sec:ch24-bounded} proves the extension theorem, in the weighted
monotone form, together with multilinearity, the negativity of the Dirichlet
form, the Cauchy--Schwarz inequality, the continuous case, and the
integrability statement that no pluripolar set is charged.
\Cref{sec:ch24-locality} proves locality and the comparison and maximum
principles at model level and says exactly which of their extensions need
capacity.  \Cref{sec:ch24-comparison} compares the three model-level
constructions and states the scope of the comparison with the CLD product.
\Cref{sec:ch24-examples} carries six examples to explicit numbers,
\Cref{sec:ch24-scope} is the hypothesis table, and
\Cref{sec:ch24-dictionary} is the chapter dictionary.

\section{The intersection measure on a model}
\label{sec:ch24-model-measure}

\subsection{Vertical intersection numbers}
\label{sec:ch24-vertical-numbers}

A model \(\cX\) has dimension \(n+1\) and is proper over \(R\), not over a
field.  Its special fibre \(\cX_s\), on the other hand, is a projective
scheme over \(k\), so \(0\)-cycles supported on \(\cX_s\) have a degree.
Every number in this chapter is obtained by producing a \(0\)-cycle on
\(\cX_s\) and taking its degree, and the mechanism that produces it is the
presence of at least one \emph{vertical} factor.  We record the resulting
formalism once, because the distinction between a vertical divisor as a
cycle and as a line-bundle class is the single most common source of sign
and mass errors in this subject; the point was already flagged in
\Cref{sec:ch20-ddc-definition} after \eqref{eq:ch20-special-fibre-degree}.

Throughout, \(\Div_0(\cX)_{\mathbf R}\) denotes the space of vertical
\(\mathbf R\)-Cartier divisors on \(\cX\), that is, \(\mathbf R\)-linear
combinations of the components of \(\cX_s\), and \(\cX_s=\sum_Eb_EE\).

\begin{lemma}
\label{lem:ch24-vertical-numbers}
Let \(\cX\) be a model, let \(1\leq m\leq n+1\), let
\(D_1,\ldots,D_m\in\Div_0(\cX)_{\mathbf R}\), and let
\(\alpha_1,\ldots,\alpha_{n+1-m}\in N^1(\cX/R)\).  Then the recipe
\begin{equation}
\label{eq:ch24-vertical-number}
 \bigl(D_1\cdots D_m\cdot\alpha_1\cdots\alpha_{n+1-m}\bigr)_{\cX}
 :=\deg_k\Bigl(
   c_1(D_2)\cdots c_1(D_m)\,\alpha_1\cdots\alpha_{n+1-m}\cap[D_1]
   \Bigr)
\end{equation}
is well defined, and:
\begin{enumerate}
\item\label{itm:ch24-vn:welldef} the \(0\)-cycle in
 \eqref{eq:ch24-vertical-number} is supported on \(\cX_s\), which is
 projective over \(k\), so its degree is a well-defined real number;
\item\label{itm:ch24-vn:sym} the number is symmetric and
 \(\mathbf R\)-multilinear in \(D_1,\ldots,D_m\), and symmetric and
 \(\mathbf R\)-multilinear in \(\alpha_1,\ldots,\alpha_{n+1-m}\);
\item\label{itm:ch24-vn:m1} for \(m=1\) and \(D=\sum_Ea_EE\) it equals
 \(\sum_Ea_E\,\bigl(\alpha_1|_E\cdots\alpha_n|_E\bigr)\);
\item\label{itm:ch24-vn:pullback} it is unchanged under pullback along a
 domination \(\rho:\cY\to\cX\), that is,
 \((\rho^*D_1\cdots\rho^*D_m\cdot\rho^*\alpha_1\cdots)_{\cY}
   =(D_1\cdots D_m\cdot\alpha_1\cdots)_{\cX}\);
\item\label{itm:ch24-vn:fibre} \((\cX_s\cdot\alpha_1\cdots\alpha_n)_{\cX}
 =(\{\alpha_1\}\cdots\{\alpha_n\})_X\), whereas for \(m\geq2\)
 \[
  \bigl(\cX_s\cdot D_2\cdots D_m\cdot\alpha_1\cdots\alpha_{n+1-m}\bigr)_{\cX}
  =0 .
 \]
\end{enumerate}
\end{lemma}

\begin{proof}
\emph{\cref{itm:ch24-vn:welldef} and \cref{itm:ch24-vn:m1}.}  By linearity we
may assume \(D_1=E\) is a prime component of \(\cX_s\).  Then
\([D_1]=[E]\) is an \(n\)-cycle on the projective \(k\)-scheme \(\cX_s\), and
capping with \(n\) first Chern classes produces a class in
\(A_0(E)_{\mathbf R}\), whose degree over \(k\) is defined because \(E\) is
proper over \(k\).  For \(m=1\) the classes \(\alpha_i\) may be restricted to
\(E\) before intersecting, which is \Cref{lem:vertical-intersection-degree},
and gives \cref{itm:ch24-vn:m1}.

\emph{\cref{itm:ch24-vn:sym}.}  Multilinearity in each slot is additivity of
\(c_1\) and of the cycle map \(D\mapsto[D]\).  Symmetry in the
\(\alpha_i\) and among \(D_2,\ldots,D_m\) is commutativity of Chern class
operators on Chow groups, \Cref{prop:appF-mixed-cartier}, applied on the
proper \(k\)-scheme carrying the cycle.  It remains to see that the
distinguished slot may be exchanged with any other, that is, that
\[
 \deg\bigl(c_1(D_2)\gamma\cap[D_1]\bigr)
 =\deg\bigl(c_1(D_1)\gamma\cap[D_2]\bigr),
 \qquad
 \gamma=c_1(D_3)\cdots\alpha_1\cdots .
\]
Both sides are computed on the integral scheme \(\cX\): by the second
assertion of \Cref{lem:vertical-intersection-degree}, \([D_i]\) represents
\(c_1(\cO(D_i))\cap[\cX]\) in the Chow group of \(|D_i|\subseteq\cX_s\), so
both sides equal the degree of
\(c_1(D_1)c_1(D_2)\gamma\cap[\cX]\), computed in
\(A_0(|D_1|\cap|D_2|)_{\mathbf R}\), by the commutativity of Chern class
operators \cite[Corollary~2.4.2]{Ful}.  Since \(|D_1|\cap|D_2|\subseteq\cX_s\)
is proper over \(k\), the degree is defined and the two computations agree.

\emph{\cref{itm:ch24-vn:pullback}.}  This is
\Cref{thm:model-projection-formula} with \(d=1\), applied with the
distinguished slot \(D_1\), together with the Chow projection formula
\Cref{thm:appF-projection} for the remaining factors; note that
\(\rho_*[\rho^*D_1]=[D_1]\) because \(\rho\) is proper and birational.

\emph{\cref{itm:ch24-vn:fibre}.}  For \(m=1\), \cref{itm:ch24-vn:m1} gives
\(\sum_Eb_E(\alpha_1|_E\cdots\alpha_n|_E)\), which is
\eqref{eq:ch20-special-fibre-degree}, that is,
\Cref{cor:special-fibre-relation}: conservation of degree in the flat family
\(\cX\to\Spec R\).  For \(m\geq2\), put \(\cX_s\) in the \emph{second} slot,
which is legitimate by \cref{itm:ch24-vn:sym}: the number then contains the
factor \(c_1(\cO_{\cX}(\cX_s))\), and \(\cO_{\cX}(\cX_s)\simeq\cO_{\cX}\)
because \(\cX_s=\divis(\pi)\) is principal, so \(c_1=0\) and the number
vanishes.
\end{proof}

\begin{warning}{The special fibre has two roles}
\Cref{lem:ch24-vertical-numbers}\,\cref{itm:ch24-vn:fibre} contains both of
them.  As a \emph{cycle}, \(\cX_s=\sum_Eb_EE\) has degree
\((\{\alpha_1\}\cdots\{\alpha_n\})\) against \(n\) numerical classes; as a
\emph{line bundle}, \(\cO(\cX_s)\) is trivial, so its first Chern class
vanishes and any intersection number containing it as a class is zero.  The
two roles are separated by the number of vertical factors: with exactly one
vertical factor the cycle interpretation is forced, with two or more the
class interpretation becomes available and gives zero.  Concretely, the
model function of \(\cX_s\) is the constant \(1\)
(\eqref{eq:ch20-sign}), so integrating it against a model measure returns the
total mass, whereas \(\ddc1=0\) by \eqref{eq:ch20-ddc-constant}.  Confusing
the two produces a total mass of zero; the warning after
\eqref{eq:ch20-special-fibre-degree} makes the same point from the other
side.
\end{warning}

\subsection{The mixed model measure}
\label{sec:ch24-mixed-model}

Choose a projective model \(\mathcal X\) on which closed forms
\(\theta_1,\ldots,\theta_n\) are represented, and write
\[
                    \mathcal X_s=\sum_E b_EE.
\]
The divisorial point associated to \(E\) is denoted \(x_E\), normalized so
that a vertical divisor \(D\) satisfies
\(\varphi_D(x_E)=b_E^{-1}\ord_E(D)\). This normalization accounts for the
factor \(b_E\) below.

\begin{definition}[Mixed model measure]
\index{Monge--Ampere measure!mixed model}
\index{Chambert--Loir measure!mixed}
\label{def:mixed-model-measure}
If \(\theta_{i,\mathcal X}\in N^1(\mathcal X/K^\circ)\) represents
\(\theta_i\), set
\begin{equation}
\label{eq:model-mixed-measure}
 \theta_1\wedge\cdots\wedge\theta_n
 =\sum_E b_E
   \bigl(\theta_{1,\mathcal X}|_E\cdots
         \theta_{n,\mathcal X}|_E\bigr)\delta_{x_E}.
\end{equation}
If \(\varphi_1,\ldots,\varphi_n\) are \(\theta\)-psh model functions, set
\[
 \MA_\theta(\varphi_1,\ldots,\varphi_n)
 :=(\theta+\ddc\varphi_1)\wedge\cdots\wedge
    (\theta+\ddc\varphi_n).
\]
The diagonal measure is
\(\MA_\theta(\varphi):=\MA_\theta(\varphi,\ldots,\varphi)\).
\end{definition}

Three features of \eqref{eq:model-mixed-measure} deserve comment before
anything is proved.  The sum is finite, because a model has finitely many
special-fibre components; the measure is therefore atomic, supported on
divisorial points, and every question about it is a question about finitely
many numbers.  The multiplicity \(b_E\) is not a decoration: it compensates
exactly for the normalization \(x_E=b_E^{-1}\ord_E\) and is what makes
\eqref{eq:model-test-intersection} below come out as a clean intersection
number.  And the right-hand side is written on a chosen model, so that
independence of the choice is the first thing to prove.

If \(f=\varphi_D\) is a model function determined on the same model, the
definition says
\begin{equation}
\label{eq:model-test-intersection}
 \int_{X^{\An}}f\,\theta_1\wedge\cdots\wedge\theta_n
 =\sum_E\ord_E(D)
   \bigl(\theta_{1,\mathcal X}|_E\cdots
         \theta_{n,\mathcal X}|_E\bigr)
 =D\cdot\theta_1\cdots\theta_n.
\end{equation}
The last expression is a vertical intersection number, in the sense of
\eqref{eq:ch24-vertical-number} with \(m=1\).  The computation is one line:
the measure is atomic with atoms at the points \(x_E\), and
\(\varphi_D(x_E)=b_E^{-1}\ord_E(D)\) by
\Cref{prop:metrics-divisor-model-function}, so the two occurrences of
\(b_E\) cancel.  \emph{This cancellation is the whole content of the
normalization conventions}, and \eqref{eq:model-test-intersection} is
identical with \eqref{eq:metrics-cl-pairing}, proved in
\Cref{prop:metrics-cl-model-function-integral}.  Every computation performed
in this chapter is an instance of it.

\begin{proposition}\label{prop:model-ma-basic}
The measure in \Cref{def:mixed-model-measure} is independent of
the common representing model. It is symmetric and multilinear. If all
\(\theta_i\) are semipositive, it is positive and
\begin{equation}
\label{eq:mixed-total-mass}
  \int_{X^{\An}}\theta_1\wedge\cdots\wedge\theta_n
  =(\{\theta_1\}\cdots\{\theta_n\}).
\end{equation}
In particular every mixed \(\theta\)-Monge--Amp\`ere measure has mass \(V\).
\end{proposition}

\begin{proof}
\emph{Step 1: independence of the model.}  Let \(\rho:\mathcal Y\to\mathcal
X\) be a domination, and let \(\mu_{\cX}\), \(\mu_{\cY}\) be the two
measures produced by \eqref{eq:model-mixed-measure} from the determinations
\(\theta_{i,\cX}\) and \(\rho^*\theta_{i,\cX}\).  Both are finite atomic
measures supported on divisorial points, so it is enough to compare
coefficients.  Write \(\cX_s=\sum_Eb_EE\) and \(\cY_s=\sum_Fb_FF\).  We use
three facts.

\emph{(a) Contracted components contribute nothing.}  If \(F\subseteq\cY_s\)
is a component whose image \(\rho(F)\) has dimension \(<n\), then
\(\rho|_F\) factors through a projective \(k\)-scheme of dimension \(<n\), so
the product of the \(n\) pulled-back classes
\(\rho^*\theta_{i,\cX}|_F\) has degree \(0\) on \(F\) by the projection
formula \Cref{thm:appF-projection}.

\emph{(b) Dominating components transform by their degree.}  If \(F\)
dominates a component \(E\subseteq\cX_s\), the induced map \(F\to E\) is
generically finite of some degree \(e_F\), and the projection formula gives
\[
 \bigl(\rho^*\theta_{1,\cX}|_F\cdots\rho^*\theta_{n,\cX}|_F\bigr)
 =e_F\bigl(\theta_{1,\cX}|_E\cdots\theta_{n,\cX}|_E\bigr).
\]
Moreover \(x_F=x_E\) as points of \(X^{\An}\): both are the normalized
valuation \(b^{-1}\ord\) of the same divisorial valuation of \(K(X)\), by
\Cref{prop:component-valuation}.

\emph{(c) The multiplicities add up.}  For each \(E\),
\[
                \sum_{F\mapsto E}e_Fb_F=b_E ,
\]
the sum being over the components of \(\cY_s\) dominating \(E\).  Indeed
\(\cY_s=\rho^*\cX_s\) as Cartier divisors, both being \(\divis(\pi)\), so by
the projection formula for cycles
\(\rho_*[\cY_s]=\rho_*\bigl(c_1(\cO(\rho^*\cX_s))\cap[\cY]\bigr)
 =c_1(\cO(\cX_s))\cap\rho_*[\cY]=[\cX_s]\), using
\(\rho_*[\cY]=[\cX]\) for a proper birational morphism.  Expanding the left
side as \(\sum_Fb_F\rho_*[F]\) and using \(\rho_*[F]=0\) for contracted
\(F\) and \(\rho_*[F]=e_F[E]\) otherwise gives the displayed identity.

Combining (a)--(c),
\[
 \begin{aligned}
 \mu_{\cY}
 &=\sum_Fb_F\bigl(\rho^*\theta_{1,\cX}|_F\cdots\bigr)\delta_{x_F}\\
 &=\sum_E\Bigl(\sum_{F\mapsto E}e_Fb_F\Bigr)
   \bigl(\theta_{1,\cX}|_E\cdots\bigr)\delta_{x_E}\\
 &=\sum_Eb_E\bigl(\theta_{1,\cX}|_E\cdots\bigr)\delta_{x_E}
  =\mu_{\cX}.
 \end{aligned}
\]
This is the computation carried out in the proof of
\Cref{thm:cl-model-properties}, written out once because the multiplicity
bookkeeping is where the normalization is tested.  Two arbitrary
determinations are compared by passing to a common domination, which exists
by \Cref{prop:common-domination}.

It follows in particular that the two measures integrate every model
function to the same number.  Conversely, that weaker statement would have
sufficed: a Radon measure on the compact space \(X^{\An}\) is determined by
its integrals against model functions, since these are dense in
\(C^0(X^{\An})\) by \Cref{prop:metrics-model-function-density} and
\Cref{cor:riesz-uniqueness} applies.  We shall use this remark repeatedly,
in the form \Cref{lem:ch24-dense-tests}.

\emph{Step 2: symmetry and multilinearity.}  Both are visible in
\eqref{eq:model-mixed-measure}: the coefficient of \(\delta_{x_E}\) is the
mixed degree \((\theta_{1}|_E\cdots\theta_{n}|_E)\) on the projective
\(k\)-variety \(E\), which is symmetric and multilinear by
\Cref{prop:appF-mixed-cartier}.  Multilinearity is understood in the
following sense: if \(\theta_1=t'\theta_1'+t''\theta_1''\) with
\(t',t''\in\mathbf R\) and all three forms determined on a common model,
then the corresponding measures satisfy the same relation.  In particular,
for \(\theta\)-psh model functions, the map
\(\varphi_1\mapsto\MA_\theta(\varphi_1,\ldots,\varphi_n)\) is \emph{affine},
because \(\varphi_1\mapsto\theta+\ddc\varphi_1\) is.

\emph{Step 3: positivity.}  Suppose every \(\theta_i\) is semipositive.  By
\Cref{prop:ch20-semipositive-well-defined} we may choose the common
determination so that each \(\theta_{i,\cX}\) is relatively nef; its
restriction \(\theta_{i,\cX}|_E\) to a component of \(\cX_s\) is then nef on
the projective \(k\)-variety \(E\)
(\Cref{lem:ch20-special-fibre-test}\,\cref{itm:ch20-sft:nef}).  A product of
\(n\) nef classes on an \(n\)-dimensional projective variety has nonnegative
degree, again by \Cref{prop:appF-mixed-cartier}.  Since \(b_E>0\), every
coefficient in \eqref{eq:model-mixed-measure} is nonnegative, which is
\eqref{eq:ch20-mixed-component-positive}.

\emph{Step 4: total mass.}  Take \(f=1=\varphi_{\cX_s}\) in
\eqref{eq:model-test-intersection}, using \eqref{eq:ch20-sign}.  The
left-hand side is the total mass; the right-hand side is
\((\cX_s\cdot\theta_1\cdots\theta_n)_{\cX}\), which is
\((\{\theta_1\}\cdots\{\theta_n\})\) by
\Cref{lem:ch24-vertical-numbers}\,\cref{itm:ch24-vn:fibre}.  This proves
\eqref{eq:mixed-total-mass}.  Taking all \(\theta_i=\theta+\ddc\varphi_i\)
and using \(\{\ddc\varphi_i\}=0\) (\Cref{thm:ch20-ddc-exact}) gives
\(\{\theta_i\}=\{\theta\}\) for every \(i\), so the mass is
\((\{\theta\}^n)=V\).
\end{proof}

\begin{corollary}
\label{cor:ch24-cl-identification}
Let \(L_1,\ldots,L_n\) be line bundles on \(X\) equipped with semipositive
model metrics \(\|\cdot\|_i\), determined by model line bundles \(\cL_i\)
extending \(L_i^{\otimes m_i}\) on a common model \(\cX\), and let
\(\theta_i\in\mathcal Z^{1,1}(X)\) be the closed model form determined by
\(m_i^{-1}c_1(\cL_i)\).  Then
\[
 \theta_1\wedge\cdots\wedge\theta_n
 =c_1(L_1,\|\cdot\|_1)\wedge\cdots\wedge c_1(L_n,\|\cdot\|_n),
\]
the mixed Chambert--Loir measure of \Cref{sec:metrics-cl-measure}.  In
particular \Cref{def:mixed-model-measure} introduces no new object: it
records the Chambert--Loir measure in the language of closed model forms.
\end{corollary}

\begin{proof}
Both sides are given by \eqref{eq:model-mixed-measure} on the model \(\cX\),
after dividing by \(m_1\cdots m_n\); the two total masses agree by
\Cref{thm:cl-model-properties} and \eqref{eq:mixed-total-mass}, and the
independence statements match.  Semipositivity of the metric \(\|\cdot\|_i\)
means relative nefness of \(\cL_i\)
(\Cref{sec:metrics-semipositive}), which is semipositivity of \(\theta_i\)
in the sense of \Cref{def:ch20-semipositive}, by
\Cref{lem:metrics-nef-independence}.
\end{proof}

\begin{proposition}
\label{prop:ch24-model-support}
Let \(\theta_1,\ldots,\theta_n\) be semipositive closed model forms
determined on a model \(\cX\), and let
\(\mu=\theta_1\wedge\cdots\wedge\theta_n\).
\begin{enumerate}
\item\label{itm:ch24-ms:atomic} \(\mu\) is a finite positive combination of
 Dirac masses at divisorial points of \(\cX\); in particular
 \(\Supp\mu\subseteq\Delta_{\cX}\), the dual complex of \(\cX\), and
 \(\mu\) is supported on the set of vertices of \(\Delta_{\cX}\) when
 \(\cX\) is SNC.
\item\label{itm:ch24-ms:refine} If \(\rho:\cY\to\cX\) is a domination, the
 atoms of \(\mu\) computed on \(\cY\) sit at the same points of
 \(X^{\An}\); no mass appears at a divisorial point of \(\cY\) that is not
 already a divisorial point of \(\cX\).
\item\label{itm:ch24-ms:pluripolar} \(\mu\) gives no mass to
 \(\{u=-\infty\}\) for any \(u\in\PSH(X,\theta)\) not identically
 \(-\infty\).
\end{enumerate}
\end{proposition}

\begin{proof}
\cref{itm:ch24-ms:atomic} is \eqref{eq:model-mixed-measure} together with
Step~3 of the previous proof; the divisorial points of an SNC model are the
vertices of \(\Delta_{\cX}\) by \Cref{lem:dual-complex-simplicial}.
\cref{itm:ch24-ms:refine} is Step~1 of the previous proof: a component of
\(\cY_s\) contracted by \(\rho\) to a subset of dimension \(<n\) receives
coefficient \(0\), and a component dominating \(E\) contributes at
\(x_E\).  For \cref{itm:ch24-ms:pluripolar}, note that by
\Cref{prop:ch21-first-consequences}\,\cref{itm:ch21-fc:quasimonomial} a
\(\theta\)-psh function that is not identically \(-\infty\) is finite at every
quasimonomial point, in particular at every divisorial point, so the polar
set misses \(\Supp\mu\).
\end{proof}

\begin{example}[Good reduction]
\label{ex:ch24-good-reduction}
If \(\cX_s\) is irreducible and reduced---the case of good reduction---then
\(\cX_s=E\) with \(b_E=1\), the dual complex is the single point \(x_E\),
and \eqref{eq:model-mixed-measure} reduces to
\[
 \theta_1\wedge\cdots\wedge\theta_n=(\{\theta_1\}\cdots\{\theta_n\})\,
 \delta_{x_E},
\]
because the restriction of \(\theta_{i,\cX}\) to \(E\) has the same degrees
as \(\{\theta_i\}\) on \(X\) by
\Cref{cor:special-fibre-relation}.  For \(X=\mathbf P^n_K\),
\(\cX=\mathbf P^n_R\) and \(\theta_i=c_1(\cO(1))\) this is
\(\delta_{x_{\cX_s}}\), of mass \(1\); for \(n=1\) it is the Dirac mass at
the Gauss point, \eqref{eq:p1-standard-measure}.  No mixed information is
present in this case: every mixed measure is a multiple of the same Dirac
mass, and the only content is the mass.
\end{example}

\subsection{Multilinearity and polarization}
\label{sec:ch24-polarization}

Multilinearity of \eqref{eq:model-mixed-measure} has a useful converse: the
mixed measures are determined by the diagonal ones.  This is not merely
formal bookkeeping.  Once the extension theorem of
\Cref{sec:ch24-bounded} is available, the polarization identity below shows
that any statement proved for diagonal products---for instance a comparison
with another theory---propagates to mixed products without further work,
provided the class of potentials involved is stable under averages.  It is
in this form that the toric computation of
\Cref{ex:ch24-toric-surface-mixed} is checked.

\begin{proposition}
\label{prop:ch24-polarization}
Let \(\varphi_1,\ldots,\varphi_n\) be \(\theta\)-psh model functions.  For a
nonempty subset \(S\subseteq\{1,\ldots,n\}\) put
\(\varphi_S:=|S|^{-1}\sum_{i\in S}\varphi_i\), again a \(\theta\)-psh model
function by \Cref{thm:ch21-stability}\,\cref{itm:ch21-stab:sum}.  Then
\begin{equation}
\label{eq:ch24-polarization}
 n!\;\MA_\theta(\varphi_1,\ldots,\varphi_n)
 =\sum_{\varnothing\neq S\subseteq\{1,\ldots,n\}}
   (-1)^{\,n-|S|}\,|S|^{\,n}\;\MA_\theta(\varphi_S).
\end{equation}
\end{proposition}

\begin{proof}
Write \(\eta_i=\theta+\ddc\varphi_i\) and let
\(M(\eta_1,\ldots,\eta_n)\) denote the symmetric \(n\)-linear
measure-valued form of \eqref{eq:model-mixed-measure}, so that
\(\MA_\theta(\varphi_1,\ldots,\varphi_n)=M(\eta_1,\ldots,\eta_n)\) and
\(Q(\eta):=M(\eta,\ldots,\eta)\).  The classical polarization identity for a
symmetric \(n\)-linear form on a real vector space reads
\[
 n!\,M(\eta_1,\ldots,\eta_n)
 =\sum_{S\subseteq\{1,\ldots,n\}}(-1)^{n-|S|}
   Q\Bigl(\sum_{i\in S}\eta_i\Bigr),
\]
the term \(S=\varnothing\) being \(Q(0)=0\); it is proved by expanding
\(Q(\sum_{i\in S}\eta_i)\) multilinearly and observing that a monomial
\(M(\eta_{j_1},\ldots,\eta_{j_n})\) whose index set omits some
\(i_0\in\{1,\ldots,n\}\) contributes with total coefficient
\(\sum_{S}(-1)^{n-|S|}\), the sum being over subsets \(S\) containing the
index set, which vanishes by pairing \(S\) with \(S\cup\{i_0\}\).  Now for
\(S\neq\varnothing\),
\[
 \sum_{i\in S}\eta_i=|S|\Bigl(\theta+\ddc\varphi_S\Bigr),
\]
because \(\ddc\) is linear, and \(Q\) is homogeneous of degree \(n\); hence
\(Q(\sum_{i\in S}\eta_i)=|S|^n\MA_\theta(\varphi_S)\).  Substituting gives
\eqref{eq:ch24-polarization}.
\end{proof}

\begin{corollary}\label{cor:ch24-mixed-from-diagonal}
Two symmetric multilinear measure-valued products defined on \(n\)-tuples of
\(\theta\)-psh model functions agree if and only if their diagonals agree.
The same holds on \(n\)-tuples of bounded \(\theta\)-psh functions, once
\Cref{thm:bounded-ma-extension} has provided the extension, since the class
of bounded \(\theta\)-psh functions is stable under averages.
\end{corollary}

\begin{proof}
Necessity is trivial and sufficiency is \eqref{eq:ch24-polarization}, whose
right-hand side involves only diagonal products.  For the bounded case,
observe that \(\varphi_S\) is bounded \(\theta\)-psh when the \(\varphi_i\)
are, by \Cref{thm:ch21-stability}\,\cref{itm:ch21-stab:sum}, and that
\eqref{eq:ch24-polarization} continues to hold there because both sides are
continuous along decreasing nets by
\Cref{thm:bounded-ma-extension}\,(2), and a decreasing net of model
approximations of the \(\varphi_i\) induces one of \(\varphi_S\).
\end{proof}

\subsection{Integration by parts and the Hodge inequalities}
\label{sec:ch24-model-ibp}

By multilinearity, \eqref{eq:model-mixed-measure} makes sense for arbitrary
closed model forms and produces a signed atomic measure.  We use this only
in the following combination.  If \(f\) is a model function and
\(\eta_1,\ldots,\eta_{n-1}\) are closed model forms determined on a common
model, we write
\[
 \ddc f\wedge\eta_1\wedge\cdots\wedge\eta_{n-1}
 :=(\theta+\ddc f)\wedge\eta_1\wedge\cdots\wedge\eta_{n-1}
   -\theta\wedge\eta_1\wedge\cdots\wedge\eta_{n-1},
\]
a signed atomic measure of \emph{total mass zero}, by
\eqref{eq:mixed-total-mass} and \(\{\ddc f\}=0\).  When \(f\) is
\(\theta\)-psh and the \(\eta_i\) are semipositive, it is a difference of
two positive measures of mass \((\{\theta\}\{\eta_1\}\cdots\{\eta_{n-1}\})\)
each, so its total variation is at most twice that number.  This crude bound
is used repeatedly below and is the reason the whole construction is stable.

\begin{proposition}\label{prop:model-integration-by-parts}
For model functions \(f,g\) and closed model forms
\(\eta_1,\ldots,\eta_{n-1}\),
\begin{equation}
\label{eq:model-ibp}
 \int f\,\ddc g\wedge\eta_1\wedge\cdots\wedge\eta_{n-1}
 =\int g\,\ddc f\wedge\eta_1\wedge\cdots\wedge\eta_{n-1}.
\end{equation}
If the \(\eta_i\) are semipositive, then
\begin{equation}
\label{eq:model-dirichlet-negativity}
 \int f\,\ddc f\wedge\eta_1\wedge\cdots\wedge\eta_{n-1}\leq0.
\end{equation}
\end{proposition}

\begin{proof}
\emph{Step 1: both sides are the same intersection number.}  Choose a model
\(\cX\) on which \(f=\varphi_D\), \(g=\varphi_{D'}\) and all the \(\eta_j\)
are determined; this is possible by \Cref{prop:common-domination} and
\Cref{cor:ch20-limit-structure}.  By definition of the signed measure and
\eqref{eq:model-test-intersection},
\[
 \begin{aligned}
 \int f\,\ddc g\wedge\eta_1\wedge\cdots\wedge\eta_{n-1}
 &=\bigl(D\cdot(\theta+\ddc g)\cdot\eta_1\cdots\eta_{n-1}\bigr)\\
 &\quad-\bigl(D\cdot\theta\cdot\eta_1\cdots\eta_{n-1}\bigr)\\
 &=\bigl(D\cdot D'\cdot\eta_1\cdots\eta_{n-1}\bigr),
 \end{aligned}
\]
where the last step uses multilinearity in the second slot and
\(\ddc g=[D']\) in \(N^1(\cX/R)\).  All three numbers are vertical
intersection numbers in the sense of \eqref{eq:ch24-vertical-number}, with
\(m=2\) in the last one.  By
\Cref{lem:ch24-vertical-numbers}\,\cref{itm:ch24-vn:sym} the number
\((D\cdot D'\cdot\eta_1\cdots\eta_{n-1})\) is symmetric in \(D\) and
\(D'\), so exchanging \(f\) and \(g\) gives the same value.  This is
\eqref{eq:model-ibp}.  Independence of the chosen model is
\Cref{lem:ch24-vertical-numbers}\,\cref{itm:ch24-vn:pullback}.

\emph{Step 2: negativity.}  Taking \(g=f\) in Step~1,
\[
 \int f\,\ddc f\wedge\eta_1\wedge\cdots\wedge\eta_{n-1}
 =\bigl(D^2\cdot\eta_1\cdots\eta_{n-1}\bigr),
\]
and if the \(\eta_j\) are semipositive, determined by relatively nef classes
on \(\cX\), this is \(\leq0\) by the mixed relative Hodge index inequality
\Cref{thm:ch20-mixed-hodge}, that is,
\eqref{eq:ch20-mixed-hodge}.  Note that \Cref{thm:ch20-mixed-hodge} is
stated on a \emph{regular} model; by \Cref{prop:snc-cofinal} the SNC, hence
regular, models are cofinal, so after passing to a domination---which
changes neither side, by Step~1---the hypothesis is met.
\end{proof}

The sign in \eqref{eq:model-dirichlet-negativity} is the algebraic origin of
concavity of the Monge--Amp\`ere energy in \Cref{chap:bfj-energy}, and of
every Cauchy--Schwarz estimate used in \Cref{chap:bfj-capacity}.  We record
the polarized form.

\begin{proposition}
\label{prop:ch24-model-cauchy-schwarz}
Let \(\eta_1,\ldots,\eta_{n-1}\) be semipositive closed model forms and let
\(f,g\) be model functions.  Then
\begin{equation}
\label{eq:ch24-model-cauchy-schwarz}
 \left|\int f\,\ddc g\wedge\eta_1\wedge\cdots\wedge\eta_{n-1}\right|^{2}
 \leq
 \left(\int f\,\ddc f\wedge\eta_\bullet\right)
 \left(\int g\,\ddc g\wedge\eta_\bullet\right),
\end{equation}
where \(\eta_\bullet=\eta_1\wedge\cdots\wedge\eta_{n-1}\).  Both factors on
the right are nonpositive, so their product is nonnegative.
\end{proposition}

\begin{proof}
By Step~1 of the previous proof the three quantities are
\((D\cdot D'\cdot\eta_\bullet)\), \((D^2\cdot\eta_\bullet)\) and
\((D'^2\cdot\eta_\bullet)\) for vertical divisors \(D,D'\) on a common
regular model.  The assertion is then
\eqref{eq:ch20-mixed-cauchy-schwarz}.
\end{proof}

\begin{remark}[Sign audit]
\label{rem:ch24-sign-audit}
Three signs have now been fixed and are mutually consistent.  A vertical
effective divisor gives a nonnegative model function
(\eqref{eq:ch20-sign}); a convex corner of a model function on a dual
complex gives \emph{positive} mass, since by
\Cref{prop:ch20-curve-laplacian} the Laplacian is the residue-degree-weighted
sum of outgoing slopes; and the Dirichlet form
\(f\mapsto\int f\,\ddc f\wedge\eta_\bullet\) is \emph{nonpositive}.  The
three are checked against the unit tests of \Cref{sec:appJ-unit-tests} as
follows: with \(f=1\) the first gives \(\varphi_{\cX_s}=1\), the second
gives \(\ddc1=0\), and the third gives \(0\), consistently.  On a curve
with the residue-rational-node two-component fibre of
\eqref{eq:ch20-two-component-matrix}, a model function with vertex
values \(u_0,u_1\) has
\(\int f\,\ddc f=-(u_0-u_1)^2\leq0\), with equality exactly for constants;
this is the case \(n=1\) of \eqref{eq:model-dirichlet-negativity} and is the
smallest complete check available.
\end{remark}

\section{Intersection-theoretic formulae}
\label{sec:ch24-formulae}

The measures of \Cref{sec:ch24-model-measure} are determined by the numbers
\(\int f\,\theta_1\wedge\cdots\wedge\theta_n\).  This section organizes those
numbers into a single symmetric object, the mixed energy pairing, and
proves the identities that will carry the whole construction to bounded
potentials in \Cref{sec:ch24-bounded}.  Nothing here is analysis: every
statement is an identity between intersection numbers on one model, and
every proof is an expansion by multilinearity.

From here to the end of \Cref{sec:ch24-bounded} we use the normalization
recorded in \Cref{sec:ch24-hypotheses}: \emph{\(\theta\) is semipositive}.
By \Cref{prop:ch20-positive-representative} this may always be arranged by
replacing \(\theta\) with \(\theta+\ddc\psi\) and each potential
\(\varphi\) with \(\varphi-\psi\), and by
\Cref{cor:ch21-reference-independence} no statement is affected.  Its only
role is to make the constant functions \(\theta\)-psh
(\Cref{prop:ch21-constants}), so that in particular the entry \(0\) and the
entries \(\pm M\) are legitimate arguments of the pairing below.

\subsection{The mixed energy pairing}
\label{sec:ch24-energy-pairing}

\begin{definition}[Mixed energy pairing]
\index{energy!mixed pairing}
\label{def:ch24-mixed-energy}
Let \(\varphi_0,\ldots,\varphi_n\) be model functions, determined on a common
model \(\cX\) by vertical \(\mathbf R\)-Cartier divisors
\(D_0,\ldots,D_n\), and let \(\theta\) be determined on \(\cX\) as well.
Put
\begin{equation}
\label{eq:ch24-mixed-energy}
 \mathfrak E(\varphi_0,\ldots,\varphi_n)
 :=\sum_{\varnothing\neq S\subseteq\{0,\ldots,n\}}
    \Bigl(\prod_{i\in S}D_i\cdot\theta^{\,n+1-|S|}\Bigr)_{\cX},
\end{equation}
each summand being a vertical intersection number
\eqref{eq:ch24-vertical-number} with \(m=|S|\geq1\).
\end{definition}

Informally, \(\mathfrak E(\varphi_0,\ldots,\varphi_n)\) is
\(\bigl((\theta+\ddc\varphi_0)\cdots(\theta+\ddc\varphi_n)\bigr)
 -(\theta^{n+1})\); we do not write it that way, because
\((\theta^{n+1})\) is \emph{not} defined---there is no degree map on
\(0\)-cycles of \(\cX\) with no vertical factor---and because the difference
is exactly what is well defined.  This is the same distinction as in the
warning after \Cref{lem:ch24-vertical-numbers}, and it is the reason the
mixed energy pairing, and not a naive top self-intersection, is the object
that extends.

\begin{proposition}
\label{prop:ch24-energy-basic}
The quantity \eqref{eq:ch24-mixed-energy} is a real number depending only on
the functions \(\varphi_0,\ldots,\varphi_n\) and on \(\theta\), and:
\begin{enumerate}
\item\label{itm:ch24-eb:model} it is independent of the common model on
 which the data are determined;
\item\label{itm:ch24-eb:sym} it is symmetric in \(\varphi_0,\ldots,\varphi_n\);
\item\label{itm:ch24-eb:slot} for every \(j\),
 \begin{equation}
 \label{eq:ch24-energy-slot}
  \begin{aligned}
  \mathfrak E(\varphi_0,\ldots,\varphi_n)
  &=\mathfrak E\bigl(\varphi_0,\ldots,\varphi_{j-1},0,
      \varphi_{j+1},\ldots,\varphi_n\bigr)\\
  &\quad+\int_{X^{\An}}\varphi_j\;
     \MA_\theta\bigl(\varphi_0,\ldots,\widehat{\varphi_j},
       \ldots,\varphi_n\bigr),
  \end{aligned}
 \end{equation}
 the hat denoting omission;
\item\label{itm:ch24-eb:const} for constants \(c_0,\ldots,c_n\),
 \begin{equation}
 \label{eq:ch24-energy-constants}
  \mathfrak E(\varphi_0+c_0,\ldots,\varphi_n+c_n)
  =\mathfrak E(\varphi_0,\ldots,\varphi_n)+\Bigl(\sum_jc_j\Bigr)V,
  \qquad \mathfrak E(0,\ldots,0)=0;
 \end{equation}
\item\label{itm:ch24-eb:affine} it is affine in each variable separately;
\item\label{itm:ch24-eb:mono} if all the \(\varphi_i\) are \(\theta\)-psh,
 then \(\mathfrak E\) is nondecreasing in each variable: if
 \(\varphi_j\leq\psi_j\) with \(\psi_j\) also \(\theta\)-psh, then
 \(\mathfrak E(\ldots,\varphi_j,\ldots)\leq\mathfrak E(\ldots,\psi_j,\ldots)\),
 and the increment is
 \(\int(\psi_j-\varphi_j)\MA_\theta(\varphi_0,\ldots,\widehat{\varphi_j},
 \ldots,\varphi_n)\);
\item\label{itm:ch24-eb:telescope} \(\mathfrak E\) is given by the
 telescoping formula
 \begin{equation}
 \label{eq:ch24-energy-telescoping}
  \mathfrak E(\varphi_0,\ldots,\varphi_n)
  =\sum_{j=0}^{n}\int_{X^{\An}}\varphi_j\;
    \MA_\theta\bigl(\varphi_0,\ldots,\varphi_{j-1},
      \underbrace{0,\ldots,0}_{n-j}\bigr).
 \end{equation}
\end{enumerate}
\end{proposition}

\begin{proof}
\emph{Well-definedness and \cref{itm:ch24-eb:model}.}  On a fixed normal model,
a vertical \(\mathbf R\)-divisor is determined by its model function, since
\(\ord_E(D)=b_E\varphi_D(x_E)\); so the \(D_i\) are determined by the
\(\varphi_i\).  Independence of the model is
\Cref{lem:ch24-vertical-numbers}\,\cref{itm:ch24-vn:pullback} applied to each
summand, together with \(\rho^*\theta_{\cX}=\theta_{\cY}\) and
\(\varphi_{\rho^*D}=\varphi_D\).

\emph{\cref{itm:ch24-eb:sym}.}  The index set of the sum
\eqref{eq:ch24-mixed-energy} is symmetric under permutations of
\(\{0,\ldots,n\}\), and each summand is symmetric in the divisors it
contains by \Cref{lem:ch24-vertical-numbers}\,\cref{itm:ch24-vn:sym}.

\emph{\cref{itm:ch24-eb:slot}.}  Split the sum according to whether
\(j\in S\).  The subsets with \(j\notin S\) reproduce
\(\mathfrak E(\varphi_0,\ldots,0,\ldots,\varphi_n)\), since replacing
\(\varphi_j\) by \(0\) replaces \(D_j\) by \(0\) and kills exactly the
subsets containing \(j\).  The subsets with \(j\in S\) sum to
\[
 \begin{aligned}
 &\sum_{T\subseteq\{0,\ldots,n\}\setminus\{j\}}
  \Bigl(D_j\cdot\prod_{i\in T}D_i\cdot\theta^{\,n-|T|}\Bigr)\\
 &\qquad=\Bigl(D_j\cdot\prod_{i\neq j}(\theta+\ddc\varphi_i)\Bigr)\\
 &\qquad=\int\varphi_j\,\MA_\theta
   (\varphi_0,\ldots,\widehat{\varphi_j},\ldots,\varphi_n),
 \end{aligned}
\]
by multilinearity in the last \(n\) slots
(\Cref{lem:ch24-vertical-numbers}\,\cref{itm:ch24-vn:sym}) and then
\eqref{eq:model-test-intersection}.

\emph{\cref{itm:ch24-eb:const}.}  Adding \(c_j\) to \(\varphi_j\) replaces
\(D_j\) by \(D_j+c_j\cX_s\), by \eqref{eq:ch20-sign}.  Expanding
\eqref{eq:ch24-mixed-energy}, a summand indexed by \(S\) acquires the extra
terms
\(c_j(\cX_s\cdot\prod_{i\in S\setminus\{j\}}D_i\cdot\theta^{n+1-|S|})\) for
\(j\in S\), and these vanish unless \(S=\{j\}\) by
\Cref{lem:ch24-vertical-numbers}\,\cref{itm:ch24-vn:fibre}; products of two
or more factors \(\cX_s\) vanish for the same reason.  The surviving terms
are \(c_j(\cX_s\cdot\theta^n)=c_jV\).  Finally
\(\mathfrak E(0,\ldots,0)=0\) because all \(D_i\) vanish.

\emph{\cref{itm:ch24-eb:affine}.}  Immediate from
\eqref{eq:ch24-energy-slot}, whose right-hand side is affine in
\(\varphi_j\): the first term does not involve \(\varphi_j\) and the second
is linear in it.

\emph{\cref{itm:ch24-eb:mono}.}  By \eqref{eq:ch24-energy-slot} applied
twice,
\[
 \mathfrak E(\ldots,\psi_j,\ldots)-\mathfrak E(\ldots,\varphi_j,\ldots)
 =\int(\psi_j-\varphi_j)\,
   \MA_\theta(\varphi_0,\ldots,\widehat{\varphi_j},\ldots,\varphi_n),
\]
and the measure is positive by \Cref{prop:model-ma-basic}, the remaining
\(\varphi_i\) being \(\theta\)-psh.

\emph{\cref{itm:ch24-eb:telescope}.}  Apply \eqref{eq:ch24-energy-slot} with
\(j=n\), then to the resulting term with \(j=n-1\), and so on; after \(n+1\)
steps the remaining pairing is \(\mathfrak E(0,\ldots,0)=0\).
\end{proof}

\begin{remark}[Relation to the Monge--Amp\`ere energy]
\label{rem:ch24-energy-vs-E}
For \(\varphi_0=\cdots=\varphi_n=\varphi\), formula
\eqref{eq:ch24-energy-telescoping} reads
\(\mathfrak E(\varphi,\ldots,\varphi)
 =\sum_{j=0}^n\int\varphi\,
   \MA_\theta(\varphi^{(j)},0^{(n-j)})\),
which is \((n+1)\) times the Monge--Amp\`ere energy of
\Cref{def:model-energy}---there written \(E_\omega\), the reference form
being called \(\omega\)---in the normalization
\eqref{eq:model-energy}, that is, \eqref{eq:appJ-energy}.  We do not use
that identification in this chapter,
and we do not use any property of \(E_\theta\); the pairing
\(\mathfrak E\) is introduced here purely as the monotone scalar through
which the extension of \(\MA_\theta\) is routed.  \Cref{chap:bfj-energy}
develops the energy proper, including the cocycle formula
\eqref{eq:energy-cocycle} and the concavity that follows from
\eqref{eq:model-dirichlet-negativity}.
\end{remark}

\subsection{Masses and mixed intersection numbers}
\label{sec:ch24-masses}

\begin{theorem}
\label{thm:ch24-mass}
Let \(\theta_1,\ldots,\theta_n\) be semipositive closed model forms.  Then
\begin{equation}
\label{eq:ch24-mass}
 \bigl\|\theta_1\wedge\cdots\wedge\theta_n\bigr\|
 =\int_{X^{\An}}\theta_1\wedge\cdots\wedge\theta_n
 =\bigl(\{\theta_1\}\cdots\{\theta_n\}\bigr)_X ,
\end{equation}
the mixed intersection number of the generic numerical classes on \(X\).  In
particular:
\begin{enumerate}
\item\label{itm:ch24-mass:psh} if \(\varphi_1,\ldots,\varphi_n\) are
 \(\theta\)-psh model functions, then
 \(\MA_\theta(\varphi_1,\ldots,\varphi_n)\) has total mass \(V\),
 independently of the potentials;
\item\label{itm:ch24-mass:pos} the number \eqref{eq:ch24-mass} is positive
 as soon as the classes \(\{\theta_i\}\) are ample, and it is the sum of the
 nonnegative componentwise degrees \(b_E(\theta_1|_E\cdots\theta_n|_E)\);
\item\label{itm:ch24-mass:mixed} the assignment
 \((\theta_1,\ldots,\theta_n)\mapsto\|\theta_1\wedge\cdots\wedge\theta_n\|\)
 is the restriction to semipositive forms of a symmetric multilinear form on
 \(\mathcal Z^{1,1}(X)\), namely the pullback of the mixed intersection form
 of \(N^1(X)\) along \(\theta\mapsto\{\theta\}\).
\end{enumerate}
\end{theorem}

\begin{proof}
The first equality is the definition of total mass and the second is
\eqref{eq:mixed-total-mass}, proved in Step~4 of
\Cref{prop:model-ma-basic}.  For \cref{itm:ch24-mass:psh}, the generic class
of \(\theta+\ddc\varphi_i\) is \(\{\theta\}\), because
\(\{\ddc\psi\}=0\) for every model function \(\psi\) by exactness of
\eqref{eq:ch20-ddc-exact} at \(\mathcal Z^{1,1}(X)\); hence the mixed number
is \((\{\theta\}^n)=V\).  Statement \cref{itm:ch24-mass:pos} is
\eqref{eq:ch20-mixed-component-positive} together with positivity of the top
self-intersection of an ample class.  Statement \cref{itm:ch24-mass:mixed} is
the observation that the right-hand side of \eqref{eq:ch24-mass} depends only
on the generic classes and is symmetric multilinear there, by
\Cref{prop:appF-mixed-cartier} applied on \(X\).
\end{proof}

\begin{corollary}
\label{cor:ch24-mass-independent}
For \(\theta\)-psh model functions the total mass of
\(\MA_\theta(\varphi_1,\ldots,\varphi_n)\) is \(V\) and does not depend on
\(\varphi_1,\ldots,\varphi_n\).  Consequently, for any two \(n\)-tuples of
\(\theta\)-psh model functions the signed measure
\(\MA_\theta(\varphi_\bullet)-\MA_\theta(\psi_\bullet)\) has total mass
\(0\), and its total variation is at most \(2V\).
\end{corollary}

\begin{proof}
The first assertion is
\Cref{thm:ch24-mass}\,\cref{itm:ch24-mass:psh}; the second follows since both
measures are positive of mass \(V\).
\end{proof}

The next statement is the model-level shadow of the entire extension theorem
and is worth isolating: it says that the map from potentials to measures is
uniformly continuous for the sup norm on potentials and the weak topology on
measures, with an explicit modulus.

\begin{proposition}
\label{prop:ch24-model-uniform-estimate}
Let \(\varphi_1,\ldots,\varphi_n\) and \(\psi_1,\ldots,\psi_n\) be
\(\theta\)-psh model functions, and let \(h\) be a \(\theta\)-psh model
function.  Then
\begin{equation}
\label{eq:ch24-model-uniform-estimate}
 \left|\int h\,\MA_\theta(\varphi_1,\ldots,\varphi_n)
      -\int h\,\MA_\theta(\psi_1,\ldots,\psi_n)\right|
 \leq 2V\sum_{i=1}^{n}\|\varphi_i-\psi_i\|_\infty .
\end{equation}
The same holds with \(h\) replaced by any difference of two \(\theta\)-psh
model functions, with the right-hand side doubled.
\end{proposition}

\begin{proof}
Telescoping in the \(n\) slots, it suffices to bound one difference in which
a single slot changes.  So let \(\varphi,\psi\) be \(\theta\)-psh model
functions and let \(\Theta\) denote a mixed product of \(n-1\) semipositive
closed model forms with generic classes \(\{\theta\}\).  Then
\[
 \int h\bigl[(\theta+\ddc\varphi)\wedge\Theta
   -(\theta+\ddc\psi)\wedge\Theta\bigr]
 =\int h\,\ddc(\varphi-\psi)\wedge\Theta
 =\int(\varphi-\psi)\,\ddc h\wedge\Theta,
\]
by \eqref{eq:model-ibp}.  The signed measure \(\ddc h\wedge\Theta\) is the
difference of the two positive measures \((\theta+\ddc h)\wedge\Theta\) and
\(\theta\wedge\Theta\), each of mass \(V\) by
\eqref{eq:ch24-mass}, so its total variation is at most \(2V\).  Hence the
displayed quantity is at most \(2V\|\varphi-\psi\|_\infty\) in absolute
value.  Summing the \(n\) telescoping terms gives
\eqref{eq:ch24-model-uniform-estimate}.  The last sentence follows by
linearity in \(h\).
\end{proof}

\section{Monotone extension to bounded psh potentials}
\label{sec:ch24-bounded}

A bounded psh function need not live on one skeleton, so
\eqref{eq:model-mixed-measure} cannot simply be applied to it. Decreasing
regularization supplies model approximants, but defining a measure by taking
a limit is legitimate only after proving independence of every choice.

\subsection{Why a limit needs an argument}
\label{sec:ch24-why-argument}

Let \(\varphi\) be a bounded \(\theta\)-psh function.  By
\Cref{thm:bfj-decreasing-regularization} there is a decreasing net
\(\varphi^\alpha\) of \(\theta\)-psh model functions with
\(\varphi^\alpha\searrow\varphi\) pointwise.  Each
\(\MA_\theta(\varphi^\alpha)\) is a positive measure of mass \(V\), and the
space of such measures is weakly compact
(\Cref{prop:measure-compactness}).  So convergent subnets exist.  That is
all that abstract nonsense supplies, and it is not enough for three separate
reasons.

\emph{First, subnets are not nets.}  A limit obtained along a subnet may
depend on the subnet.  Nothing so far forbids two subnets of
\(\bigl(\MA_\theta(\varphi^\alpha)\bigr)_\alpha\) from converging to
different measures.

\emph{Second, different regularizing nets must give the same answer.}  Even
if the limit along one net exists, a second net \(\psi^\beta\searrow\varphi\)
is not comparable to the first index by index; the two index sets are
unrelated directed sets, and there is no monotonicity to exploit between
them.  The comparison has to be made on \(X^{\An}\), and this is where
compactness of \(X^{\An}\) and Dini's theorem for nets,
\Cref{prop:dini-nets}, enter; see \Cref{lem:ch24-sandwich}.

\emph{Third, sequences do not suffice.}  One might hope to avoid nets by
extracting a decreasing \emph{sequence} of model functions converging to
\(\varphi\).  This is possible when \(\varphi\) is continuous
(\Cref{cor:ch22-continuous-decreasing-sequence}) but not in general:
\Cref{ex:ch22-countable-blind-spot} produces a \(\theta\)-psh function that
no countable family of model functions can approximate from above, and
\Cref{lem:ch21-no-cofinal-sequence} explains the structural reason.  Every
monotone statement in this chapter is therefore formulated for nets, and no
sequential monotone convergence theorem is used.

The strategy that resolves all three at once has already been announced in
\Cref{sec:ch24-why}.  Instead of extending the measure directly, we extend
the \emph{number} \(\mathfrak E\) of \Cref{def:ch24-mixed-energy}.  A
decreasing net of real numbers converges, by
\Cref{lem:monotone-net-converges}; independence of the net is a two-line
sandwich argument; symmetry is inherited from
\Cref{prop:ch24-energy-basic}\,\cref{itm:ch24-eb:sym}; and the measure is then
recovered by the Riesz representation theorem.  The whole of the analysis is
concentrated in \Cref{lem:ch24-sandwich}.

\subsection{Two density lemmas}
\label{sec:ch24-density}

The test space for the Riesz theorem is \(C^0(X^{\An})\).  We shall only be
able to compute integrals of \(\theta\)-psh model functions, so we must know
that these span a dense subspace.

\begin{lemma}
\label{lem:ch24-psh-difference}
Assume \eqref{eq:ch24-standing}.  For every model function \(f\) there are
\(t>0\) and \(\theta\)-psh model functions \(f_+,f_-\) with
\[
              f=\tfrac1{2t}\,(f_+-f_-).
\]
More precisely, after choosing a common model for \(f\), one can find a
model function \(\psi_f\), generally depending on that model, such that
\(f_\pm=\psi_f\pm tf\) for all sufficiently small \(t>0\).
\end{lemma}

\begin{proof}
By \Cref{prop:ch20-positive-representative} there are a model
\(\mathcal W\) and a model function \(\psi_0\) such that
\(\theta+\ddc\psi_0\) is determined on \(\mathcal W\) by a relatively
ample class.  Choose a model \(\cX\) dominating \(\mathcal W\) on which
\(f=\varphi_D\) is determined.  The pullback of the chosen ample class can
lose strict positivity on exceptional curves; applying
\Cref{lem:ch20-ample-propagation} supplies an exceptional correction and a
model function \(\psi_f\) for which \(\theta+\ddc\psi_f\) is represented on
\(\cX\) by a relatively ample class \(\mathcal A_f\).  This is why the
auxiliary potential is allowed to depend on \(f\).  By
\Cref{lem:ch20-special-fibre-test}\,\cref{itm:ch20-sft:ample} relative
ampleness of a class in \(N^1(\cX/R)\) is ampleness of its image in
\(N^1(\cX_s)\), and the ample cone of the projective \(k\)-scheme
\(\cX_s\) is open (Kleiman's criterion, \cite[\S1.4]{PAG1}; this is the
import recorded in \Cref{sec:ch20-imported},
item~\cref{itm:ch20-import:kleiman}).  Hence \(\mathcal A_f\pm t[D]\) is
relatively ample for all sufficiently small \(t>0\), so
\(\theta+\ddc(\psi_f\pm tf)\) is semipositive, that is, \(\psi_f\pm tf\) is
\(\theta\)-psh.  Finally
\((\psi_f+tf)-(\psi_f-tf)=2tf\).
\end{proof}

\begin{lemma}
\label{lem:ch24-dense-tests}
Let \(\mathcal C\subseteq C^0(X^{\An})\) be the real linear span of the
\(\theta\)-psh model functions.  Then \(\mathcal C\) is dense in
\(C^0(X^{\An})\) for the uniform norm and contains the constants.
Consequently:
\begin{enumerate}
\item\label{itm:ch24-dt:unique} two positive Radon measures on \(X^{\An}\)
 agreeing on \(\mathcal C\) are equal;
\item\label{itm:ch24-dt:conv} if \((\mu_\alpha)\) is a net of positive Radon
 measures with \(\sup_\alpha\mu_\alpha(X^{\An})=:C<\infty\) and
 \(\lim_\alpha\int h\,\rd\mu_\alpha\) exists for every
 \(h\in\mathcal C\), then \((\mu_\alpha)\) converges weakly to a positive
 Radon measure \(\mu\), and \(\mu(X^{\An})=\lim_\alpha\mu_\alpha(X^{\An})\).
\end{enumerate}
\end{lemma}

\begin{proof}
Model functions are dense in \(C^0(X^{\An})\) by
\Cref{prop:metrics-model-function-density}, and each of them lies in
\(\mathcal C\) by \Cref{lem:ch24-psh-difference}; the constants are
\(\theta\)-psh by \Cref{prop:ch21-constants}.  Statement
\cref{itm:ch24-dt:unique} is \Cref{cor:riesz-uniqueness} together with
density.

For \cref{itm:ch24-dt:conv}, let \(g\in C^0(X^{\An})\) and
\(\varepsilon>0\), and choose \(h\in\mathcal C\) with
\(\|g-h\|_\infty\leq\varepsilon\).  For all indices \(\alpha,\beta\),
\[
 \left|\int g\,\rd\mu_\alpha-\int g\,\rd\mu_\beta\right|
 \leq 2C\varepsilon
 +\left|\int h\,\rd\mu_\alpha-\int h\,\rd\mu_\beta\right| .
\]
Since \(\int h\,\rd\mu_\alpha\) converges, the right-hand side is at most
\(2C\varepsilon+\varepsilon\) for \(\alpha,\beta\) large.  Hence
\(\bigl(\int g\,\rd\mu_\alpha\bigr)_\alpha\) is a Cauchy net in \(\mathbf R\)
and converges; write \(\Lambda(g)\) for its limit.  \(\Lambda\) is linear,
and \(\Lambda(g)\geq0\) for \(g\geq0\), so \Cref{thm:riesz-interface} gives a
unique positive Radon measure \(\mu\) with \(\Lambda(g)=\int g\,\rd\mu\).
By construction \(\mu_\alpha\to\mu\) weakly, and taking \(g=1\) gives the
mass statement.
\end{proof}

\begin{remark}[Why not test against all continuous functions directly]
\label{rem:ch24-why-psh-tests}
The reason for insisting on the smaller family \(\mathcal C\) is that the
only formula available for computing \(\int h\,\MA_\theta(\varphi_\bullet)\)
is \eqref{eq:ch24-energy-slot}, and it requires \(h\) to be a legitimate
entry of the pairing \(\mathfrak E\), that is, a \(\theta\)-psh function.
\Cref{lem:ch24-dense-tests} is what converts information about that small
family into weak convergence of measures.  This is also why the constant
\(2V\) and not a norm of \(h\) appears in
\eqref{eq:ch24-model-uniform-estimate}.
\end{remark}

\subsection{The mixed energy pairing on bounded potentials}
\label{sec:ch24-energy-bounded}

The next lemma is the exact point at which the ambiguity of the regularizing
net is removed, and the only place where compactness of \(X^{\An}\) is used
to compare two unrelated nets.  Everything else in this section is either
intersection theory or the Riesz representation theorem.

\begin{lemma}
\label{lem:ch24-sandwich}
Let \(u\) be a bounded \(\theta\)-psh function, let
\((u^\alpha)_{\alpha\in A}\) be a decreasing net of \(\theta\)-psh functions
with \(u^\alpha\searrow u\) pointwise, and let \(\psi\) be a
\emph{continuous} function on \(X^{\An}\) with \(\psi\geq u\).  Then for
every \(\varepsilon>0\) there is \(\alpha_0\in A\) with
\begin{equation}
\label{eq:ch24-sandwich}
       u^{\alpha}\leq\psi+\varepsilon
       \qquad\text{on }X^{\An},\ \text{for all }\alpha\geq\alpha_0 .
\end{equation}
\end{lemma}

\begin{proof}
Each \(u^\alpha\) is upper semicontinuous
(\Cref{def:ch21-bfj-psh}) and \(\psi\) is continuous, so
\(u^\alpha-\psi\) is upper semicontinuous and
\[
 U_\alpha:=\{x\in X^{\An}\ :\ u^\alpha(x)-\psi(x)<\varepsilon\}
\]
is open.  The net \((u^\alpha)\) is decreasing, so \((U_\alpha)\) is
increasing.  Their union is all of \(X^{\An}\): for \(x\in X^{\An}\) we have
\(\inf_\alpha u^\alpha(x)=u(x)\leq\psi(x)<\psi(x)+\varepsilon\), so
\(u^{\alpha}(x)<\psi(x)+\varepsilon\) for some \(\alpha\).  Since
\(X^{\An}\) is compact (\Cref{thm:berk-proper-compact}), finitely many
\(U_{\alpha_1},\ldots,U_{\alpha_m}\) cover it; directedness of \(A\) supplies
\(\alpha_0\geq\alpha_1,\ldots,\alpha_m\), and then
\(U_{\alpha_0}=X^{\An}\), whence \(U_\alpha=X^{\An}\) for
\(\alpha\geq\alpha_0\).
\end{proof}

\begin{remark}[This is Dini's theorem, in the only form that survives]
\label{rem:ch24-dini-form}
If \(u\) were continuous, \Cref{lem:ch24-sandwich} with \(\psi=u\) would say
that \(u^\alpha\to u\) uniformly, which is exactly
\Cref{prop:dini-nets}.  For merely bounded \(u\), uniform convergence is
false in general---see \Cref{ex:ch24-branch-sum}, where the limit is
discontinuous---and \eqref{eq:ch24-sandwich} with a continuous majorant
\(\psi\) of \(u\) is the correct replacement.  The proof is the proof of
Dini's theorem, and it uses only upper semicontinuity of the members of the
net, not continuity.
\end{remark}

\begin{theorem}
\label{thm:ch24-energy-extension}
There is a unique extension of the mixed energy pairing
\eqref{eq:ch24-mixed-energy} to \((n+1)\)-tuples of bounded \(\theta\)-psh
functions such that
\(\mathfrak E(\varphi^\alpha_0,\ldots,\varphi^\alpha_n)\to
 \mathfrak E(\varphi_0,\ldots,\varphi_n)\)
whenever \(\varphi_i^\alpha\searrow\varphi_i\) are decreasing nets of
\(\theta\)-psh model functions.  It is given by
\begin{equation}
\label{eq:ch24-energy-inf}
 \begin{aligned}
 \mathfrak E(\varphi_0,\ldots,\varphi_n)
 =\inf\bigl\{&\mathfrak E(\psi_0,\ldots,\psi_n):\\
 &\psi_i\ \text{is a \(\theta\)-psh model function and }
   \psi_i\geq\varphi_i\ \text{for all }i\bigr\},
 \end{aligned}
\end{equation}
and it has the following properties.
\begin{enumerate}
\item\label{itm:ch24-ee:finite} It is finite: if \(m_i\leq\varphi_i\leq M_i\)
 then
 \(\bigl(\sum_im_i\bigr)V\leq\mathfrak E(\varphi_0,\ldots,\varphi_n)
   \leq\bigl(\sum_iM_i\bigr)V\).
\item\label{itm:ch24-ee:sym} It is symmetric in its \(n+1\) arguments.
\item\label{itm:ch24-ee:mono} It is nondecreasing in each argument, and
 \(\mathfrak E(\varphi_0+c_0,\ldots,\varphi_n+c_n)
   =\mathfrak E(\varphi_0,\ldots,\varphi_n)+(\sum_jc_j)V\).
\item\label{itm:ch24-ee:affine} It is affine in each argument: for bounded
 \(\theta\)-psh \(\varphi_0,\varphi_0'\) and \(t\in[0,1]\),
 \(\mathfrak E(t\varphi_0+(1-t)\varphi_0',\varphi_1,\ldots)
  =t\,\mathfrak E(\varphi_0,\varphi_1,\ldots)
   +(1-t)\,\mathfrak E(\varphi_0',\varphi_1,\ldots)\).
\item\label{itm:ch24-ee:nets} It is continuous along decreasing nets of
 \emph{bounded \(\theta\)-psh} functions: if
 \(\varphi_i^\alpha\searrow\varphi_i\) with all functions bounded
 \(\theta\)-psh and \(\varphi_i\) bounded, then
 \(\mathfrak E(\varphi_\bullet^\alpha)\to\mathfrak E(\varphi_\bullet)\).
\item\label{itm:ch24-ee:lip} It is \(V\)-Lipschitz in each argument for the
 uniform norm:
 \(|\mathfrak E(\varphi_0,\varphi_1,\ldots)
  -\mathfrak E(\varphi_0',\varphi_1,\ldots)|
  \leq V\|\varphi_0-\varphi_0'\|_\infty\).
\end{enumerate}
\end{theorem}

\begin{proof}
\emph{Step 1: the infimum is finite and monotone.}  Write \(\mathfrak E^-\)
for the right-hand side of \eqref{eq:ch24-energy-inf}.  The admissible set is
nonempty: by \Cref{thm:bfj-decreasing-regularization} each \(\varphi_i\)
admits \(\theta\)-psh model majorants.  By
\Cref{prop:ch24-energy-basic}\,\cref{itm:ch24-eb:mono} the pairing is
nondecreasing in each slot, so if \(m_i\leq\varphi_i\leq M_i\), then for any
admissible \(\psi_\bullet\) we may compare with the constant tuples: on the
one hand \(\psi_i\geq\varphi_i\geq m_i\) gives
\(\mathfrak E(\psi_\bullet)\geq\mathfrak E(m_0,\ldots,m_n)
 =(\sum_im_i)V\) by \eqref{eq:ch24-energy-constants}; on the other hand the
constant tuple \((M_0,\ldots,M_n)\) is itself admissible, so
\(\mathfrak E^-\leq(\sum_iM_i)V\).  Hence \(\mathfrak E^-\) is finite and
\cref{itm:ch24-ee:finite} holds.  Monotonicity of \(\mathfrak E^-\) in each
argument is clear, since enlarging \(\varphi_i\) shrinks the admissible set.

\emph{Step 2: \(\mathfrak E^-\) is the limit along any decreasing net of
model majorants.}  Let \(\varphi_i^\alpha\searrow\varphi_i\) be decreasing
nets of \(\theta\)-psh model functions.  We may index them by one directed
set: replace the individual index sets \(A_0,\ldots,A_n\) by their product
\(A:=A_0\times\cdots\times A_n\) with the componentwise order, which is
directed, and set \(\varphi_i^{(\alpha_0,\ldots,\alpha_n)}
:=\varphi_i^{\alpha_i}\).  Each of the resulting nets is decreasing with the
same limit, so no generality is lost.  The
numbers \(\mathfrak E(\varphi_\bullet^\alpha)\) form a decreasing net of
reals, by \cref{itm:ch24-eb:mono} of \Cref{prop:ch24-energy-basic}, so they
converge to their infimum, by \Cref{lem:monotone-net-converges}.  Each
\(\varphi_\bullet^\alpha\) is admissible in \eqref{eq:ch24-energy-inf}, so
\(\lim_\alpha\mathfrak E(\varphi^\alpha_\bullet)\geq\mathfrak E^-\).

Conversely let \(\psi_\bullet\) be admissible and \(\varepsilon>0\).  Each
\(\psi_i\) is a continuous function on \(X^{\An}\) with
\(\psi_i\geq\varphi_i\), so \Cref{lem:ch24-sandwich} supplies
\(\alpha_i\) with \(\varphi_i^\alpha\leq\psi_i+\varepsilon\) for
\(\alpha\geq\alpha_i\).  Choosing \(\alpha_0\) above all \(\alpha_i\) and
using monotonicity and \eqref{eq:ch24-energy-constants},
\[
 \mathfrak E(\varphi^\alpha_0,\ldots,\varphi^\alpha_n)
 \leq\mathfrak E(\psi_0+\varepsilon,\ldots,\psi_n+\varepsilon)
 =\mathfrak E(\psi_0,\ldots,\psi_n)+(n+1)\varepsilon V
\]
for \(\alpha\geq\alpha_0\).  Letting \(\alpha\) increase, then
\(\varepsilon\downarrow0\), then taking the infimum over admissible
\(\psi_\bullet\), gives
\(\lim_\alpha\mathfrak E(\varphi^\alpha_\bullet)\leq\mathfrak E^-\).  The two
inequalities give equality, which proves both the existence of the limit and
its independence of the net, hence also uniqueness of the extension.  We
write \(\mathfrak E\) for \(\mathfrak E^-\) from now on.

\emph{Step 3: symmetry, affineness, constants.}  Each of these holds for
model tuples by \Cref{prop:ch24-energy-basic} and passes to the limit along
decreasing nets by Step~2.  For symmetry and \eqref{eq:ch24-energy-constants}
this is immediate.  For affineness, note that if
\(\varphi_0^\alpha\searrow\varphi_0\) and
\(\varphi_0'^\alpha\searrow\varphi_0'\), then
\(t\varphi_0^\alpha+(1-t)\varphi_0'^\alpha\) is again a decreasing net of
\(\theta\)-psh model functions
(\Cref{thm:ch21-stability}\,\cref{itm:ch21-stab:sum}) decreasing to
\(t\varphi_0+(1-t)\varphi_0'\); apply Step~2 to all three nets.

\emph{Step 4: Lipschitz continuity.}  If
\(\|\varphi_0-\varphi_0'\|_\infty\leq\delta\), then
\(\varphi_0\leq\varphi_0'+\delta\) and \(\varphi_0'\leq\varphi_0+\delta\),
so monotonicity and \eqref{eq:ch24-energy-constants} give
\(|\mathfrak E(\varphi_0,\ldots)-\mathfrak E(\varphi_0',\ldots)|\leq\delta V\).
This is \cref{itm:ch24-ee:lip}.

\emph{Step 5: monotone continuity for bounded nets.}  Let
\(\varphi_i^\alpha\searrow\varphi_i\) be decreasing nets of \emph{bounded
\(\theta\)-psh} functions with bounded limits.  Monotonicity gives
\(\mathfrak E(\varphi^\alpha_\bullet)\geq\mathfrak E(\varphi_\bullet)\) for
every \(\alpha\), so
\(\lim_\alpha\mathfrak E(\varphi^\alpha_\bullet)\geq
 \mathfrak E(\varphi_\bullet)\), the limit existing because the net of
numbers is decreasing.  For the reverse inequality, fix an admissible
model tuple \(\psi_\bullet\geq\varphi_\bullet\) and \(\varepsilon>0\).  Each
\(\psi_i\) is continuous and \(\geq\varphi_i\), so
\Cref{lem:ch24-sandwich}---applied to the decreasing net
\((\varphi_i^\alpha)_\alpha\) of \(\theta\)-psh functions, whose members are
upper semicontinuous---gives \(\alpha_i\) with
\(\varphi_i^\alpha\leq\psi_i+\varepsilon\) for \(\alpha\geq\alpha_i\).
Hence, for \(\alpha\) above all the \(\alpha_i\),
\[
 \mathfrak E(\varphi^\alpha_\bullet)
 \leq\mathfrak E(\psi_0+\varepsilon,\ldots,\psi_n+\varepsilon)
 =\mathfrak E(\psi_\bullet)+(n+1)\varepsilon V .
\]
Taking the limit in \(\alpha\), then \(\varepsilon\downarrow0\), then the
infimum over \(\psi_\bullet\), yields
\(\lim_\alpha\mathfrak E(\varphi^\alpha_\bullet)\leq
 \mathfrak E(\varphi_\bullet)\).  This proves \cref{itm:ch24-ee:nets} and
completes the proof.
\end{proof}

\begin{remark}[What was and was not used]
\label{rem:ch24-energy-extension-audit}
The proof used: monotonicity and the constant-shift rule for model tuples,
which are intersection-theoretic; the existence of \emph{some} decreasing
net of \(\theta\)-psh model majorants, which is
\Cref{thm:bfj-decreasing-regularization} and is the only place where
residue characteristic zero enters; compactness of \(X^{\An}\) and upper
semicontinuity of \(\theta\)-psh functions, through
\Cref{lem:ch24-sandwich}; and convergence of monotone nets of real numbers,
\Cref{lem:monotone-net-converges}.  It did \emph{not} use weak compactness
of measures, any Choquet or countability reduction, or any property of the
measures beyond their masses.  In particular no sequential monotone
convergence theorem was invoked, and no statement was made about nets of
unbounded potentials.
\end{remark}

\subsection{The extension theorem}
\label{sec:ch24-extension-theorem}

\begin{theorem}\label{thm:bounded-ma-extension}
There is a unique symmetric map
\[
 (\varphi_1,\ldots,\varphi_n)\longmapsto
 \MA_\theta(\varphi_1,\ldots,\varphi_n)
\]
from \(n\)-tuples of bounded \(\theta\)-psh functions to positive Radon
measures of mass \(V\) with the following properties.
\begin{enumerate}
\item It agrees with \Cref{def:mixed-model-measure} on model
functions.
\item If \(\varphi_i^\alpha\searrow\varphi_i\) and
\(h^\alpha\searrow h\) are decreasing nets of bounded
\(\theta\)-psh functions, then
\begin{equation}
\label{eq:weighted-monotone-ma}
 \int h^\alpha\,
   \MA_\theta(\varphi_1^\alpha,\ldots,\varphi_n^\alpha)
 \longrightarrow
 \int h\,\MA_\theta(\varphi_1,\ldots,\varphi_n).
\end{equation}
\end{enumerate}
This is \cite[Theorem~3.1]{BFJ15}.
\end{theorem}

\begin{proof}[Construction and independence]
Throughout, \(\mathfrak E\) is the extended mixed energy pairing of
\Cref{thm:ch24-energy-extension}.

\emph{Step 1: the defining functional.}  Let \(\varphi_1,\ldots,\varphi_n\)
be bounded \(\theta\)-psh functions.  For a \(\theta\)-psh model function
\(h\) put
\begin{equation}
\label{eq:ch24-lambda}
 \Lambda(h):=\mathfrak E(h,\varphi_1,\ldots,\varphi_n)
             -\mathfrak E(0,\varphi_1,\ldots,\varphi_n).
\end{equation}
By \eqref{eq:ch24-energy-slot} this is exactly
\(\int h\,\MA_\theta(\varphi_1,\ldots,\varphi_n)\) when all the
\(\varphi_i\) are model functions, so \(\Lambda\) is the only candidate.
Note that \(\Lambda\) is defined without reference to any approximating net.

\emph{Step 2: \(\Lambda\) is the limit of model integrals.}  Choose
decreasing nets \(\varphi_i^\alpha\searrow\varphi_i\) of \(\theta\)-psh
model functions, indexed by a common directed set, and put
\(\mu_\alpha:=\MA_\theta(\varphi_1^\alpha,\ldots,\varphi_n^\alpha)\), a
positive measure of mass \(V\) by \Cref{prop:model-ma-basic}.  For every
\(\theta\)-psh model function \(h\), \eqref{eq:ch24-energy-slot} gives
\[
 \int h\,\rd\mu_\alpha
 =\mathfrak E(h,\varphi^\alpha_1,\ldots,\varphi^\alpha_n)
  -\mathfrak E(0,\varphi^\alpha_1,\ldots,\varphi^\alpha_n)
 \xrightarrow[\ \alpha\ ]{}\ \Lambda(h),
\]
by \Cref{thm:ch24-energy-extension} applied to the two constant-in-one-slot
nets \((h,\varphi^\alpha_\bullet)\) and \((0,\varphi^\alpha_\bullet)\); a
constant net is decreasing.

\emph{Step 3: existence of the measure.}  The masses \(\mu_\alpha(X^{\An})\)
are all equal to \(V\).  By Step~2 the limit
\(\lim_\alpha\int h\,\rd\mu_\alpha\) exists for every \(\theta\)-psh model
function \(h\); since \(h\mapsto\int h\,\rd\mu_\alpha\) is linear, the limit
then exists for every \(h\) in the linear span \(\mathcal C\) of the
\(\theta\)-psh model functions.  (No linearity of \(\Lambda\) is asserted or
needed: it is the integrals, not the defining formula, that are linear in
\(h\).)  By
\Cref{lem:ch24-dense-tests}\,\cref{itm:ch24-dt:conv} the net
\((\mu_\alpha)\) converges weakly to a positive Radon measure of mass
\(V\).  Define
\[
 \MA_\theta(\varphi_1,\ldots,\varphi_n):=\lim_\alpha\mu_\alpha .
\]
Because \(\Lambda\) does not involve the net, the limit is the same for
every choice of regularizing nets, which is the required independence.

\emph{Step 4: symmetry and compatibility.}  Compatibility with
\Cref{def:mixed-model-measure} is Step~1 read for model \(\varphi_i\).
Symmetry follows from
\Cref{thm:ch24-energy-extension}\,\cref{itm:ch24-ee:sym}: the functional
\eqref{eq:ch24-lambda} is symmetric in \(\varphi_1,\ldots,\varphi_n\), hence
so is the measure it determines, by
\Cref{lem:ch24-dense-tests}\,\cref{itm:ch24-dt:unique}.

\emph{Step 5: the identity \eqref{eq:ch24-energy-slot} for bounded
entries.}  We claim that for bounded \(\theta\)-psh
\(\varphi_0,\ldots,\varphi_n\),
\begin{equation}
\label{eq:ch24-energy-slot-bounded}
 \mathfrak E(\varphi_0,\varphi_1,\ldots,\varphi_n)
 -\mathfrak E(0,\varphi_1,\ldots,\varphi_n)
 =\int_{X^{\An}}\varphi_0\,\MA_\theta(\varphi_1,\ldots,\varphi_n).
\end{equation}
Indeed, choose \(\theta\)-psh model functions
\(\varphi_0^\beta\searrow\varphi_0\).  The left-hand side is the limit of the
same expression with \(\varphi_0^\beta\) in place of \(\varphi_0\), by
\Cref{thm:ch24-energy-extension}\,\cref{itm:ch24-ee:nets}, and that
expression equals \(\int\varphi_0^\beta\,
\MA_\theta(\varphi_1,\ldots,\varphi_n)\) by Step~1 and Step~3.  The
right-hand side is the limit of these integrals by monotone convergence
along a decreasing net of continuous functions against the \emph{fixed}
Radon measure \(\MA_\theta(\varphi_1,\ldots,\varphi_n)\), which is
\Cref{cor:usc-net-monotone-integral}.  This proves
\eqref{eq:ch24-energy-slot-bounded}.

\emph{Step 6: weighted monotone continuity.}  Let
\(\varphi_i^\alpha\searrow\varphi_i\) and \(h^\alpha\searrow h\) be
decreasing nets of bounded \(\theta\)-psh functions, indexed by a common
directed set.  By Step~5,
\[
 \int h^\alpha\,\MA_\theta(\varphi^\alpha_1,\ldots,\varphi^\alpha_n)
 =\mathfrak E(h^\alpha,\varphi^\alpha_1,\ldots,\varphi^\alpha_n)
  -\mathfrak E(0,\varphi^\alpha_1,\ldots,\varphi^\alpha_n),
\]
and both terms converge, to
\(\mathfrak E(h,\varphi_1,\ldots,\varphi_n)\) and
\(\mathfrak E(0,\varphi_1,\ldots,\varphi_n)\) respectively, by
\Cref{thm:ch24-energy-extension}\,\cref{itm:ch24-ee:nets} applied to the
\((n+1)\)-tuples \((h^\alpha,\varphi^\alpha_\bullet)\) and
\((0,\varphi^\alpha_\bullet)\).  By Step~5 again the limit is
\(\int h\,\MA_\theta(\varphi_1,\ldots,\varphi_n)\).  This is
\eqref{eq:weighted-monotone-ma}.  Taking \(h^\alpha=h\) constant in
\(\alpha\), with \(h\) a \(\theta\)-psh model function, and then extending to
\(\mathcal C\) by linearity of the integral in \(h\), one obtains from
\Cref{lem:ch24-dense-tests}\,\cref{itm:ch24-dt:conv} the weak convergence
\(\MA_\theta(\varphi^\alpha_\bullet)\to\MA_\theta(\varphi_\bullet)\).

\emph{Step 7: uniqueness.}  Suppose \(N\) is another map with properties (1)
and (2).  Given bounded \(\theta\)-psh \(\varphi_1,\ldots,\varphi_n\),
choose decreasing nets of \(\theta\)-psh model functions
\(\varphi_i^\alpha\searrow\varphi_i\).  Property (1) gives
\(N(\varphi^\alpha_\bullet)=\MA_\theta(\varphi^\alpha_\bullet)\) for every
\(\alpha\), and property (2), applied with the constant test net
\(h^\alpha=h\) for a \(\theta\)-psh model function \(h\), forces
\(\int h\,\rd N(\varphi_\bullet)=\lim_\alpha\int h\,
\rd\MA_\theta(\varphi^\alpha_\bullet)
=\int h\,\rd\MA_\theta(\varphi_\bullet)\), the last equality by Step~6.
Both sides being linear in \(h\), the identity extends to \(\mathcal C\).  By
\Cref{lem:ch24-dense-tests}\,\cref{itm:ch24-dt:unique},
\(N(\varphi_\bullet)=\MA_\theta(\varphi_\bullet)\).
\end{proof}

\begin{remark}[Where each hypothesis is used]
\label{rem:ch24-extension-audit}
Boundedness of the potentials is used twice and is not removable by any
formal device.  It is used in \Cref{thm:ch24-energy-extension} to keep
\(\mathfrak E\) finite, and it is used in Step~5 through
\Cref{cor:usc-net-monotone-integral}, whose conclusion may be \(-\infty\) for
unbounded integrands.  Ampleness of \(\{\theta\}\) is used through
\Cref{lem:ch24-psh-difference}, hence through the density
\Cref{lem:ch24-dense-tests}; it also guarantees \(V>0\).  Residue
characteristic zero is used only through
\Cref{thm:bfj-decreasing-regularization}, that is, only to know that
decreasing model regularizations exist at all; see \Cref{sec:ch24-scope}.
Nothing in the proof requires \(X\) to be a curve, or the measures to be
atomic, or the index sets to be countable.
\end{remark}

\subsection{Two formulae from the construction}
\label{sec:ch24-two-formulae}

The construction just given routes everything through the scalar
\(\mathfrak E\).  The original argument of \cite[\S3]{BFJ15} instead
constructs the measure one slot at a time, by \emph{forcing} the model
integration-by-parts formula to hold.  Both formulae used there are now
theorems, and we record them, both because they are the shape in which the
result is usually quoted and because the second is the quantitative
statement behind the qualitative monotone continuity.

\begin{proposition}
\label{prop:ch24-forced-ibp}
Let \(1\leq p\leq n\), let \(\varphi_1,\ldots,\varphi_p\) be bounded
\(\theta\)-psh functions, let \(\varphi_{p+1},\ldots,\varphi_n\) be
\(\theta\)-psh model functions, and write
\(M_p(\varphi_1,\ldots,\varphi_p):=
\MA_\theta(\varphi_1,\ldots,\varphi_n)\) and
\(T:=\MA_\theta(\varphi_1,\ldots,\varphi_{p-1},
\varphi_{p+1},\ldots,\varphi_n)\) for the corresponding
\((n-1)\)-fold product.  Then, for every \(\theta\)-psh model function
\(h\),
\begin{equation}
\label{eq:inductive-ma-definition}
\begin{split}
 \int h\,M_p(\varphi_1,\ldots,\varphi_p)
 ={}&\int\varphi_p(\theta+\ddc h)\wedge T\\
     &+\int(h-\varphi_p)\theta\wedge T,
\end{split}
\end{equation}
where \((\theta+\ddc h)\wedge T\) and \(\theta\wedge T\) denote the mixed
products in which the \(p\)-th slot carries the model potential \(h\),
respectively \(0\).  Equivalently,
\[
 \int h\,M_p(\varphi_1,\ldots,\varphi_p)
 =\int h\,\theta\wedge T+\int\varphi_p\,\ddc h\wedge T .
\]
\end{proposition}

\begin{proof}
The two displayed forms are equivalent: subtract and add
\(\int\varphi_p\,\theta\wedge T\).  For the second form, apply
\eqref{eq:ch24-energy-slot-bounded} twice.  With
\(\Phi:=(\varphi_1,\ldots,\varphi_{p-1},\varphi_{p+1},\ldots,\varphi_n)\)
denoting the remaining \(n-1\) entries,
\[
 \begin{aligned}
 \int h\,M_p(\varphi_1,\ldots,\varphi_p)
   &=\mathfrak E(h,\varphi_p,\Phi)-\mathfrak E(0,\varphi_p,\Phi),\\
 \int\varphi_p\,\ddc h\wedge T
   &=\mathfrak E(\varphi_p,h,\Phi)-\mathfrak E(\varphi_p,0,\Phi),
 \end{aligned}
\]
the second because \(\ddc h\wedge T\) is by definition the difference of the
mixed measures with \(h\) and with \(0\) in the \(p\)-th slot, and because
\eqref{eq:ch24-energy-slot-bounded} may be applied to each.  Also
\(\int h\,\theta\wedge T=\mathfrak E(h,0,\Phi)-\mathfrak E(0,0,\Phi)\).
Symmetry of \(\mathfrak E\)
(\Cref{thm:ch24-energy-extension}\,\cref{itm:ch24-ee:sym}) gives
\(\mathfrak E(\varphi_p,h,\Phi)=\mathfrak E(h,\varphi_p,\Phi)\) and
\(\mathfrak E(\varphi_p,0,\Phi)=\mathfrak E(0,\varphi_p,\Phi)\), so
  \begin{align*}
 \int h\,\theta\wedge T+\int\varphi_p\,\ddc h\wedge T
 ={}&\bigl[\mathfrak E(h,0,\Phi)-\mathfrak E(0,0,\Phi)\bigr]\\
 &+\bigl[\mathfrak E(h,\varphi_p,\Phi)-\mathfrak E(h,0,\Phi)\bigr]\\
 &-\bigl[\mathfrak E(0,\varphi_p,\Phi)-\mathfrak E(0,0,\Phi)\bigr],
  \end{align*}
which telescopes to
\(\mathfrak E(h,\varphi_p,\Phi)-\mathfrak E(0,\varphi_p,\Phi)\), as
required.
\end{proof}

\begin{remark}[Reading \eqref{eq:inductive-ma-definition} as a definition]
\label{rem:ch24-forced-ibp-as-definition}
The right-hand side of \eqref{eq:inductive-ma-definition} involves
\(\varphi_p\) only through integration against the two \emph{fixed} positive
measures \((\theta+\ddc h)\wedge T\) and \(\theta\wedge T\), which are built
from the other entries.  This is why it can be used as a definition: if
\(\varphi_p^\alpha\searrow\varphi_p\) is a decreasing net of bounded
\(\theta\)-psh functions, the right-hand side converges by
\Cref{cor:usc-net-monotone-integral}, applied separately to the two positive
measures, and the limit is manifestly independent of the net.  What such an
argument does not give directly is symmetry in the \(n\) slots, since it
singles out the slot \(p\); this is the reason we have routed the
construction through \(\mathfrak E\), where symmetry is built in.
\end{remark}

\begin{proposition}
\label{prop:ch24-approximate-monotonicity}
Let \(1\leq p\leq n\), let \(\varphi_i\leq\chi_i\) be \(\theta\)-psh model
functions for \(1\leq i\leq p\), let the remaining entries
\(\varphi_{p+1},\ldots,\varphi_n\) be \(\theta\)-psh model functions, and let
\(h\) be a \(\theta\)-psh model function.  Write \(M_p\) as in
\Cref{prop:ch24-forced-ibp}, the entry \(0\) meaning that the corresponding
slot carries the form \(\theta\).  Then
\begin{equation}
\label{eq:approximate-monotonicity}
\begin{split}
 \int h\,M_p(\chi_1,\ldots,\chi_p)
 \geq{}&\int h\,M_p(\varphi_1,\ldots,\varphi_p)\\
 &+\sum_{i=1}^p\int(\varphi_i-\chi_i)
       M_p(\varphi_1,\ldots,\varphi_{i-1},0,
           \chi_{i+1},\ldots,\chi_p).
\end{split}
\end{equation}
In particular
\(\int h\,M_p(\chi_\bullet)\geq\int h\,M_p(\varphi_\bullet)
 -V\sum_{i=1}^p\|\chi_i-\varphi_i\|_\infty\).
\end{proposition}

\begin{proof}
Telescope.  For \(0\leq i\leq p\) put
\(A_i:=M_p(\varphi_1,\ldots,\varphi_i,\chi_{i+1},\ldots,\chi_p)\), so that
\(A_0=M_p(\chi_\bullet)\) and \(A_p=M_p(\varphi_\bullet)\).  Fix \(i\) and
let \(S_i\) denote the \((n-1)\)-fold mixed product obtained by deleting the
\(i\)-th slot from \(A_{i-1}\), that is, the product with entries
\(\varphi_1,\ldots,\varphi_{i-1},\chi_{i+1},\ldots,\chi_p\) and
\(\varphi_{p+1},\ldots,\varphi_n\).  Then
\[
 A_{i-1}-A_i=\ddc(\chi_i-\varphi_i)\wedge S_i ,
\]
by multilinearity in the \(i\)-th slot, and
\eqref{eq:model-ibp} gives
\[
 \begin{aligned}
 \int h\,\rd(A_{i-1}-A_i)
 &=\int(\chi_i-\varphi_i)\,\ddc h\wedge S_i\\
 &=\int(\chi_i-\varphi_i)(\theta+\ddc h)\wedge S_i\\
 &\quad-\int(\chi_i-\varphi_i)\,\theta\wedge S_i .
 \end{aligned}
\]
The first term on the right is nonnegative, since \(\chi_i-\varphi_i\geq0\)
and \((\theta+\ddc h)\wedge S_i\) is a positive measure by
\Cref{prop:model-ma-basic}.  Hence
\[
 \begin{aligned}
 \int h\,\rd(A_{i-1}-A_i)
 &\geq-\int(\chi_i-\varphi_i)\,\theta\wedge S_i\\
 &=\int(\varphi_i-\chi_i)\,
   M_p(\varphi_1,\ldots,\varphi_{i-1},0,
       \chi_{i+1},\ldots,\chi_p),
 \end{aligned}
\]
since \(\theta\wedge S_i\) is precisely the mixed product with \(0\) in the
\(i\)-th slot.  Summing over \(i=1,\ldots,p\) gives
\eqref{eq:approximate-monotonicity}.  The final inequality follows because
each measure occurring there is positive of mass \(V\) by
\Cref{thm:ch24-mass}.
\end{proof}

\begin{remark}[Two monotonicities, not one]
\label{rem:ch24-two-monotonicities}
\Cref{prop:ch24-approximate-monotonicity} is an inequality between integrals
of a \emph{fixed} test function against measures attached to different
potentials; the error term is controlled by the sup-norm distance of the
potentials.  It is therefore adapted to \emph{uniform} approximation, and is
what one uses to treat continuous potentials in
\Cref{sec:ch24-continuous}.  It is \emph{not} what proves monotone
continuity along a general decreasing net, where the sup-norm distance need
not tend to zero: there the monotone quantity is
\(\mathfrak E\), a number, and the mechanism is
\Cref{thm:ch24-energy-extension}.  Conflating the two is the commonest
error in this area, and the reason the statement
\eqref{eq:weighted-monotone-ma} is weighted.
\end{remark}

\subsection{Multilinearity, convexity, and the Dirichlet form}
\label{sec:ch24-bounded-calculus}

\begin{corollary}\label{cor:bounded-ma-multilinearity}
Mixed products are symmetric and affine in each psh entry. In particular,
for \(0\leq t\leq1\),
\begin{equation}
\label{eq:ma-convex-lower-bound}
 \MA_\theta(t\varphi+(1-t)\psi)
 \geq t^n\MA_\theta(\varphi)+(1-t)^n\MA_\theta(\psi).
\end{equation}
Model integration by parts and its Cauchy--Schwarz inequality extend to
bounded psh differences.
\end{corollary}

\begin{proof}
\emph{Symmetry} is
\Cref{thm:bounded-ma-extension}, Step~4.

\emph{Affineness.}  Let \(\varphi,\varphi'\) be bounded \(\theta\)-psh and
\(t\in[0,1]\); we must show
\[
 \begin{aligned}
 &\MA_\theta\bigl(t\varphi+(1-t)\varphi',
     \varphi_2,\ldots,\varphi_n\bigr)\\
 &\qquad=t\,\MA_\theta(\varphi,\varphi_2,\ldots,\varphi_n)\\
 &\qquad\quad+(1-t)\,\MA_\theta(\varphi',\varphi_2,\ldots,\varphi_n).
 \end{aligned}
\]
By \Cref{lem:ch24-dense-tests}\,\cref{itm:ch24-dt:unique} it suffices to test
against \(\theta\)-psh model functions \(h\), and by
\eqref{eq:ch24-energy-slot-bounded} the required identity is
\[
 \mathfrak E\bigl(h,t\varphi+(1-t)\varphi',\varphi_2,\ldots\bigr)
 -\mathfrak E\bigl(0,t\varphi+(1-t)\varphi',\varphi_2,\ldots\bigr)
 =t\bigl[\cdots\bigr]+(1-t)\bigl[\cdots\bigr],
\]
which is
\Cref{thm:ch24-energy-extension}\,\cref{itm:ch24-ee:affine} applied in the
second slot.

\emph{Inequality \eqref{eq:ma-convex-lower-bound}.}  Put
\(u=t\varphi+(1-t)\psi\), a bounded \(\theta\)-psh function by
\Cref{thm:ch21-stability}\,\cref{itm:ch21-stab:sum}.  Expanding the diagonal
product by affineness in each of the \(n\) slots,
\[
 \MA_\theta(u)=\sum_{k=0}^{n}\binom nk t^k(1-t)^{n-k}\,
   \MA_\theta\bigl(\varphi^{(k)},\psi^{(n-k)}\bigr),
\]
where \(\varphi^{(k)}\) means \(k\) entries equal to \(\varphi\).  All terms
are positive measures, and the terms \(k=n\) and \(k=0\) are
\(t^n\MA_\theta(\varphi)\) and \((1-t)^n\MA_\theta(\psi)\).  Discarding the
remaining \(n-1\) nonnegative terms gives
\eqref{eq:ma-convex-lower-bound}.

\emph{Integration by parts and Cauchy--Schwarz} are
\Cref{prop:ch24-bounded-dirichlet} and
\Cref{cor:ch24-bounded-cauchy-schwarz} below, which are stated separately
because the correct formulation for bounded potentials is in terms of
differences and not of individual entries.
\end{proof}

The naive extension of \eqref{eq:model-ibp} to bounded potentials is
problematic for a reason worth naming.  A difference \(u_1-u_2\) of bounded
\(\theta\)-psh functions is neither upper nor lower semicontinuous, so an
expression such as \(\int(u_1-u_2)\,\rd\nu_\alpha\) does not behave well
along a net \(\nu_\alpha\to\nu\) of measures.  The remedy is to write the
\emph{whole} quadratic expression in terms of \(\mathfrak E\), where every
entry is a genuine \(\theta\)-psh function and the monotone continuity
theorem applies.

\begin{proposition}
\label{prop:ch24-bounded-dirichlet}
Let \(u_1,u_2\) and \(w_1,\ldots,w_{n-1}\) be bounded \(\theta\)-psh
functions.  Then
\begin{equation}
\label{eq:ch24-bounded-dirichlet}
\begin{split}
 Q_w(u_1,u_2)
 :={}&\int(u_1-u_2)\,
   \bigl[\MA_\theta(u_1,w_1,\ldots,w_{n-1})
        -\MA_\theta(u_2,w_1,\ldots,w_{n-1})\bigr]\\
 ={}&\mathfrak E(u_1,u_1,w_\bullet)-2\mathfrak E(u_1,u_2,w_\bullet)
     +\mathfrak E(u_2,u_2,w_\bullet)\ \leq\ 0 .
\end{split}
\end{equation}
Equality holds when \(u_1-u_2\) is constant.
\end{proposition}

\begin{proof}
\emph{Step 1: the two expressions agree.}  Apply
\eqref{eq:ch24-energy-slot-bounded} four times, with the two possible
entries in the zeroth and first slots:
\[
 \int u_i\,\MA_\theta(u_j,w_\bullet)
 =\mathfrak E(u_i,u_j,w_\bullet)-\mathfrak E(0,u_j,w_\bullet)
 \qquad (i,j\in\{1,2\}).
\]
Expanding the first line of \eqref{eq:ch24-bounded-dirichlet} into four such
integrals, the four terms \(\mathfrak E(0,u_j,w_\bullet)\) cancel in pairs,
and what remains is
\(\mathfrak E(u_1,u_1,w_\bullet)-\mathfrak E(u_1,u_2,w_\bullet)
 -\mathfrak E(u_2,u_1,w_\bullet)+\mathfrak E(u_2,u_2,w_\bullet)\), which by
symmetry (\Cref{thm:ch24-energy-extension}\,\cref{itm:ch24-ee:sym}) is the
second line.  In particular the left-hand side is well defined, even though
\(u_1-u_2\) is only a bounded Borel function: it is an integral of a bounded
Borel function against a finite signed measure.

\emph{Step 2: the model case.}  If all the entries are \(\theta\)-psh model
functions, then by \eqref{eq:ch24-energy-slot} and
\eqref{eq:model-test-intersection} the first line of
\eqref{eq:ch24-bounded-dirichlet} equals
\(\int(u_1-u_2)\,\ddc(u_1-u_2)\wedge\Theta_w\), where
\(\Theta_w=(\theta+\ddc w_1)\wedge\cdots\wedge(\theta+\ddc w_{n-1})\), and
this is \(\leq0\) by \eqref{eq:model-dirichlet-negativity}.

\emph{Step 3: passage to bounded entries.}  Choose decreasing nets of
\(\theta\)-psh model functions \(u_i^\alpha\searrow u_i\) and
\(w_j^\alpha\searrow w_j\) on a common directed set.  Every one of the four
terms in the second line of \eqref{eq:ch24-bounded-dirichlet} converges, by
\Cref{thm:ch24-energy-extension}\,\cref{itm:ch24-ee:nets}, and by Step~2 the
combination is \(\leq0\) at each index.  Hence the limit is \(\leq0\).  This
is the point of writing \(Q_w\) through \(\mathfrak E\): no integral of a
difference against a moving measure ever has to be passed to the limit.

\emph{Step 4: the equality case.}  If \(u_1=u_2+c\), then by affineness and
\eqref{eq:ch24-energy-constants} all four terms differ from
\(\mathfrak E(u_2,u_2,w_\bullet)\) by multiples of \(cV\) which cancel; the
same is seen from the first line, since \(\MA_\theta(u_2+c,w_\bullet)
=\MA_\theta(u_2,w_\bullet)\).
\end{proof}

\begin{corollary}
\label{cor:ch24-bounded-cauchy-schwarz}
Let \(u_1,u_2,v_1,v_2,w_1,\ldots,w_{n-1}\) be bounded \(\theta\)-psh
functions and set
\[
 \begin{aligned}
 B_w\bigl((u_1,u_2),(v_1,v_2)\bigr)
 :={}&\mathfrak E(u_1,v_1,w_\bullet)-\mathfrak E(u_1,v_2,w_\bullet)\\
 &-\mathfrak E(u_2,v_1,w_\bullet)+\mathfrak E(u_2,v_2,w_\bullet).
 \end{aligned}
\]
Then \(B_w\) is symmetric and bilinear in the two differences, and
\begin{equation}
\label{eq:ch24-bounded-cauchy-schwarz}
 \bigl|B_w\bigl((u_1,u_2),(v_1,v_2)\bigr)\bigr|^2
 \leq Q_w(u_1,u_2)\;Q_w(v_1,v_2),
\end{equation}
both factors on the right being nonpositive.
\end{corollary}

\begin{proof}
For model entries, Step~2 of the previous proof identifies
\(B_w\) with \(\bigl(D\cdot D'\cdot\Theta_w\bigr)\), where
\(\ddc(u_1-u_2)=[D]\) and \(\ddc(v_1-v_2)=[D']\), and
\eqref{eq:ch24-bounded-cauchy-schwarz} is then
\eqref{eq:ch20-mixed-cauchy-schwarz}, that is,
\Cref{prop:ch24-model-cauchy-schwarz}.  Both sides converge along decreasing
nets of model approximants by
\Cref{thm:ch24-energy-extension}\,\cref{itm:ch24-ee:nets}, so the inequality
persists.  Symmetry and bilinearity follow from
\Cref{thm:ch24-energy-extension}\,\cref{itm:ch24-ee:sym} and
\cref{itm:ch24-ee:affine}.
\end{proof}

\subsection{Continuous potentials}
\label{sec:ch24-continuous}

Between the model and the bounded theory sits the continuous one, where
approximation is uniform and no net is needed.  It is worth isolating,
because it is the level at which the comparison with Chambert--Loir's
original construction and with the CLD calculus takes place
(\Cref{sec:ch24-comparison}), and because the mechanism is completely
different from that of \Cref{thm:bounded-ma-extension}.

\begin{proposition}
\label{prop:ch24-continuous-uniform}
Let \(\varphi_1,\ldots,\varphi_n\) be continuous \(\theta\)-psh functions.
\begin{enumerate}
\item\label{itm:ch24-cu:seq} There are \(\theta\)-psh model functions
 \(\varphi_i^m\) with \(\varphi_i^m\searrow\varphi_i\) \emph{uniformly}, and
 for any such choice
 \(\MA_\theta(\varphi_1^m,\ldots,\varphi_n^m)\to
   \MA_\theta(\varphi_1,\ldots,\varphi_n)\) weakly.
\item\label{itm:ch24-cu:unif} More generally, if \(\varphi_i^m\) are
 \(\theta\)-psh model functions with
 \(\|\varphi_i^m-\varphi_i\|_\infty\to0\), not necessarily monotonically,
 then \(\MA_\theta(\varphi_1^m,\ldots,\varphi_n^m)\to
 \MA_\theta(\varphi_1,\ldots,\varphi_n)\) weakly.  In particular the
 bounded product restricted to continuous potentials is the unique
 extension that is continuous for uniform convergence.
\item\label{itm:ch24-cu:cont} The operator is continuous on continuous
 potentials: if \(\|\varphi_i^{(k)}-\varphi_i\|_\infty\to0\) with all
 functions continuous \(\theta\)-psh, then the measures converge weakly.
\end{enumerate}
\end{proposition}

\begin{proof}
\cref{itm:ch24-cu:seq}.  Existence of a decreasing sequence of \(\theta\)-psh
model functions converging uniformly to a continuous \(\theta\)-psh function
is \Cref{cor:ch22-continuous-decreasing-sequence}.  A decreasing sequence is
a decreasing net, so weak convergence is
\Cref{thm:bounded-ma-extension}, Step~6.

\cref{itm:ch24-cu:unif}.  Set \(\varepsilon_m:=\max_i\|\varphi_i^m
-\varphi_i\|_\infty\).  By \Cref{prop:ch24-model-uniform-estimate}, for any
\(\theta\)-psh model function \(h\) and any \(m,m'\),
\[
 \left|\int h\,\MA_\theta(\varphi^m_\bullet)
      -\int h\,\MA_\theta(\varphi^{m'}_\bullet)\right|
 \leq 2Vn\,(\varepsilon_m+\varepsilon_{m'}),
\]
so \(\bigl(\int h\,\MA_\theta(\varphi^m_\bullet)\bigr)_m\) is Cauchy.  Since
all the measures have mass \(V\),
\Cref{lem:ch24-dense-tests}\,\cref{itm:ch24-dt:conv} gives a weak limit
\(\mu\), and the same estimate shows that \(\mu\) does not depend on the
approximating sequence.  Comparing with a decreasing sequence as in
\cref{itm:ch24-cu:seq} identifies \(\mu\) with
\(\MA_\theta(\varphi_1,\ldots,\varphi_n)\).

\cref{itm:ch24-cu:cont}.  First note that
\eqref{eq:ch24-model-uniform-estimate} persists for \emph{continuous}
\(\theta\)-psh entries: approximating each entry uniformly by model
functions and using \cref{itm:ch24-cu:unif}, both sides pass to the limit, so
that for continuous \(\theta\)-psh \(\varphi_\bullet,\psi_\bullet\) and
\(\theta\)-psh model \(h\),
\begin{equation}
\label{eq:ch24-continuous-uniform-estimate}
 \left|\int h\,\MA_\theta(\varphi_\bullet)
      -\int h\,\MA_\theta(\psi_\bullet)\right|
 \leq 2V\sum_{i=1}^{n}\|\varphi_i-\psi_i\|_\infty .
\end{equation}
Now let \(\varphi_i^{(k)}\to\varphi_i\) uniformly with all functions
continuous \(\theta\)-psh.  By \eqref{eq:ch24-continuous-uniform-estimate},
\(\int h\,\MA_\theta(\varphi^{(k)}_\bullet)\to
\int h\,\MA_\theta(\varphi_\bullet)\) for every \(\theta\)-psh model
function \(h\), and all the measures have mass \(V\); so
\Cref{lem:ch24-dense-tests}\,\cref{itm:ch24-dt:conv} gives weak convergence.
\end{proof}

\begin{corollary}
\label{cor:ch24-continuous-cl}
Let \(L\) be an ample line bundle with a model metric \(\|\cdot\|_0\) of
curvature \(\theta\), and let \(\varphi\) be a continuous \(\theta\)-psh
function, so that \(\|\cdot\|_\varphi=\rme^{-\varphi}\|\cdot\|_0\) is a
continuous semipositive metric in the sense of Zhang
(\Cref{thm:ch21-continuous-zhang}).  Then
\[
 \MA_\theta(\varphi)=c_1\bigl(L,\|\cdot\|_\varphi\bigr)^{\wedge n},
\]
the Chambert--Loir measure of \Cref{thm:cl-continuous-extension}.  The same
holds for mixed products of \(n\) continuous semipositive metrics.
\end{corollary}

\begin{proof}
Both sides are the weak limit of the same model measures.  On the left this
is \Cref{prop:ch24-continuous-uniform}\,\cref{itm:ch24-cu:unif} together with
\Cref{cor:ch24-cl-identification}; on the right it is the definition of the
continuous Chambert--Loir measure,
\Cref{thm:cl-continuous-extension}.
\end{proof}

\subsection{Integrability and pluripolar sets}
\label{sec:ch24-integrability}

The last property of the bounded product needed downstream is that it is
insensitive to the polar sets of psh functions.  This is what makes capacity
work in \Cref{chap:bfj-capacity}, and it is also the precise place at which
the bounded theory differs from the finite-energy theory of
\Cref{chap:bfj-energy}: there the total mass may drop, and the drop is
exactly the mass that a naive limit would place on a polar set.  Note that
the unbounded function below is integrated \emph{against} the product; it is
never an entry of it.

\begin{proposition}\label{prop:bounded-ma-no-polar-mass}
If \(\varphi_1,\ldots,\varphi_n\) are bounded \(\theta\)-psh and
\(u\in\PSH(X,\theta)\), then \(u\) is integrable with respect to
\(\MA_\theta(\varphi_1,\ldots,\varphi_n)\). Consequently the mixed measure
gives mass zero to every pluripolar set.
\end{proposition}

\begin{proof}
\emph{Step 0: normalizations.}  Replacing \(\theta\) by
\(\theta+\ddc\psi\) and every potential by its difference with \(\psi\), we
may assume by \Cref{prop:ch20-positive-representative} that \(\theta\) is
semipositive, so that \(0\in\PSH(X,\theta)\) and the entry \(0\) is
admissible in \(\mathfrak E\); by
\Cref{cor:ch21-reference-independence} this changes neither the hypothesis
nor the conclusion.  Adding constants to the \(\varphi_i\) changes neither
\(\MA_\theta(\varphi_1,\ldots,\varphi_n)\) nor the class of \(\theta\)-psh
functions, and adding a constant to \(u\) does not affect integrability; so
we may assume
\[
 -M\leq\varphi_i\leq0\ \text{ for all }i,
 \qquad \sup_{X^{\An}}u=0 ,
\]
the supremum being finite by
\Cref{prop:ch21-first-consequences}\,\cref{itm:ch21-fc:bounded}.  Since
\(u\leq0\), integrability of \(u\) means
\(\int(-u)\,\MA_\theta(\varphi_\bullet)<\infty\).

\emph{Step 1: truncation.}  For \(j\in\mathbf N\) put
\(u_j:=\max\{u,-j\}\).  By
\Cref{thm:ch21-stability}\,\cref{itm:ch21-stab:max} and
\cref{itm:ch21-stab:const} each \(u_j\) is \(\theta\)-psh, and it is bounded,
with \(-j\leq u_j\leq0\).  Also \(u_j\searrow u\) pointwise, and
\(-u_j=\min\{-u,j\}\) is a nonnegative, lower semicontinuous function
increasing to \(-u\).

\emph{Step 2: the basic estimate.}  Write
\(\theta^{\wedge n}:=\MA_\theta(0,\ldots,0)\), the model measure of
\(\theta\) itself, a positive atomic measure of mass \(V\).  By
\eqref{eq:ch24-energy-slot-bounded},
\[
 \int(-u_j)\,\MA_\theta(\varphi_1,\ldots,\varphi_n)
 =\mathfrak E(0,\varphi_1,\ldots,\varphi_n)
  -\mathfrak E(u_j,\varphi_1,\ldots,\varphi_n).
\]
We bound the two terms separately, using monotonicity of \(\mathfrak E\)
(\Cref{thm:ch24-energy-extension}\,\cref{itm:ch24-ee:mono}) and the
constant-shift rule.  Since \(\varphi_i\leq0\),
\[
 \mathfrak E(0,\varphi_1,\ldots,\varphi_n)\leq\mathfrak E(0,0,\ldots,0)=0 ;
\]
since \(\varphi_i\geq-M\),
\[
 \mathfrak E(u_j,\varphi_1,\ldots,\varphi_n)
 \geq\mathfrak E(u_j,-M,\ldots,-M)
 =\mathfrak E(u_j,0,\ldots,0)-nMV .
\]
Combining, and using \eqref{eq:ch24-energy-slot-bounded} once more in the
form \(\mathfrak E(u_j,0,\ldots,0)=\int u_j\,\theta^{\wedge n}\),
\begin{equation}
\label{eq:ch24-integrability-estimate}
 0\leq\int(-u_j)\,\MA_\theta(\varphi_1,\ldots,\varphi_n)
 \leq\int(-u_j)\,\theta^{\wedge n}+nMV
 \leq\int(-u)\,\theta^{\wedge n}+nMV .
\end{equation}

\emph{Step 3: the right-hand side is finite.}  The measure
\(\theta^{\wedge n}\) is a finite positive combination of Dirac masses at
divisorial points of a model determining \(\theta\)
(\Cref{prop:ch24-model-support}\,\cref{itm:ch24-ms:atomic}).  A
\(\theta\)-psh function is finite at every quasimonomial point, in
particular at every divisorial point, by
\Cref{prop:ch21-first-consequences}\,\cref{itm:ch21-fc:quasimonomial}; here
we use that \(u\) is not identically \(-\infty\), which is part of
\Cref{def:ch21-bfj-psh}.  Hence
\(\int(-u)\theta^{\wedge n}=\sum_Eb_E(\theta|_E^{\,n})(-u(x_E))<\infty\).

\emph{Step 4: passing to the limit.}  The functions \(-u_j\) are
nonnegative, lower semicontinuous, and increase to \(-u\); by
\Cref{prop:lsc-net-monotone} applied to the fixed Radon measure
\(\MA_\theta(\varphi_\bullet)\),
\[
 \begin{aligned}
 \int(-u)\,\MA_\theta(\varphi_1,\ldots,\varphi_n)
 &=\sup_j\int(-u_j)\,\MA_\theta(\varphi_1,\ldots,\varphi_n)\\
 &\leq\int(-u)\,\theta^{\wedge n}+nMV<\infty
 \end{aligned}
\]
by \eqref{eq:ch24-integrability-estimate}.  Thus \(u\) is integrable.

\emph{Step 5: pluripolar sets.}  Let \(E\subseteq X^{\An}\) be a Borel set
contained in \(\{v=-\infty\}\) for some \(v\in\PSH(X,\theta)\).  By Step~4
the function \(-v\) is \(\MA_\theta(\varphi_\bullet)\)-integrable, so the set
where it equals \(+\infty\) is null; hence
\(\MA_\theta(\varphi_\bullet)(E)=0\).
\end{proof}

\begin{corollary}
\label{cor:ch24-no-charge}
For bounded \(\theta\)-psh \(\varphi_1,\ldots,\varphi_n\) and
\(u\in\PSH(X,\theta)\),
\[
 \MA_\theta(\varphi_1,\ldots,\varphi_n)\bigl(\{u=-\infty\}\bigr)=0 ,
\]
and more generally every Borel pluripolar set is null for every mixed
Monge--Amp\`ere measure of bounded potentials.  Consequently
\(\MA_\theta(\varphi_1,\ldots,\varphi_n)\) does not charge the
\(G_\delta\) set \(\{u=-\infty\}\) of
\Cref{prop:ch21-polar-structure}\,\cref{itm:ch21-pol:gdelta}, and in
particular assigns no mass to any countable union of polar sets.
\end{corollary}

\begin{proof}
This is Step~5 of the preceding proof, followed by countable subadditivity
of the fixed measure for the last assertion.  Note that no outer-regularity
or capacity statement is used; the argument is a bare integrability
statement.  The corresponding capacity-theoretic assertion---that pluripolar
sets have capacity zero---is \Cref{prop:capacity-elementary}, which is
deduced from \Cref{prop:bounded-ma-no-polar-mass} in
\Cref{chap:bfj-capacity}.
\end{proof}

\begin{warning}{Nets, boundedness, and what monotone convergence means here}
\Cref{thm:bounded-ma-extension} concerns decreasing nets whose limits
remain bounded. Ordinary monotone convergence for measurable functions is
not valid for arbitrary nets. The weighted assertion
\eqref{eq:weighted-monotone-ma} is a theorem, not a formal use of the
sequential monotone convergence theorem. For an unbounded limit the bounded
operator is not used directly; truncation and finite energy enter in
\Cref{chap:bfj-energy}.

Three separate monotone statements have been used above and should be kept
apart.  (i) \Cref{lem:monotone-net-converges}: a decreasing net of
\emph{real numbers} converges to its infimum---no topology on \(X^{\An}\) is
involved, and this is what drives
\Cref{thm:ch24-energy-extension}.  (ii)
\Cref{cor:usc-net-monotone-integral}: a decreasing net of upper
semicontinuous \emph{functions}, integrated against one \emph{fixed} Radon
measure, has convergent integrals---this is used in Step~5 of
\Cref{thm:bounded-ma-extension} and in
\Cref{rem:ch24-forced-ibp-as-definition}.  (iii) The statement that a
decreasing net of functions integrated against a \emph{moving} net of
measures has convergent integrals is \emph{false} in general; only the
inequality \(\limsup\leq\) survives
(\Cref{lem:usc-along-decreasing-net}).  It is precisely because (iii) fails
that \eqref{eq:weighted-monotone-ma} cannot be proved by a portmanteau
argument, and that the proof passes through \(\mathfrak E\).
\end{warning}

\section{Locality and maximum principles}
\label{sec:ch24-locality}

\subsection{Locality: the model checkpoint}
\label{sec:ch24-model-locality}

The full bounded maximum principle requires capacity and is proved in
\Cref{thm:bounded-maximum-principle}.  Here we retain real model inputs,
but must distinguish them from rational model inputs: the maximum of
real model functions can be continuous without being a model function.
The measure of that maximum is already available from the bounded
construction and \Cref{prop:ch24-continuous-uniform}.

On the open set where \(\varphi\) strictly exceeds \(\psi\), the measure
of the maximum should see only \(\varphi\).  For rational potentials we
prove this by choosing a common SNC determination; the resulting
measures are finite sums of Dirac masses, and their coefficients are
intersection numbers, which can be real because \(\theta\) is real.
The geometric mechanism is that the coefficient at \(x_E\) depends only
on the divisor classes restricted to \(E\).  We then pass to real model
inputs by uniform approximation at the fixed form \(\theta\), testing on
compact subsets of the strict inequality locus.  Neither the bounded
locality theorem nor capacity is used in this argument.

\begin{proposition}\label{prop:model-maximum-principle}
For \(\theta\)-psh model functions \(\varphi,\psi\),
\begin{equation}
\label{eq:model-maximum-principle}
 \one_{\{\varphi>\psi\}}
 \MA_\theta(\max\{\varphi,\psi\})
 =\one_{\{\varphi>\psi\}}\MA_\theta(\varphi).
\end{equation}
\end{proposition}

\begin{proof}
The inputs may be real model functions.  Their maximum
\(\sigma:=\max\{\varphi,\psi\}\) is continuous and \(\theta\)-psh,
but need not be a model function.  Its measure is the bounded product of
\Cref{thm:bounded-ma-extension}, restricted to continuous potentials as
in \Cref{prop:ch24-continuous-uniform}.  We first prove the formula when
the two potentials are rational; the reference form \(\theta\) remains
the given real form throughout.

\emph{Step 1: rational inputs on one common SNC model.}
Suppose \(\varphi,\psi\) are rational model functions.  Then \(\sigma\)
is a rational \(\theta\)-psh model function by
\Cref{prop:ch20-model-max}.  Choose an SNC model \(\cX\) determining
\(\theta\) and all three potentials, using \Cref{prop:snc-cofinal}.
By \Cref{lem:model-function-affine}, all three factor through
\(p_{\cX}\) and are affine on every face of \(\Delta_{\cX}\).
The two measures are supported at its vertices by
\eqref{eq:model-mixed-measure}.  We therefore compare their coefficients
at each \(x_E\) with \(\varphi(x_E)>\psi(x_E)\).

If another component \(F\) meets \(E\), consider an edge
\([x_E,x_F]\).  The restrictions of \(\varphi\), \(\psi\), and their
maximum \(\sigma\) are all affine.  The maximum of two affine functions
on a segment cannot be affine if their strict order reverses along that
segment.  Since \(\varphi(x_E)>\psi(x_E)\), it follows that
\(\varphi(x_F)\geq\psi(x_F)\) and \(\sigma=\varphi\) on the whole
edge.  Write \(\sigma=\varphi_{D_\sigma}\) and
\(\varphi=\varphi_D\), with vertical rational Cartier divisors on
\(\cX\).  The coefficient of any component \(G\) in \(D\) is
\(b_G\varphi(x_G)\), by
\Cref{prop:metrics-divisor-model-function}; the same formula holds for
\(D_\sigma\).  Thus \(D_\sigma-D\) has zero coefficient along \(E\)
and along every component meeting \(E\).  Its support is disjoint from
\(E\).

Choose a positive integer \(m\) clearing the denominators of these
divisors.  The Cartier divisor \(m(D_\sigma-D)\) vanishes on a
neighborhood of \(E\), so its associated line bundle restricts trivially
to \(E\).  Dividing its numerical class by \(m\) gives
\([D_\sigma-D]|_E=0\) in \(N^1(E)_{\mathbf R}\).  Consequently
\[
 (\theta+\ddc\sigma)|_E=(\theta+\ddc\varphi)|_E,
\]
and the two coefficients
\(b_E\bigl((\theta+\ddc\sigma)|_E\bigr)^n\) and
\(b_E\bigl((\theta+\ddc\varphi)|_E\bigr)^n\) in
\eqref{eq:model-mixed-measure} agree.  Summing over all strict vertices
proves \eqref{eq:model-maximum-principle} for rational inputs, without
any rationality assumption on \(\theta\).

\emph{Step 2: approximation at the fixed real form.}
For real \(\theta\)-psh model inputs, apply
\Cref{lem:ch20-fixed-form-rational-approximation} to obtain rational
\(\theta\)-psh model functions \(\varphi_m,\psi_m\) with
\[
 \varepsilon_m:=\max\{\|\varphi_m-\varphi\|_\infty,
                         \|\psi_m-\psi\|_\infty\}\longrightarrow0.
\]
The ample generic class needed for this approximation is part of
\eqref{eq:ch24-standing}; the form itself is not perturbed.
Put \(\sigma_m=\max\{\varphi_m,\psi_m\}\), a rational
\(\theta\)-psh model function.  The elementary maximum inequality gives
\(\|\sigma_m-\sigma\|_\infty\leq\varepsilon_m\).
In particular \(\sigma\) is \(\theta\)-psh by uniform stability,
\Cref{thm:ch21-stability}\,\cref{itm:ch21-stab:unif}.
By \Cref{prop:ch24-continuous-uniform}\,\cref{itm:ch24-cu:unif},
\[
 \MA_\theta(\sigma_m)\longrightarrow\MA_\theta(\sigma),
 \qquad
 \MA_\theta(\varphi_m)\longrightarrow\MA_\theta(\varphi)
 \quad\text{weakly}.
\]

\emph{Step 3: localization before taking the limit.}
The set \(U=\{\varphi>\psi\}\) is open.  Let \(h\) be a continuous
function on \(X^{\An}\) whose support \(K\) is a compact subset of
\(U\).  If \(K\) is nonempty, continuity and compactness give
\(\delta:=\min_K(\varphi-\psi)>0\).  For every sufficiently large
\(m\), we have \(2\varepsilon_m<\delta\), hence
\(K\subseteq\{\varphi_m>\psi_m\}\).  The rational formula from
Step~1 therefore gives
\[
 \int h\,\MA_\theta(\sigma_m)
       =\int h\,\MA_\theta(\varphi_m).
\]
The equality is automatic if \(K\) is empty.  Weak convergence from
Step~2 now yields
\(\int h\,\MA_\theta(\sigma)=\int h\,\MA_\theta(\varphi)\).
The open subspace \(U\) of the compact Hausdorff space \(X^{\An}\)
is locally compact Hausdorff.  Every continuous compactly supported
function on \(U\) extends by zero to a continuous function of the type
just used.  Uniqueness of Radon measures from these tests shows that the
two measures restricted to \(U\) agree, proving
\eqref{eq:model-maximum-principle}.  We have not passed masses of moving
strict-inequality sets through a weak limit.
\end{proof}

\begin{corollary}
\label{cor:ch24-model-locality-open}
Let \(\varphi,\psi\) be \(\theta\)-psh model functions and let
\(U\subseteq X^{\An}\) be open.
\begin{enumerate}
\item\label{itm:ch24-mlo:agree} If \(\varphi=\psi\) on \(U\), then
 \(\MA_\theta(\varphi)|_U=\MA_\theta(\psi)|_U\).
\item\label{itm:ch24-mlo:mixed} The same holds for mixed measures: if
 \(\varphi_i=\psi_i\) on \(U\) for all \(i\), then
 \(\MA_\theta(\varphi_1,\ldots,\varphi_n)|_U
   =\MA_\theta(\psi_1,\ldots,\psi_n)|_U\).
\end{enumerate}
\end{corollary}

\begin{proof}
We prove \cref{itm:ch24-mlo:mixed}, of which \cref{itm:ch24-mlo:agree} is the
diagonal case.  Choose an SNC model \(\cX\) on which \(\theta\) and all
\(2n\) real model potentials are determined (\Cref{prop:snc-cofinal}).
Both measures are supported on the vertex set \(\{x_E\}\) of
\(\Delta_{\cX}\) by
\Cref{prop:ch24-model-support}\,\cref{itm:ch24-ms:atomic}, so it suffices
to show that the coefficients at each \(x_E\in U\) agree.

Fix such an \(E\).  The set \(U\cap\Delta_{\cX}\) is open in
\(\Delta_{\cX}\) and contains \(x_E\), so for every component \(F\)
meeting \(E\) it contains an initial segment of an edge
\([x_E,x_F]\) of positive length.  By
\Cref{lem:model-function-affine}, \(\varphi_i\) and \(\psi_i\) are
affine on that edge.  Since they agree on an initial segment, they agree
on the whole edge.  Thus \(\varphi_i(x_F)=\psi_i(x_F)\), and also
\(\varphi_i(x_E)=\psi_i(x_E)\).

Let \(D_i,D_i'\) be the vertical real Cartier divisors determining these
potentials.  By \Cref{prop:metrics-divisor-model-function}, their
coefficients agree along \(E\) and every component meeting \(E\).
Their difference is therefore supported away from \(E\).  On the
regular SNC model each component \(G\) is a Cartier divisor.  Writing
\(D_i-D_i'=\sum_G c_{iG}G\), every term with nonzero coefficient is
disjoint from \(E\), so \(\cO_{\cX}(G)|_E\) is trivial.  Taking the
real linear combination of these numerical classes gives
\([D_i-D_i']|_E=0\) in \(N^1(E)_{\mathbf R}\).  This is a statement
about a real numerical class, not an assertion that an arbitrary real
divisor has an associated honest line bundle.  Hence
\[
 (\theta+\ddc\varphi_i)|_E=(\theta+\ddc\psi_i)|_E
 \qquad\text{for every }i.
\]
The coefficients of \(\delta_{x_E}\) in the two mixed measures
\eqref{eq:model-mixed-measure} are \(b_E\) times the products of these
restrictions, and are consequently equal.
\end{proof}

\begin{remark}[The mass on the contact set is genuinely new]
\label{rem:ch24-contact-mass}
The strict-locus formula \eqref{eq:model-maximum-principle} does not
determine the measure on the contact set \(\{\varphi=\psi\}\).
Nor does \Cref{cor:ch24-model-locality-open} apply merely because the
potentials agree on a set that is not open.  The measure
\(\MA_\theta(\max\{\varphi,\psi\})\) can have an atom on the contact set
that is present in neither input measure.
\Cref{ex:ch24-locality-check} exhibits this on a chain of three
components, where a new atom of mass \(1\) appears at the crossing point.
The maximum is continuous even when it is not model, and its total
mass is still \(V\) by \Cref{thm:bounded-ma-extension}.
Thus the new atom represents redistribution of mass, not an increase
of total mass.
\end{remark}

\subsection{The comparison principle on a model}
\label{sec:ch24-model-comparison}

The comparison principle is the first genuinely nonlinear consequence of
locality, and it is what makes uniqueness statements possible.  We prove the
model version here; the bounded version is
\Cref{cor:bounded-comparison-principle}, deduced in
\Cref{chap:bfj-capacity} from the bounded locality theorem by the same
argument.

\begin{proposition}
\label{prop:ch24-model-comparison}
For \(\theta\)-psh model functions \(\varphi,\psi\),
\begin{equation}
\label{eq:ch24-model-comparison}
 \int_{\{\varphi<\psi\}}\MA_\theta(\psi)
 \leq
 \int_{\{\varphi<\psi\}}\MA_\theta(\varphi).
\end{equation}
\end{proposition}

\begin{proof}
Fix \(\varepsilon>0\) and put
\(\sigma_\varepsilon:=\max\{\varphi,\psi-\varepsilon\}\).
The shifted input \(\psi-\varepsilon\) is still a real \(\theta\)-psh
model function, while \(\sigma_\varepsilon\) is continuous and
\(\theta\)-psh and need not be a model function.  Its bounded
Monge--Amp\`ere measure is positive and has mass \(V\), by
\Cref{thm:bounded-ma-extension}.  The two sets
\[
 U:=\{\varphi<\psi-\varepsilon\},
 \qquad W:=\{\varphi>\psi-\varepsilon\}
\]
are disjoint.  The full real-input formula
\eqref{eq:model-maximum-principle}, applied to
\((\psi-\varepsilon,\varphi)\) and
\((\varphi,\psi-\varepsilon)\), gives
\[
 \MA_\theta(\sigma_\varepsilon)|_U
       =\MA_\theta(\psi-\varepsilon)|_U
       =\MA_\theta(\psi)|_U,
 \qquad
 \MA_\theta(\sigma_\varepsilon)|_W
       =\MA_\theta(\varphi)|_W.
\]
Here adding a constant does not change \(\ddc\) or the model measure.
Discarding the nonnegative mass of \(\MA_\theta(\sigma_\varepsilon)\)
on the complement of \(U\cup W\) yields
\[
 V\geq\int_U\MA_\theta(\psi)+\int_W\MA_\theta(\varphi).
\]
Since \(\MA_\theta(\varphi)\) also has mass \(V\), subtraction gives
\[
 \int_{\{\varphi<\psi-\varepsilon\}}\MA_\theta(\psi)
 \leq\int_{\{\varphi\leq\psi-\varepsilon\}}\MA_\theta(\varphi)
 \leq\int_{\{\varphi<\psi\}}\MA_\theta(\varphi).
\]
Along \(\varepsilon=1/m\), the sets on the left increase to
\(\{\varphi<\psi\}\).  Continuity from below for the one fixed finite
measure \(\MA_\theta(\psi)\) gives
\eqref{eq:ch24-model-comparison}.  No model determination of
\(\sigma_\varepsilon\) is used.
\end{proof}

\begin{corollary}
\label{cor:ch24-model-uniqueness-curve}
Let \(n=1\), let \(\cX\) be a semistable model of the curve \(X\), and let
\(\varphi,\psi\) be \(\theta\)-psh model functions determined on \(\cX\)
with \(\MA_\theta(\varphi)=\MA_\theta(\psi)\).  Then \(\varphi-\psi\) is
constant.
\end{corollary}

\begin{proof}
Write \(f:=\varphi-\psi\), a model function determined on \(\cX\).  By
\eqref{eq:model-mixed-measure} with \(n=1\), the hypothesis says that the
signed measure \(\ddc f\) vanishes; by
\Cref{prop:ch20-curve-laplacian} this means that the weighted graph
Laplacian of \(f|_{\Delta_{\cX}}\) vanishes at every vertex,
\[
 \sum_{F\neq E}b_F(E\cdot F)\bigl(f(x_F)-f(x_E)\bigr)=0
 \qquad\text{for every }E .
\]
Multiplying by \(f(x_E)\) and summing gives, by
\eqref{eq:ch20-hodge-laplacian},
\(\sum_{E<F}w_{EF}\bigl(f(x_E)-f(x_F)\bigr)^2=0\) with
\(w_{EF}=b_Eb_F(E\cdot F)\geq0\); since the dual graph of \(\cX_s\) is
connected (\Cref{thm:ch20-ddc-exact}, Step~1) and \(w_{EF}>0\) exactly when
\(E\cap F\neq\varnothing\), all the values \(f(x_E)\) are equal.  A model
function determined on \(\cX\) is affine on each face of
\(\Delta_{\cX}\) and factors through \(p_{\cX}\)
(\Cref{lem:model-function-affine}), so \(f\) is constant on \(X^{\An}\).
\end{proof}

\begin{remark}[Uniqueness in higher dimension]
\label{rem:ch24-uniqueness-scope}
\Cref{cor:ch24-model-uniqueness-curve} is a statement about the
\emph{linear} equation \(\ddc f=0\) and is available only because for
\(n=1\) the Monge--Amp\`ere operator is linear.  For \(n\geq2\) uniqueness of
solutions of \(\MA_\theta(\varphi)=\mu\) modulo constants is a genuine
theorem, proved in \Cref{chap:bfj-continuity} for the continuous solutions
of the original BFJ theorem and in \Cref{chap:bfj-finite-energy-measures} in
the finite-energy setting; the mechanism there is the domination principle,
which uses capacity.  Nothing in the present chapter proves uniqueness in
dimension \(\geq2\), and \eqref{eq:ch24-model-comparison} alone does not
suffice.
\end{remark}

\subsection{Maximum principles}
\label{sec:ch24-max}

Three statements go under the name of a maximum principle in this theory,
and they have different depths.

\begin{proposition}
\label{prop:ch24-skeletal-maximum}
Let \(\cX\) be an SNC model on which \(\theta\) is determined.
\begin{enumerate}
\item\label{itm:ch24-sm:model} If \(f\) is a model function determined on
 \(\cX\), then \(\max_{X^{\An}}f=\max_{E}f(x_E)\) and
 \(\min_{X^{\An}}f=\min_Ef(x_E)\), both attained at divisorial points of
 \(\cX\).
\item\label{itm:ch24-sm:psh} If \(u\in\PSH(X,\theta)\) and \(f\) is a model
 function determined on \(\cX\), then
 \(\sup_{X^{\An}}(u-f)=\max_{\Delta_{\cX}}(u-f)\).
\item\label{itm:ch24-sm:supp} If \(\varphi\) is a \(\theta\)-psh model
 function determined on \(\cX\) and \(u\in\PSH(X,\theta)\), then
 \(\sup_{X^{\An}}(u-\varphi)\) is attained at a point of \(\Delta_{\cX}\),
 and \(\Supp\MA_\theta(\varphi)\subseteq\Delta_{\cX}\).
\end{enumerate}
\end{proposition}

\begin{proof}
\cref{itm:ch24-sm:model} is \Cref{prop:ch20-maximum-principle}.  For
\cref{itm:ch24-sm:psh}, \(f=f\circ p_{\cX}\) by
\Cref{lem:model-function-affine} and \(u\leq u\circ p_{\cX}\) by
\eqref{eq:ch21-retraction-inequality}, so
\(u-f\leq(u-f)\circ p_{\cX}\); hence the supremum over \(X^{\An}\) is the
supremum over \(\Delta_{\cX}\), which is attained because
\(u|_{\Delta_{\cX}}\) is continuous
(\Cref{prop:ch21-first-consequences}\,\cref{itm:ch21-fc:finite}) and
\(\Delta_{\cX}\) is compact.  This is
\Cref{prop:ch21-sup-on-skeleton} applied to \(u-f\), which is
\((\theta+\ddc f)\)-psh.  Statement \cref{itm:ch24-sm:supp} is
\cref{itm:ch24-sm:psh} together with
\Cref{prop:ch24-model-support}\,\cref{itm:ch24-ms:atomic}.
\end{proof}

\begin{remark}[The domination principle is not proved here]
\label{rem:ch24-domination-scope}
The strongest statement of this type---if \(\varphi,\psi\) are bounded
\(\theta\)-psh and \(\varphi\leq\psi\) holds
\(\MA_\theta(\psi)\)-almost everywhere, then \(\varphi\leq\psi\)
everywhere---is the \emph{domination principle}.  It is not proved in this
chapter.  It requires the comparison principle for \emph{bounded}
potentials, hence quasicontinuity, and is
\Cref{prop:bounded-capacity-domination} in \Cref{chap:bfj-capacity}.  We
flag this because \Cref{prop:ch24-skeletal-maximum} looks superficially like
a domination statement and is not one: it locates a supremum on a skeleton,
which is a consequence of the retraction inequality alone and carries no
information about \(\Supp\MA_\theta(\psi)\).
\end{remark}

\subsection{What needs capacity, and what does not}
\label{sec:ch24-needs-capacity}

The following table records, for the four statements of this section, the
model version proved here and the bounded version proved later, together
with the extra input.

\begin{center}
\small
\begin{tabular}{p{.25\textwidth}|p{.21\textwidth}|p{.21\textwidth}|p{.24\textwidth}}
\toprule
\textbf{Statement}&\textbf{Model level}&\textbf{Bounded level}
 &\textbf{Extra input needed}\\
\midrule
locality of \(\MA\) on \(\{\varphi>\psi\}\)
 &\Cref{prop:model-maximum-principle}
 &\Cref{thm:bounded-maximum-principle}
 &quasicontinuity, \Cref{thm:bounded-quasicontinuity}\\
agreement on an open set
 &\Cref{cor:ch24-model-locality-open}
 &\Cref{thm:bounded-maximum-principle}, last sentence
 &as above\\
comparison principle
 &\Cref{prop:ch24-model-comparison}
 &\Cref{cor:bounded-comparison-principle}
 &bounded locality only\\
domination principle
 &---
 &\Cref{prop:bounded-capacity-domination}
 &capacity estimate \Cref{lem:fundamental-capacity-estimate}\\
\bottomrule
\end{tabular}
\end{center}

The reason the model level does not already give the bounded level is that
the set \(\{\varphi>\psi\}\) is open when \(\varphi,\psi\) are continuous but
merely a difference of upper semicontinuous functions in general; weak
convergence of measures says nothing about the mass of such a set.
Quasicontinuity replaces it by an open set up to small capacity, and this is
the whole content of the passage from
\Cref{prop:model-maximum-principle} to
\Cref{thm:bounded-maximum-principle}.

\section{Comparison with the Chambert--Loir and CLD constructions}
\label{sec:ch24-comparison}

\subsection{The three model-level measures}
\label{sec:ch24-three-measures}

Let \(L\) be ample, choose a model metric \(\|\cdot\|_0\) with curvature
\(\theta\), and write a continuous semipositive metric as
\(\|\cdot\|_\varphi=\rme^{-\varphi}\|\cdot\|_0\).

\begin{theorem}\label{thm:bfj-cld-ma-common-locus}
For semipositive model metrics, the intersection measure
\eqref{eq:model-mixed-measure}, the Chambert--Loir measure, and the CLD
product of the corresponding formal Chern currents agree. For continuous
semipositive metrics, the BFJ measure \(\MA_\theta(\varphi)\) agrees with the
Chambert--Loir measure obtained by uniform approximation by semipositive
model metrics.
\end{theorem}

\begin{proof}
The theorem is the specialization to the present notation of
\Cref{thm:ch19-cld-bfj-comparison}, proved in \Cref{chap:cld-examples}; we
indicate how the two statements match and which parts of that proof are used
here, so that no circularity arises.

\emph{The model case} is
\Cref{thm:ch19-cld-bfj-comparison}\,(A), and specifically the chain of
identities \eqref{eq:ch19-cld-bfj-model}.  In one line: for a model metric
the coefficient at \(x_E\) in all three constructions is
\(b_E\bigl(c_1(\mathcal L)|_E\bigr)^n\).  That this is the BFJ coefficient is
\eqref{eq:model-mixed-measure}; that it is the Chambert--Loir coefficient is
the defining formula of \cite{CL06}, recorded here as
\Cref{cor:ch24-cl-identification}; and that it is the mass of the CLD
Bedford--Taylor product of the corresponding formal Chern currents is
\Cref{thm:cld-recovers-cl}, whose proof rests on the formal intersection
formula \cite[Theorem~18.5.5 and Proposition~18.5.8]{CLD25}.  The
identification of the divisorial normalizations on the two sides is
\Cref{rem:divisorial-normalization}.  All of this is purely algebraic: no
statement of the present chapter is used, so the argument may be quoted here
without circularity.

\emph{The continuous case} is
\Cref{thm:ch19-cld-bfj-comparison}\,(B), and it is here that the analysis of
the present chapter is used.  The mechanism is the reduction of uniform
approximation to \emph{decreasing} uniform approximation.  Choose model
weights converging uniformly to \(\varphi\), and pass to a subsequence with
\(\|\varphi_m-\varphi\|_\infty\leq\varepsilon_m\), where
\(\sum_m\varepsilon_m<\infty\). Put
\[
 d_m:=\sum_{j\geq m}(\varepsilon_j+\varepsilon_{j+1}),
 \qquad \psi_m:=\varphi_m+d_m.
\]
Then \(\psi_m\) is a semipositive model weight and
\(\psi_m\searrow\varphi\), since
\(\|\varphi_m-\varphi_{m+1}\|_\infty
 \leq\varepsilon_m+\varepsilon_{m+1}\).
BFJ monotone continuity and Chambert--Loir's continuity under uniform
semipositive approximation identify both limits with the weak limit of the
same model measures \(\MA_\theta(\psi_m)\).  On the BFJ side this is
\Cref{prop:ch24-continuous-uniform}\,\cref{itm:ch24-cu:seq} together with
\Cref{cor:ch24-continuous-cl}; on the CLD side it is
\Cref{thm:cld-bt-continuity}, whose hypotheses require exactly the local
boundedness and uniform convergence supplied by the decreasing sequence
\((\psi_m)\).  The mixed case follows either by repeating the argument with
different potentials in different slots or, more economically, from the
diagonal case by the polarization identity
\eqref{eq:ch24-polarization} and \Cref{cor:ch24-mixed-from-diagonal}.
\end{proof}

\subsection{Exact scope of the comparison}
\label{sec:ch24-scope-comparison}

\begin{remark}[Boundary of the comparison]
\label{rem:ch24-comparison-boundary}
The theorem does not identify every bounded global BFJ-psh function with a
locally CLD-psh function. Nor does it identify the BFJ monotone product with
a CLD Bedford--Taylor current product outside the common regularizable
locus. Boucksom--Eriksson give broader comparisons under their semiample and
CLD hypotheses \cite[Theorems~7.14 and~8.4]{BE21}; these do not prove an
unrestricted singular BFJ--CLD equivalence.
\end{remark}

It is worth spelling out what \Cref{rem:ch24-comparison-boundary} excludes,
because the temptation to extend \Cref{thm:bfj-cld-ma-common-locus} by
``continuity'' is strong and each of the four possible extensions fails for a
different reason.  The list below is
\Cref{thm:ch19-cld-bfj-comparison}\,(C), repeated here with the reasons
attached; the limitations table of \Cref{sec:ch19-limitations} is the full
account.

\begin{warning}{Four statements this book does not prove}
\begin{enumerate}
\item\label{itm:ch24-notproved:psh}
 That a bounded \(\theta\)-psh function is locally CLD-psh
 (\Cref{thm:ch19-cld-bfj-comparison}, item~\cref{itm:ch19-cb:scope1}).  The
 model identification propagates only along \emph{uniformly} convergent
 sequences, whereas a bounded \(\theta\)-psh function is only a decreasing
 pointwise limit of model functions and need not be continuous; see
 \Cref{ex:ch24-branch-sum} for such a function.
\item\label{itm:ch24-notproved:converse}
 That a locally CLD-psh function on \(X^{\An}\) is \(\theta\)-psh for some
 closed model form \(\theta\) (item~\cref{itm:ch19-cb:scope2} there).  The
 CLD class is a sheaf and the BFJ class is defined by a global condition on
 all models; \Cref{prop:ch21-bfj-vs-cld} isolates the two structural
 differences.
\item\label{itm:ch24-notproved:unbounded}
 That the BFJ monotone product of an \emph{unbounded} \(\theta\)-psh
 function equals a CLD Bedford--Taylor product
 (item~\cref{itm:ch19-cb:scope3} there).  \Cref{thm:cld-bt-continuity}
 requires local boundedness, and the BFJ extension beyond bounded
 potentials is by a different mechanism altogether, the nonpluripolar
 convention \Cref{def:nonpluripolar-ma}.
\item\label{itm:ch24-notproved:polar}
 That the two constructions assign the same mass to a pluripolar set
 (item~\cref{itm:ch19-cb:scope4} there).  This is not a gap but a genuine
 difference: \Cref{ex:ch16-polar-failure} exhibits a CLD product carrying
 mass on a polar set, while \Cref{prop:bounded-ma-no-polar-mass} shows the
 bounded BFJ product never does, and
 \Cref{def:nonpluripolar-ma} discards such mass by construction.
\end{enumerate}
\end{warning}

\begin{remark}[Which product the comparison concerns]
\label{rem:ch24-which-product}
\Cref{thm:bfj-cld-ma-common-locus} concerns the \emph{model} product in its
first sentence and the \emph{continuous} product in its second.  It says
nothing about the bounded product of \Cref{thm:bounded-ma-extension} beyond
the continuous locus, and nothing at all about the finite-energy product of
\Cref{chap:bfj-energy}.  On the CLD side it concerns the Bedford--Taylor
product of \Cref{def:ch16-bt-current} for continuous psh weights, and not the
extension to the wider classes of \Cref{def:ch16-psh-filtration}.  The
three-level delimitation on the CLD side is
\Cref{sec:ch16-three-levels}, and the corresponding table is
\Cref{sec:ch16-delimitation-table}; the summary of everything that
\emph{is} transferable is \Cref{cor:ch19-transfer}, and the synthesis at the
end of \Cref{part:bfj} is
\Cref{thm:bfj-cld-common-locus-synthesis}.
\end{remark}

\begin{remark}[Delta-forms and the third construction]
\label{rem:ch24-delta-forms}
On the model locus a fourth object agrees with the three of
\Cref{thm:bfj-cld-ma-common-locus}: the product of Chern delta-forms of
\Cref{chap:delta-forms}, by
\Cref{thm:gk-delta-cl-measure}.  We mention it because it is the
construction under which positivity of a formal metric is most easily
recognized (\Cref{thm:ch19-cld-bfj-comparison},
item~\cref{itm:ch19-cb:delta}), but it plays no role in
\Cref{part:bfj} and no statement of this chapter depends on it.
\end{remark}

\section{Worked examples}
\label{sec:ch24-examples}

The six computations below test every normalization introduced in this
chapter: the multiplicity \(b_E\), the sign of \(\ddc\), the total mass, the
polarization identity, the monotone extension, and the locality identity.
Each is carried to explicit numbers, and each is cross-checked against an
independent computation wherever one is available.

\subsection{Mixed measures on a curve}
\label{sec:ch24-ex-curve}

For \(n=1\) the Monge--Amp\`ere operator has only one slot, so ``mixed''
means multilinear in the metrized line bundle rather than in the potential.
That is nevertheless a genuine test: it checks the additivity of
\eqref{eq:model-mixed-measure} in the classes, the additivity of masses, and
the symmetry of the energy pairing \eqref{eq:ch24-mixed-energy}, all against
the graph Laplacian of \Cref{chap:model-calculations}.  We begin with the
running two-component curve.

\begin{example}[A semistable curve with two components]
\label{ex:two-vertex-ma}
Let a regular semistable model have special fibre \(E_0+E_1\), with one
node rational over the residue field, and let the reference measure be
\(\theta=m_0\delta_{x_0}+m_1\delta_{x_1}\), where
\(m_0+m_1=\deg L\). Give the edge length one and let a model function have
vertex values \(a_0,a_1\). With the graph Laplacian convention forced by
\(\ddc\varphi_D=[D]\),
\[
 \ddc\varphi=(a_1-a_0)\delta_{x_0}+(a_0-a_1)\delta_{x_1}.
\]
Hence
\[
 \MA_\theta(\varphi)
 =(m_0+a_1-a_0)\delta_{x_0}
 +(m_1+a_0-a_1)\delta_{x_1}.
\]
The metric is semipositive precisely when both coefficients are nonnegative,
and their sum is always \(\deg L\).
\end{example}

The last sentence is \eqref{eq:ch24-mass} in dimension one: the mass
\(\deg L\) is blind to the potential, as
\Cref{cor:ch24-mass-independent} asserts, and the two nonnegativity
conditions are the two componentwise inequalities
\eqref{eq:ch20-mixed-component-positive}.  Explicitly, with
\(m_0=m_1=1\) and \(\deg L=2\), the admissible pairs \((a_0,a_1)\) are those
with \(|a_0-a_1|\leq1\), which is the nef interval
\eqref{eq:running-nef-interval} of the running conic; at the two extremes the
whole mass sits at one vertex.

\begin{example}[Two metrized bundles on the same curve: bilinearity and the
energy pairing]
\label{ex:ch24-curve-mixed}
Keep the two-component model of \Cref{ex:two-vertex-ma} with
\((E_0^2)=(E_1^2)=-1\), \((E_0\cdot E_1)=1\), both multiplicities one, and
edge length and residue weight one.  Let \(L_1,L_2\) be line bundles on \(X\) of degrees
\(d_1=2\) and \(d_2=3\), extended to model line bundles
\(\cL_1,\cL_2\) with component degrees
\[
 \bigl(\deg\cL_1|_{E_0},\deg\cL_1|_{E_1}\bigr)=(1,1),
 \qquad
 \bigl(\deg\cL_2|_{E_0},\deg\cL_2|_{E_1}\bigr)=(2,1).
\]
Write \(\theta_1,\theta_2\) for the corresponding closed model forms.  Then
\[
 \theta_1=\delta_{x_0}+\delta_{x_1},
 \qquad
 \theta_2=2\delta_{x_0}+\delta_{x_1},
\]
of masses \(2=d_1\) and \(3=d_2\), as \eqref{eq:ch24-mass} requires.

\emph{Bilinearity.}  The model line bundle \(\cL_1\otimes\cL_2\) extends
\(L_1\otimes L_2\), of degree \(5\), with component degrees \((3,2)\), so
\[
 \theta_1+\theta_2=3\delta_{x_0}+2\delta_{x_1},
\]
of mass \(5=d_1+d_2\).  This is additivity of
\eqref{eq:model-mixed-measure} in the classes, and it is the case \(n=1\) of
the multilinearity in \Cref{prop:model-ma-basic}.

\emph{The energy pairing and its symmetry.}  Let \(\varphi=\varphi_D\) and
\(\psi=\varphi_{D'}\) be model functions with vertex values \((a_0,a_1)\) and
\((c_0,c_1)\); thus \(D=a_0E_0+a_1E_1\) and \(D'=c_0E_0+c_1E_1\).  Using
\eqref{eq:model-test-intersection} and
\(\ddc\varphi=(a_1-a_0)(\delta_{x_0}-\delta_{x_1})\),
\[
 \int\psi\,\ddc\varphi
 =(c_0-c_1)(a_1-a_0)
 =\int\varphi\,\ddc\psi ,
\]
which is \eqref{eq:model-ibp}, and directly
\((D\cdot D')=-(a_0-a_1)(c_0-c_1)\) from the intersection matrix
\eqref{eq:ch20-two-component-matrix}.  With \(\theta=\theta_1\), the mixed
energy pairing of \Cref{def:ch24-mixed-energy} is, by
\eqref{eq:ch24-mixed-energy} with \(n=1\),
\[
 \mathfrak E(\varphi,\psi)
 =(D\cdot\theta_1)+(D'\cdot\theta_1)+(D\cdot D')
 =(a_0+a_1)+(c_0+c_1)-(a_0-a_1)(c_0-c_1),
\]
manifestly symmetric, affine in each variable, and satisfying
\(\mathfrak E(\varphi+t,\psi+s)=\mathfrak E(\varphi,\psi)+2(t+s)\), which is
\eqref{eq:ch24-energy-constants} with \(V=2\).  Monotonicity
(\Cref{prop:ch24-energy-basic}\,\cref{itm:ch24-eb:mono}) is visible: for
\(\varphi\) fixed, the derivative of \(\mathfrak E\) in \(c_0\) is
\(1-(a_0-a_1)\) and in \(c_1\) is \(1+(a_0-a_1)\), and these are exactly the
two coefficients of \(\MA_{\theta_1}(\varphi)\), both nonnegative precisely
when \(\varphi\) is \(\theta_1\)-psh.  This is
\eqref{eq:ch24-energy-slot} written out.

\emph{A numerical instance.}  For \((a_0,a_1)=(0,-1)\) and
\((c_0,c_1)=(0,1)\) one gets
\(\MA_{\theta_1}(\varphi)=0\cdot\delta_{x_0}+2\delta_{x_1}\),
\(\int\psi\,\MA_{\theta_1}(\varphi)=2\), and
\(\mathfrak E(\varphi,\psi)=(-1)+(1)-(1)(-1)=1\); the check
\(\mathfrak E(\varphi,\psi)-\mathfrak E(\varphi,0)
 =1-(-1)=2=\int\psi\,\MA_{\theta_1}(\varphi)\) confirms
\eqref{eq:ch24-energy-slot}.
\end{example}

\subsection{A mixed measure on a toric surface}
\label{sec:ch24-ex-toric}

\begin{example}[A genuinely mixed Chambert--Loir measure on
\(\mathbf P^2\)]
\label{ex:ch24-toric-surface-mixed}
We continue \Cref{ex:metrics-cl-toric-surface}, now computing a
\emph{mixed} measure and checking it twice: once by intersection numbers on
the two components of a special fibre, once by mixed volumes of subgradient
polytopes.

\emph{Setup.}  Let \(K\) be discretely valued with \(v(\pi)=1\), let
\(X=\mathbf P^2_K\) with torus coordinates \(x=Z_1/Z_0\), \(y=Z_2/Z_0\), and
let \(L=\cO(1)\), so \((L^2)=1\).  Write \(u=(v(x),v(y))\) for the
tropical coordinate \eqref{eq:toric-trop}.  Consider two invariant model
metrics on \(L\):
\[
 F_1(u)=\max\{0,-u_1,-u_2\},
 \qquad
 F_2(u)=\max\{-1,-u_1,-u_2\},
\]
the first the standard model metric determined by \(\mathbf P^2_R\), the
second the Fubini--Study metric of \eqref{eq:metrics-toric-weight}
determined by the sections \(\pi Z_0,Z_1,Z_2\).  Let \(\theta_1,\theta_2\)
be the corresponding closed model forms; both have generic class
\(c_1(\cO(1))\), so \((\{\theta_1\}\{\theta_2\})=(L^2)=1\).

\emph{The common model.}  Let \(\rho:\cX\to\mathbf P^2_R\) be the blowup of
the torus-fixed point \(p\) of the special fibre, with exceptional divisor
\(E\simeq\mathbf P^2_k\) and strict transform \(\widetilde{X_s}\), both of
multiplicity one:
\(\cX_s=\widetilde{X_s}+E\).  As in
\Cref{ex:metrics-cl-toric-surface}, \(\theta_1\) is determined by
\(\rho^*\cO(1)\) and \(\theta_2\) by
\(\rho^*\cO(1)\otimes\cO_{\cX}(-E)\).  Writing \(h\) for the pullback of the
hyperplane class to \(\widetilde{X_s}\) and \(e\) for its exceptional curve,
so \((h^2)=1\), \((h\cdot e)=0\), \((e^2)=-1\), the restrictions are
\[
 \theta_1|_{E}=0,\quad \theta_1|_{\widetilde{X_s}}=h;
 \qquad
 \theta_2|_{E}=\cO_{\mathbf P^2}(1),\quad
 \theta_2|_{\widetilde{X_s}}=h-e .
\]
The first vanishes because \(\rho\) contracts \(E\).

\emph{The three measures.}  By \eqref{eq:model-mixed-measure},
\[
 \theta_1\wedge\theta_1
 =1\cdot(h^2)\,\delta_{x_{\widetilde{X_s}}}+1\cdot0\,\delta_{x_E}
 =\delta_{x_0},
\]
where \(x_0=x_{\widetilde{X_s}}\) is the Gauss point, at \(u=(0,0)\);
\[
 \theta_2\wedge\theta_2
 =1\cdot\bigl((h-e)^2\bigr)\delta_{x_0}
  +1\cdot\bigl(\cO_{\mathbf P^2}(1)^2\bigr)\delta_{x_E}
 =0\cdot\delta_{x_0}+1\cdot\delta_{x_E}
 =\delta_{x_E},
\]
with \(x_E\) at \(u=(1,1)\), which is
\eqref{eq:metrics-toric-measure}; and the mixed measure is
\begin{equation}
\label{eq:ch24-toric-mixed}
 \theta_1\wedge\theta_2
 =1\cdot\bigl(h\cdot(h-e)\bigr)\delta_{x_0}
  +1\cdot\bigl(0\cdot\cO_{\mathbf P^2}(1)\bigr)\delta_{x_E}
 =\delta_{x_0}.
\end{equation}
All three have mass \(1\), in accordance with \eqref{eq:ch24-mass}:
\((\{\theta_1\}^2)=(\{\theta_1\}\{\theta_2\})=(\{\theta_2\}^2)=(L^2)=1\).
Note that the mixed measure coincides with \(\theta_1\wedge\theta_1\) and
not with \(\theta_2\wedge\theta_2\); mixed measures are not interpolations.

\emph{Cross-check by convex geometry.}  By
\Cref{thm:toric-ma-comparison} the pushforward under tropicalization of the
diagonal measure of an invariant model metric with weight \(F\) is
\(n!\operatorname{MA}_{\mathbf R}(F)=2\operatorname{MA}_{\mathbf R}(F)\),
where \(\operatorname{MA}_{\mathbf R}\) is the Alexandrov measure
\eqref{eq:alexandrov-ma}.  Applying this to \(F_1\), \(F_2\), and
\(F_1+F_2\)---the last being the weight of the tensor product metric on
\(\cO(2)\)---and using bilinearity of \(\wedge\), one obtains
\[
 \Trop_*\bigl(\theta_1\wedge\theta_2\bigr)
 =\tfrac12\Bigl[
   2\operatorname{MA}_{\mathbf R}(F_1+F_2)
   -2\operatorname{MA}_{\mathbf R}(F_1)
   -2\operatorname{MA}_{\mathbf R}(F_2)\Bigr].
\]
Now \(F_1\) is the maximum of three affine functions with slopes
\(0,-e_1^*,-e_2^*\) coinciding only at \(u=(0,0)\), where
\(\partial F_1=\operatorname{conv}\{0,-e_1,-e_2\}\) has area \(\tfrac12\);
so \(\operatorname{MA}_{\mathbf R}(F_1)=\tfrac12\delta_{(0,0)}\) and
\(2\operatorname{MA}_{\mathbf R}(F_1)=\delta_{(0,0)}\), matching
\(\theta_1\wedge\theta_1=\delta_{x_0}\).  Likewise
\(\operatorname{MA}_{\mathbf R}(F_2)=\tfrac12\delta_{(1,1)}\).  For the sum,
\(\partial(F_1+F_2)(u)=\partial F_1(u)+\partial F_2(u)\), so a point
contributes only if that Minkowski sum is two-dimensional.  Each
\(\partial F_i(u)\) is a point, a segment, or a triangle; the triangles occur
only at \(u=(0,0)\) for \(F_1\) and \(u=(1,1)\) for \(F_2\), and the two
segment loci that could combine---the sets where exactly two affine pieces
are active---are \(\{u_1=0,u_2>0\}\cup\{u_2=0,u_1>0\}\cup\{u_1=u_2\leq0\}\)
for \(F_1\) and \(\{u_1=1,u_2>1\}\cup\{u_2=1,u_1>1\}
\cup\{u_1=u_2\leq1\}\) for \(F_2\); they meet only along
\(\{u_1=u_2\leq0\}\), where both subdifferentials are the \emph{same}
segment \(\operatorname{conv}\{-e_1,-e_2\}\) and the sum is one-dimensional.
Hence only \((0,0)\) and \((1,1)\) contribute.  At \((0,0)\),
\[
 \partial F_1(0,0)=\operatorname{conv}\{0,-e_1,-e_2\},
 \qquad
 \partial F_2(0,0)=\operatorname{conv}\{-e_1,-e_2\},
\]
and the Minkowski sum, of area \(\tfrac32\), has vertices
\[
 (-2,0),\quad(-1,0),\quad(0,-1),\quad(0,-2).
\]
At \((1,1)\),
\(\partial F_1(1,1)=\{0\}\) and
\(\partial F_2(1,1)=\operatorname{conv}\{0,-e_1,-e_2\}\), of area
\(\tfrac12\).  Hence
\[
 \operatorname{MA}_{\mathbf R}(F_1+F_2)
 =\tfrac32\delta_{(0,0)}+\tfrac12\delta_{(1,1)}.
\]
Consequently,
\[
 \Trop_*(\theta_1\wedge\theta_2)
 =\tfrac12\bigl[(3-1-0)\delta_{(0,0)}+(1-0-1)\delta_{(1,1)}\bigr]
 =\delta_{(0,0)},
\]
which is \eqref{eq:ch24-toric-mixed} transported to \(N_{\mathbf R}\).  The
two computations---one by mixed degrees on the two components of a special
fibre, one by mixed areas of subgradient polytopes---agree, including the
factor \(n!=2\) of \Cref{lem:toric-real-total-mass}, and both give mass
\(1\).

\emph{Polarization check.}  \Cref{prop:ch24-polarization} with \(n=2\) reads
\(2\MA_\theta(\varphi_1,\varphi_2)
 =4\MA_\theta\bigl(\tfrac12(\varphi_1+\varphi_2)\bigr)
  -\MA_\theta(\varphi_1)-\MA_\theta(\varphi_2)\).
Here the potentials are the relative weights of the two metrics with respect
to a common reference, and
\(\tfrac12(F_1+F_2)\) is the weight of the metric on \(\cO(1)\) obtained as
the square root of the tensor product.  Substituting the three measures
computed above,
\[
 2\delta_{x_0}
 =4\cdot\tfrac14\bigl(3\delta_{x_0}+\delta_{x_E}\bigr)
  -\delta_{x_0}-\delta_{x_E}
 =3\delta_{x_0}+\delta_{x_E}-\delta_{x_0}-\delta_{x_E}
 =2\delta_{x_0},
\]
as required, where
\(\MA(\tfrac12(F_1+F_2))=\tfrac14(3\delta_{x_0}+\delta_{x_E})\) by
homogeneity of the real Monge--Amp\`ere operator.
\end{example}

\subsection{A bounded psh function that is not a model function}
\label{sec:ch24-ex-branch}

The next example is the smallest complete illustration of
\Cref{thm:bounded-ma-extension}: a bounded \(\theta\)-psh function that is
\emph{not} a model function, is not even continuous, and whose
Monge--Amp\`ere measure has infinitely many atoms---so that it is not a model
measure either---and yet is computed exactly, by a decreasing sequence.

\begin{example}[A branch sum on the Berkovich projective line]
\label{ex:ch24-branch-sum}
Let \(K\) be complete and discretely valued with \(v(\pi)=1\) and residue
field \(k\) of characteristic \(0\); in particular \(k\supseteq\mathbf Q\) is
infinite.  Let \(X=\mathbf P^1_K\), \(L=\cO(1)\), and let \(\theta\) be the
closed model form of the standard model \(\mathbf P^1_R\), so that by
\eqref{eq:p1-standard-measure}
\[
 \theta=\delta_{\zeta},\qquad \zeta:=\zeta_{0,1}
 \ \text{the Gauss point},\qquad V=(L^1)=1 .
\]
We use the tree structure of \(\mathbf P^{1,\An}_K\) and the path metric
\eqref{eq:ch08-path-metric} of \Cref{def:ch08-path-metric}.

\emph{The pieces.}  By \Cref{prop:ch08-branching} the tangent directions at
\(\zeta\) are in bijection with \(\mathbf P^1(k)\), an infinite set.  Choose
pairwise distinct residues \(\bar a_1,\bar a_2,\ldots\in k\) with lifts
\(a_k\in R\), put \(\ell_k:=2^k\), and let \(x_k\) be the point of
\(\mathbf P^{1,\An}_K\) at distance \(\ell_k\) from \(\zeta\) in the
direction \(\bar a_k\); concretely \(x_k=\zeta_{a_k,|\pi|^{\ell_k}}\), a
divisorial point, obtained by blowing up the ideal
\((\pi^{\ell_k},T-a_k)\).  Define
\[
 \psi_k:=-\frac1{\ell_k}\,
   \min\bigl\{\ell_k,\ \max\{0,v_\bullet(T-a_k)\}\bigr\}
   \ \in\ \mathcal D(X),
\]
so that \(\psi_k(\zeta)=0\), \(\psi_k\) decreases linearly with slope
\(-1/\ell_k\) along the segment \([\zeta,x_k]\), equals \(-1\) on the whole
branch beyond \(x_k\), and vanishes off that branch.  By the graph
Laplacian convention \eqref{eq:graph-laplacian}, forced by
\(\ddc\varphi_D=[D]\),
\[
 \ddc\psi_k=\frac1{\ell_k}\bigl(\delta_{x_k}-\delta_{\zeta}\bigr)
 =2^{-k}\bigl(\delta_{x_k}-\delta_{\zeta}\bigr),
\]
of total mass zero as \eqref{eq:ch19-lap-mass-zero} requires.  In
particular each \(\psi_k\) is \(\theta\)-psh, since
\(\theta+\ddc\psi_k=(1-2^{-k})\delta_\zeta+2^{-k}\delta_{x_k}\geq0\).

\emph{The function.}  Set
\[
 \varphi:=\sum_{k\geq1}\psi_k,
 \qquad
 \varphi_N:=\sum_{k=1}^{N}\psi_k .
\]
Because the branches are pairwise disjoint away from \(\zeta\), at every
point at most one summand is nonzero, so the sum converges pointwise and
\(-1\leq\varphi\leq0\).  Each \(\varphi_N\) is a \(\theta\)-psh model
function: it is a finite \(\mathbf Q\)-linear combination of model functions,
and
\[
 \theta+\ddc\varphi_N
 =\Bigl(1-\sum_{k\leq N}2^{-k}\Bigr)\delta_\zeta
  +\sum_{k\leq N}2^{-k}\delta_{x_k}
 =2^{-N}\delta_\zeta+\sum_{k\leq N}2^{-k}\delta_{x_k}\ \geq\ 0 .
\]
Moreover \(\varphi_N-\varphi_{N+1}=-\psi_{N+1}\geq0\), so
\(\varphi_N\searrow\varphi\) is a \emph{decreasing sequence} of
\(\theta\)-psh model functions.  Hence \(\varphi\) is bounded and
\(\theta\)-psh, by \Cref{thm:ch21-decreasing-stability}.

\emph{\(\varphi\) is not continuous, hence not a model function.}  The
points \(x_k\) converge to \(\zeta\).  Indeed, by the tree description
\Cref{prop:berk-a1-tree} a basic neighbourhood of \(\zeta\) is obtained by
deleting, in each of finitely many tangent directions at \(\zeta\), the
closed sub-branch at distance \(\geq\delta\) from \(\zeta\); each such
sub-branch is a closed disc (or, in the direction of \(\infty\), the
complement of an open disc), hence closed, so the complement of finitely
many of them is open.  A point \(x_k\) lies in this neighbourhood as soon as
its direction \(\bar a_k\) is not among the finitely many deleted ones, which
excludes only finitely many indices \(k\), the residues \(\bar a_k\) being
pairwise distinct.
But \(\varphi(x_k)=\psi_k(x_k)=-1\) for every \(k\), while
\(\varphi(\zeta)=0\).  So \(\varphi\) is discontinuous at \(\zeta\); it is
of course upper semicontinuous, being a decreasing limit of continuous
functions, and indeed \(\limsup_k\varphi(x_k)=-1\leq0=\varphi(\zeta)\).

\emph{The Monge--Amp\`ere measure.}  Since
\(\sum_{k\geq1}2^{-k}=1\), we get \(\ddc\varphi=-\delta_\zeta
+\sum_{k\geq1}2^{-k}\delta_{x_k}\) as measures, and
\Cref{thm:bounded-ma-extension} gives
\begin{equation}
\label{eq:ch24-branch-measure}
 \MA_\theta(\varphi)
 =\lim_{N\to\infty}\MA_\theta(\varphi_N)
 =\lim_{N\to\infty}\Bigl(2^{-N}\delta_\zeta
   +\sum_{k\leq N}2^{-k}\delta_{x_k}\Bigr)
 =\sum_{k\geq1}2^{-k}\delta_{x_k},
\end{equation}
the limit being weak.  The total mass is
\(\sum_{k\geq1}2^{-k}=1=V\) at every stage and in the limit, as
\Cref{thm:bounded-ma-extension} requires.  Note that the mass \(2^{-N}\) at
the Gauss point escapes in the limit, not to infinity---the space is
compact---but into the accumulating atoms; and that the limit measure has
infinitely many atoms, so it is not the measure of any model metric
(\Cref{prop:ch24-model-support}\,\cref{itm:ch24-ms:atomic}).  This is the
precise sense in which the bounded theory is strictly larger than the model
theory.

\emph{The energy pairing, computed.}  With \(n=1\),
\eqref{eq:ch24-mixed-energy} gives, for model functions \(f=\varphi_D\),
\(g=\varphi_{D'}\), the value
\(\mathfrak E(f,g)=(D\cdot\theta)+(D'\cdot\theta)+(D\cdot D')\).  Taking
\(f=g=\varphi_N\) and using
\(\int\varphi_N\,\theta=\varphi_N(\zeta)=0\) and
\(\int\varphi_N\,\ddc\varphi_N
 =\sum_{k\leq N}2^{-k}\bigl(\varphi_N(x_k)-\varphi_N(\zeta)\bigr)
 =-\sum_{k\leq N}2^{-k}\),
we find
\[
 \mathfrak E(\varphi_N,\varphi_N)=0+0-\sum_{k\leq N}2^{-k}
 =-(1-2^{-N})
 \ \searrow\ -1 .
\]
By \Cref{thm:ch24-energy-extension} the limit is
\(\mathfrak E(\varphi,\varphi)=-1\); and indeed
\eqref{eq:ch24-energy-telescoping} gives
\(\mathfrak E(\varphi,\varphi)=\int\varphi\,\theta
 +\int\varphi\,\MA_\theta(\varphi)
 =0+\sum_{k\geq1}2^{-k}\varphi(x_k)=-1\), which is the same number
computed from the extended operator.  This is a complete numerical check of
\eqref{eq:ch24-energy-slot-bounded}.

\emph{A remark on nets.}  Here a decreasing \emph{sequence} sufficed,
because the construction was countable by design.  That is a feature of the
example, not of the theory: by
\Cref{ex:ch22-countable-blind-spot} there are \(\theta\)-psh functions that
no countable family of model functions approximates from above, and for
those the net formulation of \Cref{thm:bounded-ma-extension} is the only one
available.  The example does show that continuity of the potential is not
needed anywhere in that theorem.
\end{example}

\subsection{A locality check}
\label{sec:ch24-ex-locality}

\begin{example}[Locality and the contact set on a chain of three
components]
\label{ex:ch24-locality-check}
Let \(n=1\) and let \(\cX\) be a regular semistable model of a curve whose
special fibre is a chain \(E_0,E_1,E_2\) of three components of
multiplicity one, with \((E_0\cdot E_1)=(E_1\cdot E_2)=1\) and
\((E_0\cdot E_2)=0\); the row-sum relation gives
\((E_0^2)=(E_2^2)=-1\) and \((E_1^2)=-2\).  The dual complex is the segment
\([0,2]\) with vertices \(v_0,v_1,v_2\) at the integers, and a model
function determined on \(\cX\) is the piecewise affine function with
prescribed vertex values.  Let \(\theta\) have component degrees
\((1,0,1)\), so \(V=2\) and \(\MA_\theta(0)=\delta_{v_0}+\delta_{v_2}\).

By \eqref{eq:ch20-curve-nef-test} a model function \(u\) with vertex values
\((u_0,u_1,u_2)\) is \(\theta\)-psh exactly when
\[
 1+(u_1-u_0)\geq0,
 \quad
 (u_0-u_1)+(u_2-u_1)\geq0,
 \quad
 1+(u_1-u_2)\geq0 .
\]

\emph{Two potentials.}  Take \(\varphi\) with values \((0,-1,-2)\) and
\(\psi:=-\tfrac12\), the constant.  Both are \(\theta\)-psh: for
\(\varphi\) the three left-hand sides are \(0,0,2\); for a constant they are
\(1,0,1\).  On the segment, parametrized by \(t\in[0,2]\),
\(\varphi(t)=-t\) and \(\psi(t)=-\tfrac12\).  Their measures are
\[
 \MA_\theta(\varphi)=0\cdot\delta_{v_0}+0\cdot\delta_{v_1}+2\delta_{v_2},
 \qquad
 \MA_\theta(\psi)=\delta_{v_0}+\delta_{v_2},
\]
both of mass \(2\).

\emph{The maximum.}  We have
\(\sigma:=\max\{\varphi,\psi\}\) equal to \(-t\) for \(t\leq\tfrac12\) and to
\(-\tfrac12\) for \(t\geq\tfrac12\).  This is a model function on the model
\(\cX'\) obtained by blowing up the node \(E_0\cap E_1\): the new component
has multiplicity \(b_0+b_1=2\), its divisorial point is the midpoint
\(v_{1/2}\) of the edge \([v_0,v_1]\), and the two halves of that edge have
length \(\tfrac12\).  The form \(\theta\) has degree \(0\) on the new
component, since \(\rho^*\) of a class is trivial on a contracted curve, and
its degrees on the remaining components are unchanged.  Summing outgoing
slopes with respect to the subdivided lengths, as in
\eqref{eq:graph-laplacian},
\[
 \MA_\theta(\sigma)
 =0\cdot\delta_{v_0}+1\cdot\delta_{v_{1/2}}
  +0\cdot\delta_{v_1}+1\cdot\delta_{v_2},
\]
of mass \(2\), as \Cref{thm:ch24-mass} requires.  Indeed at \(v_0\) the
degree is \(1+\bigl(\sigma(\tfrac12)-\sigma(0)\bigr)/\tfrac12=1-1=0\); at
\(v_{1/2}\) it is
\(0+\bigl(\sigma(0)-\sigma(\tfrac12)\bigr)/\tfrac12
 +\bigl(\sigma(1)-\sigma(\tfrac12)\bigr)/\tfrac12=1+0=1\); at \(v_1\) it is
\(0+0+0=0\), because \(\sigma\) is constant from \(t=\tfrac12\) onwards; and
at \(v_2\) it is \(1+0=1\).

\emph{The locality identity.}  Here \(\{\psi>\varphi\}=\{t>\tfrac12\}\) and
\(\{\varphi>\psi\}=\{t<\tfrac12\}\), both understood as their preimages
under the retraction.  Then
\[
 \one_{\{\psi>\varphi\}}\MA_\theta(\sigma)=\delta_{v_2}
 =\one_{\{\psi>\varphi\}}\MA_\theta(\psi),
\]
and
\[
 \one_{\{\varphi>\psi\}}\MA_\theta(\sigma)=0
 =\one_{\{\varphi>\psi\}}\MA_\theta(\varphi),
\]
both instances of \eqref{eq:model-maximum-principle}.  Note that
\(\MA_\theta(\sigma)\) carries an atom of mass \(1\) at the contact point
\(v_{1/2}\), which is present in neither \(\MA_\theta(\varphi)\) nor
\(\MA_\theta(\psi)\); this is the phenomenon of
\Cref{rem:ch24-contact-mass}, and it is exactly compensated by the mass
\(1\) that \(\MA_\theta(\psi)\) carries at \(v_0\) and
\(\MA_\theta(\sigma)\) does not.

\emph{The comparison principle.}  With
\(\{\varphi<\psi\}=\{t>\tfrac12\}\),
\[
 \int_{\{\varphi<\psi\}}\MA_\theta(\psi)=1
 \ \leq\
 \int_{\{\varphi<\psi\}}\MA_\theta(\varphi)=2 ,
\]
an instance of \eqref{eq:ch24-model-comparison}, with strict inequality.
\end{example}

\subsection{A maximum-principle application}
\label{sec:ch24-ex-maxprin}

\begin{example}[Solving the model equation on a curve, with uniqueness]
\label{ex:ch24-uniqueness-curve}
Keep the two-component model of \Cref{ex:two-vertex-ma}: special fibre
\(E_0+E_1\), one node, both multiplicities one, edge length one,
\(\theta=m_0\delta_{x_0}+m_1\delta_{x_1}\) with \(m_0+m_1=d=\deg L=V\).
Let
\[
 \mu=c_0\delta_{x_0}+c_1\delta_{x_1},
 \qquad c_0,c_1\geq0,\quad c_0+c_1=d ,
\]
be a target measure supported on the two divisorial points.

\emph{Existence.}  A model function \(\varphi\) with vertex values
\((a_0,a_1)\) satisfies
\(\MA_\theta(\varphi)=(m_0+a_1-a_0)\delta_{x_0}+(m_1+a_0-a_1)\delta_{x_1}\),
by \Cref{ex:two-vertex-ma}.  Hence \(\MA_\theta(\varphi)=\mu\) if and only
if
\[
 a_1-a_0=c_0-m_0 ,
\]
the second equation being the same one, because both measures have mass
\(d\).  Such \(\varphi\) exists, and it is \(\theta\)-psh precisely because
\(c_0,c_1\geq0\).  Thus the model Monge--Amp\`ere equation is solvable on
this model exactly for the measures of mass \(V\) supported on the two
divisorial points with nonnegative coefficients---which is
\Cref{thm:ch24-mass} read as a necessary condition, and the computation
above read as a sufficient one.

\emph{Uniqueness.}  Any two solutions differ by a constant, by
\Cref{cor:ch24-model-uniqueness-curve}; here this is visible, since the
solution is determined by \(a_1-a_0\).

\emph{A maximum-principle argument for the same conclusion.}  It is
instructive to obtain uniqueness from
\Cref{prop:ch24-model-comparison} instead, since that argument is the one
that generalizes.  Suppose \(\MA_\theta(\varphi)=\MA_\theta(\psi)=\mu\) and,
after adding a constant, that \(\max(\varphi-\psi)=0\).  If
\(\varphi-\psi\) is not identically \(0\), then
\(U:=\{\varphi<\psi\}\) is a nonempty open set, and
\eqref{eq:ch24-model-comparison} gives
\(\mu(U)=\int_U\MA_\theta(\psi)\leq\int_U\MA_\theta(\varphi)=\mu(U)\), which
is no contradiction by itself.  What does give a contradiction is the
sharper information available here: by
\Cref{prop:ch24-skeletal-maximum}\,\cref{itm:ch24-sm:model} the maximum of
the model function \(\varphi-\psi\) is attained at a divisorial point, say
\(x_0\), so \(\varphi-\psi\) attains its maximum \(0\) there and is
\(<0\) somewhere; then the linear equation \(\ddc(\varphi-\psi)=0\) forces
the outgoing slope of \(\varphi-\psi\) at \(x_0\) to vanish, hence
\((\varphi-\psi)(x_1)=(\varphi-\psi)(x_0)=0\), and \(\varphi-\psi\) is
constant after all.  This is the argument of
\Cref{cor:ch24-model-uniqueness-curve} written out on two vertices, and
\Cref{rem:ch24-uniqueness-scope} explains why it does not generalize as it
stands.

\emph{A numerical instance.}  For \(d=2\), \((m_0,m_1)=(1,1)\), and
\(\mu=2\delta_{x_1}\), one needs \(a_1-a_0=0-1=-1\); the normalized solution
is \((a_0,a_1)=(0,-1)\).  Its measure is
\((1-1)\delta_{x_0}+(1+1)\delta_{x_1}=2\delta_{x_1}=\mu\), and the energy
pairing of \Cref{ex:ch24-curve-mixed} evaluates to
\(\mathfrak E(\varphi,\varphi)=(a_0+a_1)+(a_0+a_1)-(a_0-a_1)^2
 =-1-1-1=-3\).
For \(\mu=\delta_{x_0}+\delta_{x_1}\) the solution is
\((a_0,a_1)=(0,0)\) and \(\mathfrak E=0\); the drop from \(0\) to \(-3\)
is the energy cost of concentrating the mass, and its sign is
\eqref{eq:model-dirichlet-negativity}.
\end{example}

\subsection{The toric projective line}
\label{sec:ch24-ex-toric-p1}

\begin{example}[The toric projective line]
\label{ex:ch24-toric-p1}
On \(\mathbf P^1\) with \(L=\mathcal O(d)\), an invariant semipositive metric
is encoded on the dense torus by a convex function with slopes in
\([0,d]\). For \(F(x)=\max\{0,dx\}\), the slope jumps from \(0\) to \(d\)
at \(x=0\). The lattice-normalized Alexandrov measure is \(d\delta_0\).
Under tropicalization this is the BFJ measure, and its mass is
\(d=(L^1)\). This checks both the absence of a hidden probability
normalization and the sign of \(\ddc\).
\end{example}

\section{Scope: which hypothesis is used where}
\label{sec:ch24-scope}

The standing hypotheses \eqref{eq:ch24-standing} are not used uniformly.  The
table records, for each of the principal statements, which of them the proof
given here actually uses, and what is known about removing it.  As always in
\Cref{part:bfj}, the original BFJ scope is kept visibly distinct from later
extensions.

\begin{center}
\small
\begin{tabular}{p{.28\textwidth}|p{.15\textwidth}|p{.15\textwidth}|p{.32\textwidth}}
\toprule
\textbf{Statement}&\textbf{Discrete valuation}&\textbf{Residue char.\ \(0\)}
 &\textbf{Other hypotheses used}\\
\midrule
\Cref{lem:ch24-vertical-numbers},
\Cref{prop:model-ma-basic},
\eqref{eq:mixed-total-mass}
 &yes, for the model theory of \Cref{chap:models}
 &no
 &\(X\) proper, flat projective models; not smoothness, not ampleness\\
\Cref{prop:model-integration-by-parts},
\Cref{prop:ch24-model-cauchy-schwarz}
 &yes
 &only through cofinality of SNC models in
  \Cref{thm:ch20-mixed-hodge}
 &regular models, relative Hodge index\\
\Cref{prop:ch24-energy-basic},
\Cref{thm:ch24-mass}
 &yes
 &no
 &multilinearity of intersection numbers only\\
\Cref{lem:ch24-psh-difference},
\Cref{lem:ch24-dense-tests}
 &yes
 &no
 &\(\{\theta\}\) ample; density
  \Cref{prop:metrics-model-function-density}\\
\Cref{thm:ch24-energy-extension},
\Cref{thm:bounded-ma-extension}
 &yes
 &\textbf{yes}, only through
  \Cref{thm:bfj-decreasing-regularization}
 &compactness of \(X^{\An}\); \(\{\theta\}\) ample\\
\Cref{prop:bounded-ma-no-polar-mass}
 &yes
 &as above
 &finiteness of psh functions at quasimonomial points\\
\Cref{prop:model-maximum-principle},
\Cref{prop:ch24-model-comparison}
 &yes
 &only through cofinality of SNC models
 &\Cref{prop:ch20-model-max}, hence smoothness of \(X\)\\
\Cref{thm:bfj-cld-ma-common-locus}, model part
 &yes, for \(b_E\) to be the multiplicity
 &no
 &\(X\) proper reduced only; see
  \Cref{rem:ch19-cld-bfj-hypotheses}\\
\Cref{thm:bfj-cld-ma-common-locus}, continuous part
 &yes
 &no
 &density of model functions\\
\bottomrule
\end{tabular}
\end{center}

Three comments on the table.

\emph{Residue characteristic zero enters once.}  It is used only through
\Cref{thm:bfj-decreasing-regularization}, that is, only to know that every
\(\theta\)-psh function is a decreasing limit of \(\theta\)-psh model
functions.  Everything else in \Cref{sec:ch24-bounded} is a formal
consequence of that statement together with intersection theory.  If
decreasing regularization is available by other means---for instance for
curves over an arbitrary complete field, or in the positive-characteristic
situations opened up by test ideals in
\Cref{sec:ch23-later-test-ideals}---then \Cref{thm:bounded-ma-extension}
follows verbatim.  This is why the two statements are separated in this
book: the analytic content of the extension theorem is
\Cref{lem:ch24-sandwich}, and the arithmetic content is elsewhere.

\emph{Ampleness is used for the test space, not for positivity.}  The
hypothesis that \(\{\theta\}\) be ample enters through
\Cref{lem:ch24-psh-difference}, whose role is to guarantee that
\(\theta\)-psh model functions span a dense subspace of
\(C^0(X^{\An})\).  If \(\{\theta\}\) is merely nef and big, one may still run
the construction on the ample locus, but the density statement, and with it
the uniqueness clause of \Cref{thm:bounded-ma-extension}, requires
modification; we do not pursue this, and \Cref{chap:bfj-energy} likewise
keeps \(\{\theta\}\) ample.

\emph{Discreteness is structural.}  Every statement of this chapter is
formulated in terms of models over \(R=K^\circ\) and the multiplicities
\(b_E\) of a special fibre.  Over a general nontrivially valued field the
Chambert--Loir measure of a formal metric still exists
(\Cref{sec:metrics-cl-measure} and \Cref{thm:cld-formal-intersection}), but
the BFJ framework of \Cref{chap:bfj-model-forms} has no substitute in the
form used here.  The reader interested in that generality should consult the
scope discussion in \Cref{rem:ch19-cld-bfj-hypotheses} and the treatment of
Zhang semipositivity over arbitrary complete fields in \cite{GM19}.

\section{Dictionary}
\label{sec:ch24-dictionary}

\subsection{Forward}
\label{sec:ch24-dict-forward}

\begin{center}
\small
\begin{tabular}{p{.32\textwidth}|p{.30\textwidth}|p{.30\textwidth}}
\toprule
\textbf{Algebraic object}&\textbf{Measure-theoretic object}
 &\textbf{Where}\\
\midrule
mixed degree \(b_E(\theta_1|_E\cdots\theta_n|_E)\)
 &atom of the mixed measure at \(x_E\)
 &\eqref{eq:model-mixed-measure},
  \Cref{def:mixed-model-measure}\\
vertical intersection number \((D\cdot\theta_1\cdots\theta_n)\)
 &\(\int\varphi_D\,\theta_1\wedge\cdots\wedge\theta_n\)
 &\eqref{eq:model-test-intersection},
  \eqref{eq:metrics-cl-pairing}\\
generic mixed intersection number
 &total mass of the mixed measure
 &\eqref{eq:mixed-total-mass},
  \Cref{thm:ch24-mass}\\
sum of vertical intersection numbers
 &mixed energy pairing \(\mathfrak E\)
 &\eqref{eq:ch24-mixed-energy},
  \Cref{prop:ch24-energy-basic}\\
\((D\cdot D'\cdot\eta_1\cdots\eta_{n-1})\), symmetric
 &integration by parts
 &\eqref{eq:model-ibp}\\
\((D^2\cdot\eta_1\cdots\eta_{n-1})\leq0\)
 &nonpositivity of the Dirichlet form
 &\eqref{eq:model-dirichlet-negativity},
  \eqref{eq:ch24-bounded-dirichlet}\\
divisors agreeing near \(E\)
 &locality of \(\MA\) at \(x_E\)
 &\Cref{prop:model-maximum-principle}\\
relative nefness of a model line bundle
 &positivity of the measure
 &\Cref{prop:model-ma-basic}, Step~3\\
decreasing net of vertical divisors
 &weak convergence of measures
 &\Cref{thm:bounded-ma-extension}\\
\bottomrule
\end{tabular}
\end{center}

\subsection{Backward}
\label{sec:ch24-dict-backward}

Four entries above are one-way, and it is worth saying which.

\emph{From a measure back to a potential} there is no arrow in this chapter.
Solving \(\MA_\theta(\varphi)=\mu\) is the subject of
\Cref{chap:bfj-finite-energy-measures,chap:bfj-continuity}, and even the
model-level solvability of \Cref{ex:ch24-uniqueness-curve} is special to
\(n=1\), where the operator is linear.

\emph{From the bounded product back to a model product} there is no arrow:
\Cref{ex:ch24-branch-sum} produces a bounded \(\theta\)-psh function whose
Monge--Amp\`ere measure has infinitely many atoms and is therefore not the
measure of any model metric.

\emph{From the BFJ product to a CLD Bedford--Taylor product} the arrow
exists only on the model and continuous loci, by
\Cref{thm:bfj-cld-ma-common-locus}; the four statements that are not proved
are listed in \Cref{sec:ch24-scope-comparison}.

\emph{From locality at model level to locality for bounded potentials} the
arrow requires capacity, and passes through
\Cref{thm:bounded-quasicontinuity}; see the table of
\Cref{sec:ch24-needs-capacity}.

\subsection{Numbers to remember}
\label{sec:ch24-dict-numbers}

\begin{center}
\small
\begin{tabular}{p{.36\textwidth}|p{.24\textwidth}|p{.32\textwidth}}
\toprule
\textbf{Quantity}&\textbf{Value}&\textbf{Where}\\
\midrule
total mass of every mixed
 \(\theta\)-Monge--Amp\`ere measure
 &\(V=(\{\theta\}^n)\)
 &\eqref{eq:mixed-total-mass},
  \Cref{thm:ch24-mass}\\
total variation of \(\ddc h\wedge\Theta\), \(h\) \(\theta\)-psh
 &\(\leq2V\)
 &\Cref{sec:ch24-model-ibp}\\
Lipschitz constant of \(\mathfrak E\) in one slot
 &\(V\)
 &\Cref{thm:ch24-energy-extension}\,\cref{itm:ch24-ee:lip}\\
modulus in
 \eqref{eq:ch24-model-uniform-estimate}
 &\(2V\sum_i\|\varphi_i-\psi_i\|_\infty\)
 &\Cref{prop:ch24-model-uniform-estimate}\\
mixed measure on the toric surface
 &\(\delta_{x_0}\), mass \(1\)
 &\eqref{eq:ch24-toric-mixed}\\
Minkowski area at \(u=(0,0)\) in the same example
 &\(\tfrac32\)
 &\Cref{ex:ch24-toric-surface-mixed}\\
measure of the branch sum
 &\(\sum_{k\geq1}2^{-k}\delta_{x_k}\), mass \(1\)
 &\eqref{eq:ch24-branch-measure}\\
energy pairing of the branch sum
 &\(-1\)
 &\Cref{ex:ch24-branch-sum}\\
contact atom in the locality check
 &\(1\) at \(t=\tfrac12\)
 &\Cref{ex:ch24-locality-check}\\
comparison-principle instance
 &\(1\leq2\)
 &\Cref{ex:ch24-locality-check}\\
\bottomrule
\end{tabular}
\end{center}

\section*{Notes and further reading}

The model formula, integration by parts, and bounded extension are
\cite[Proposition--Definition~2.18, Propositions~2.19--2.21, and
Theorem~3.1]{BFJ15}. The noncharging statement is
\cite[Proposition~3.7]{BFJ15}. Agreement with the original Chambert--Loir
measure is \cite[\S3.2]{BFJ15}. Only model and continuous semipositive
comparison loci have been asserted here.

\paragraph{Provenance, statement by statement.}
\Cref{def:mixed-model-measure} and \Cref{prop:model-ma-basic} are
\cite[\S2.7]{BFJ15}; the measure itself goes back to Chambert--Loir
\cite{CL06}, and the version for arbitrary complete fields, together with
the comparison of algebraic, formal, and piecewise linear semipositivity, is
in \cite{GM19}.  \Cref{prop:model-integration-by-parts} and the negativity
\eqref{eq:model-dirichlet-negativity} are \cite[\S2]{BFJ15}, resting on the
relative Hodge index theorem proved here as \Cref{thm:ch20-mixed-hodge}.
\Cref{thm:bounded-ma-extension} is \cite[Theorem~3.1]{BFJ15}; the earlier
paper \cite{BFJ16} contains the psh formalism it uses, and the survey
\cite{BFJ16b} is a short entry point.  \Cref{prop:model-maximum-principle}
is the model case of the locality statement of \cite[\S3]{BFJ15}, whose
bounded form is
\Cref{thm:bounded-maximum-principle} here.  The trivially valued analogue of
the entire chapter is developed in \cite{BJ22}, and the synthetic treatment of
measures of finite energy that supersedes several of the later arguments is
\cite{BJ25FiniteEnergy}.  The comparison with the Chambert--Loir--Ducros
calculus, \Cref{thm:bfj-cld-ma-common-locus}, is proved in this book as
\Cref{thm:ch19-cld-bfj-comparison}, on the basis of
\cite[Theorem~18.5.5 and Proposition~18.5.8]{CLD25}; the delta-form version
is due to Gubler--K\"unnemann \cite{GK17,GK19}, and the broader comparisons
available under semiampleness hypotheses are
\cite[Theorems~7.14 and~8.4]{BE21}.  The toric identification used in
\Cref{ex:ch24-toric-surface-mixed} is classical in the arithmetic form of
\cite{BGPS14} and in the tropical form of \cite{BGJKK25}.

\paragraph{What is organized differently here.}
Three choices depart from the arrangement in the sources.  First, the
extension theorem is proved by extending the \emph{scalar} mixed energy
pairing \(\mathfrak E\) of \Cref{def:ch24-mixed-energy} and recovering the
measure by the Riesz theorem, rather than by an induction on the number of
non-model slots.  The advantage is that symmetry, multilinearity, monotone
continuity and the weighted statement \eqref{eq:weighted-monotone-ma} all
become properties of a monotone function of \(n+1\) real variables, and that
the only analytic input is the sandwich lemma
\Cref{lem:ch24-sandwich}.  The inductive formulae of the original argument
are retained as \Cref{prop:ch24-forced-ibp} and
\Cref{prop:ch24-approximate-monotonicity}, now with proofs, because they are
the shape in which the result is usually quoted and because the second is the
quantitative statement behind the qualitative one.  Second, the distinction
between a vertical divisor as a cycle and as a line-bundle class has been
made into a lemma, \Cref{lem:ch24-vertical-numbers}, rather than left as a
convention; without it the total mass \eqref{eq:mixed-total-mass} can be
computed to be zero, a mistake the sources warn against implicitly and the
warning after \Cref{lem:ch24-vertical-numbers} warns against explicitly.
Third, the four Monge--Amp\`ere products of \Cref{sec:ch24-four-products} are
tabulated at the outset, and every statement of the chapter names the one it
concerns.

\paragraph{Interfaces.}
The measures of this chapter are the input to capacity
(\Cref{chap:bfj-capacity}), whose very definition
\eqref{eq:capacity-definition} is a supremum of the masses constructed here,
and to the energy functional of \Cref{chap:bfj-energy}, which by
\Cref{rem:ch24-energy-vs-E} is \((n+1)^{-1}\mathfrak E\) on the diagonal.
The extension to finite energy, in which the mass may drop and the
nonpluripolar convention \Cref{def:nonpluripolar-ma} is imposed, is
\Cref{thm:e1-mixed-products}; that extension is \emph{not} a limit of the
present one, and the two are compared in
\Cref{rem:cma-no-universal-e1}.  Functoriality of the construction under
finite and birational morphisms and under base change is
\Cref{thm:finite-ma-projection} and \Cref{thm:base-change-ma}.  The complete
worked catalogue of curves, toric varieties and abelian varieties is
\Cref{chap:bfj-examples}, and the theorem-by-theorem compatibility table
with the CLD calculus is \Cref{sec:ch19-limitations} together with
\Cref{chap:compare-ma}.

\paragraph{What is new here, and what is not.}
No theorem of this chapter is new.  \Cref{lem:ch24-vertical-numbers},
\Cref{prop:ch24-polarization}, \Cref{prop:ch24-energy-basic},
\Cref{prop:ch24-model-uniform-estimate},
\Cref{prop:ch24-model-comparison} and
\Cref{cor:ch24-model-uniqueness-curve} are elementary consequences of
statements in the sources; they have been written out because each is used
later in a form that is easy to misremember.  The proof of
\Cref{thm:bounded-ma-extension} given here is a rearrangement, not a new
argument: every ingredient---monotonicity of intersection numbers, Dini's
theorem for nets, the Riesz theorem---is in \cite{BFJ15}, and only the order
is different.  \Cref{ex:ch24-branch-sum} has no counterpart in the sources
and is the smallest example we know in which the bounded operator is computed
exactly on a potential that is neither a model function nor continuous.
